% !TeX program = xelatex
% ============================================================
%  SY02-19 《几何朗兰兹的一个算术应用》
%  底本：Sam Raskin, "An arithmetic application of geometric Langlands",
%        Proc. Sympos. Pure Math. 112.1 (2025), 42 pp.
%  编译：xelatex -> xelatex -> xelatex
% ============================================================
\documentclass[UTF8, a4paper, 12pt, fontset=fandol, scheme=chinese]{ctexart}
% ============================================================
%  SY02-19 《几何朗兰兹的一个算术应用》译稿导言区
%  底本：Sam Raskin, An Arithmetic Application of Geometric Langlands,
%        Proc. Sympos. Pure Math. 112.1 (2025), 42 pp.
%  体例：中文数学论文；定理类编号「节.小节.子小节.序号」，公式编号「节.序号」
% ============================================================

\usepackage{geometry}
\geometry{a4paper, left=2.7cm, right=2.7cm, top=2.8cm, bottom=2.7cm}
\linespread{1.22}
\setlength{\parskip}{0.25em}
\setlength{\emergencystretch}{2em}
\sloppy
\hyphenpenalty=10000
\tolerance=3000

% ---------- 数学 ----------
\usepackage{amsmath}
\usepackage{amsthm}
\usepackage{amssymb}
\usepackage{amsfonts}
\usepackage{mathtools}
\usepackage{mathrsfs}
\usepackage{stmaryrd}
\usepackage{bm}
\allowdisplaybreaks[2]

% ---------- 图表 ----------
\usepackage{graphicx}
\usepackage{xcolor}
\usepackage{tikz}
\usetikzlibrary{cd, arrows.meta, positioning, calc}
\usepackage[most]{tcolorbox}
\usepackage{enumitem}
\setlist{itemsep=2pt, topsep=3pt, parsep=0pt}

% ---------- 超链接 ----------
\usepackage[hyperfootnotes=false, colorlinks=true, linkcolor=blue!45!black, citecolor=green!35!black,
            urlcolor=teal!60!black, bookmarksnumbered=true]{hyperref}
\usepackage{bookmark}

% ---------- 中文微调 ----------
\xeCJKsetup{CJKecglue={\hskip 0.24em plus 0.05em minus 0.05em}}
\ctexset{
  section    = {name={,}, number=\arabic{section},
                aftername={.\enspace}, format=\Large\bfseries\centering},
  subsection = {number=\arabic{section}.\arabic{subsection},
                aftername={.\enspace}, format=\large\bfseries},
  subsubsection = {number=\arabic{section}.\arabic{subsection}.\arabic{subsubsection},
                aftername={.\enspace}, format=\normalsize\bfseries}
}

% ---------- 定理环境：共用计数器，按子小节重置（复现 1.2.4.1、4.3.3.1 体例） ----------
\theoremstyle{plain}
\newtheorem{thm}{定理}[subsubsection]
\newtheorem{prop}[thm]{命题}
\newtheorem{lem}[thm]{引理}
\newtheorem{cor}[thm]{推论}
\newtheorem{conj}{猜想}[subsection]
\newtheorem{notation}[thm]{记号}
\newtheorem{question}[thm]{问题}
\newtheorem{hypothesis}[thm]{假设}
\theoremstyle{definition}
\newtheorem{defn}[thm]{定义}
\newtheorem{example}[thm]{例}
\newtheorem{convention}[thm]{约定}
\theoremstyle{remark}
\newtheorem{remark}[thm]{注}
\newtheorem{warning}[thm]{警告}

% 定理 A 是原文专名，不参与自动编号
\newtheorem*{thmA}{定理 A}

% 原文把「1.1.1. 背景」这类编号段落作为正文编号，而非标题；
% \pnum 递增 subsubsection 计数器并run-in 输出，定理编号随之对齐
\newcommand{\pnum}[1]{\par\refstepcounter{subsubsection}%
  \noindent\textbf{\thesubsubsection.}\ #1\ }
\newcommand{\pnumt}[2]{\par\refstepcounter{subsubsection}%
  \noindent\textbf{\thesubsubsection.}\ \textbf{#1}\ #2\ }

\numberwithin{equation}{section}
\renewcommand{\proofname}{证明}
\renewcommand{\qedsymbol}{$\square$}
\renewcommand{\refname}{参考文献}
\renewcommand{\contentsname}{目\quad 录}

% ---------- 译注 ----------
\usepackage{manyfoot}
\DeclareNewFootnote[plain]{trans}
\renewcommand{\thefootnotetrans}{译\arabic{footnotetrans}}
\newcommand{\zy}[1]{\footnotetrans{#1}}
\newtcolorbox{yizhu}[1][]{
  enhanced, breakable, colback=orange!4!white, colframe=orange!55!black,
  boxrule=0.5pt, left=6pt, right=6pt, top=4pt, bottom=4pt,
  fonttitle=\bfseries\small, title=#1,
  before skip=6pt, after skip=6pt
}

% ---------- 本文记号 ----------
% ============================================================
%  SY02-19 记号宏（按原文用法整理）
% ============================================================

% 常用缩写
\newcommand{\LS}{\mathrm{LS}}
\newcommand{\Bun}{\mathrm{Bun}}
\newcommand{\Shv}{\mathrm{Shv}}
\newcommand{\IndCoh}{\mathrm{IndCoh}}
\newcommand{\QCoh}{\mathrm{QCoh}}
\newcommand{\Coh}{\mathrm{Coh}}
\newcommand{\Perf}{\mathrm{Perf}}
\newcommand{\Rep}{\mathrm{Rep}}
\newcommand{\Nilp}{\mathrm{Nilp}}
\newcommand{\qLisse}{\mathrm{qLisse}}
\newcommand{\Eis}{\mathrm{Eis}}
\newcommand{\Frob}{\mathrm{Frob}}
\newcommand{\Dmod}{\mathrm{D}\text{-}\mathrm{mod}}
\newcommand{\DGCat}{\mathrm{DGCat}}
\newcommand{\Vect}{\mathrm{Vect}}
\newcommand{\Spectra}{\mathrm{Spectra}}
\newcommand{\Exc}{\mathrm{Exc}}
\newcommand{\Sht}{\mathrm{Sht}}
\newcommand{\Eisspec}{\mathrm{Eis}^{\mathrm{spec}}}
\newcommand{\Eissoyt}{\mathrm{Eis}^{\mathrm{toy}}}
\newcommand{\arthm}{\mathrm{arthm}}
\newcommand{\restr}{\mathrm{restr}}
\newcommand{\cusp}{\mathrm{cusp}}
\newcommand{\unr}{\mathrm{unr}}
\newcommand{\nonres}{\mathrm{non}\text{-}\mathrm{res}}
\newcommand{\triv}{\mathrm{triv}}
\newcommand{\op}{\mathrm{op}}
\newcommand{\et}{\text{ét}}
\newcommand{\dR}{\mathrm{dR}}

% 算子
\DeclareMathOperator{\Aut}{Aut}
\DeclareMathOperator{\Hom}{Hom}
\DeclareMathOperator{\Ext}{Ext}
\DeclareMathOperator{\Spec}{Spec}
\DeclareMathOperator{\ad}{ad}
\DeclareMathOperator{\Ad}{Ad}
\DeclareMathOperator{\ch}{ch}
\DeclareMathOperator{\vol}{vol}
\DeclareMathOperator{\cl}{cl}
\DeclareMathOperator{\tr}{tr}
\DeclareMathOperator{\id}{id}
\DeclareMathOperator{\Sect}{Sect}
\DeclareMathOperator{\colim}{colim}

% 常用对象
\newcommand{\bG}{\check{G}}
\newcommand{\bT}{\check{T}}
\newcommand{\bB}{\check{B}}
\newcommand{\bN}{\check{N}}
\newcommand{\bg}{\check{\mathfrak g}}
\newcommand{\bt}{\check{\mathfrak t}}
\newcommand{\hG}{\widehat{G}}
\newcommand{\hT}{\widehat{T}}
\newcommand{\hB}{\widehat{B}}
\newcommand{\hZ}{\widehat{Z}}
\newcommand{\hf}{\widehat{f}}
\newcommand{\fh}{\widehat{h}}
\newcommand{\wt}{\widetilde}
\newcommand{\Om}{\Omega}
\newcommand{\om}{\omega}
\newcommand{\e}{\mathrm{e}}
\newcommand{\F}{\mathbb{F}}
\newcommand{\Z}{\mathbb{Z}}
\newcommand{\Q}{\mathbb{Q}}
\newcommand{\A}{\mathbb{A}}
\newcommand{\OO}{\mathcal{O}}
\newcommand{\ZZ}{\mathbb{Z}}
\newcommand{\QQ}{\mathbb{Q}}
\newcommand{\Fq}{\F_q}
\newcommand{\Fb}{\overline{\F}_q}



\begin{document}
\pagestyle{plain}

% ============ 扉页 ============
\begin{titlepage}
  \centering
  \vspace*{2.2cm}
  {\Huge\bfseries 几何朗兰兹的一个算术应用\par}
  \vspace{0.8cm}
  {\LARGE An Arithmetic Application of Geometric Langlands\par}
  \vspace{1.2cm}
  {\large Sam Raskin\par}
  \vspace{0.4cm}
  {\normalsize Yale University\par}
  \vspace{1.6cm}
  {\large 原文：Proc. Sympos. Pure Math. 112.1 (2025), 42 页\par}
  \vspace{0.4cm}
  {\large DOI: 10.1090/pspum/112.1/02062\par}
  \vspace{0.5cm}
  {\large 译者：数译（DSH agent）\par}
  \vspace{0.3cm}
  {\large 版本：v1.0\qquad 日期：2026 年 11 月\par}
  \vfill
  {\small 本译稿遵循本库《学术翻译总纲》与《数学文献翻译纪律 v2.1》。\par}
\end{titlepage}

% ============ 译者说明 ============
\section*{译者说明}
\addcontentsline{toc}{section}{译者说明}

\textbf{原文性质。}\ 本文是 Sam Raskin 2025 年发表的研究论文，正文 42 页，刊于《Proceedings of Symposia in Pure Mathematics》第 112.1 卷。它不是综述：AGKRRV 的几何朗兰兹理论（[2][3][4]）是全文的外部支柱，作者在此之上构造谱侧 Eisenstein 级数、证明 \verb|LS^arthm| 在平凡局部系统附近的显式坐标，并最终证明函数域上平凡 Langlands 参数处无无分歧尖形式（定理 A）。译文保留全部证明与交叉引用，不省略技术段落；作者说明可独立阅读的 §4.6.5 也照译，并加提示。

\textbf{编号约定。}\ 原文的定理类对象按"节.小节.子小节.序号"编号（如 Theorem 4.3.3.1），公式按节编号（如 (4.7)）。本译稿用 \verb|\newtheorem{...}[subsubsection]| 与 \verb|\numberwithin{equation}{section}| 复现这两套编号；"定理 A"是原文给主定理的专名，不是自动编号，译文照写"定理 A"。正文中的交叉引用一律按原文编号排印（如"定理 4.3.3.1"），不依赖 \verb|\ref| 计数，以免分块编译时节号偏移导致错位。

\textbf{译注标记。}\ 解释性译注记作"译注 $n$"，用独立脚注序列（不占作者脚注编号）；较长的说明放在橙色"译注"框中。凡解释性文字一律附来源（\verb|[来源：...]|），或明确标为译者推断（\verb|[译者按...]|）；来源按手册 A/B/I 三级执行，不使用黑名单站点。原文的少量笔误在正文中保留原样，并以源文件注释 \verb|% [Tutor Note: ...]| 标出，同时在 \verb|资料卡/04-术语文献勘误.md| 与 \verb|工程文件/work/疑点清单.tsv| 中登记，不悄悄改动原文。参考文献中的法语重音按原刊补全；§4.6.5 的局部 $\varepsilon$-因子按上下文排印并登记。

\textbf{术语约定。}\ 术语先冻结后使用，关键几条：geometric Langlands program 几何朗兰兹纲领；restricted variation 受限变差；restricted GLC 受限几何朗兰兹猜想；local system 局部系统；lisse sheaf 光滑层；perverse sheaf 反常层；nilpotent singular support 幂零奇异支撑；categorical trace 范畴迹；shtuka 保留原文 shtuka；unipotent 幂幺；non-resonant 非共振；formal étale 形式 étale；Grothendieck--Springer theory Grothendieck--Springer 理论；Eisenstein series Eisenstein 级数；pseudo-Eisenstein series 伪 Eisenstein 级数。缩写 \verb|Bun|、\verb|LS|、\verb|Shv|、\verb|IndCoh|、\verb|QCoh|、\verb|qLisse| 保留英文。完整对照见 \verb|工程文件/work/glossary.tsv| 与 \verb|资料卡/04-术语文献勘误.md|。

\newpage
\tableofcontents
\newpage

% ============ 正文 ============
% ============================================================
%  ch00：标题、摘要、§1 引言（原文 pp.1–4）
% ============================================================
\begin{center}
  {\LARGE\bfseries 几何朗兰兹的一个算术应用}\\[6pt]
  {\large Sam Raskin}\\[2pt]
  {\normalsize Yale University}
\end{center}

\vspace{0.6em}
\begin{center}
\begin{minipage}{0.92\textwidth}
\small
\noindent\textbf{摘要.}\quad Vincent Lafforgue 为函数域上的尖自守函数空间构造了 Langlands 分解。在我们与 Arinkin、Gaitsgory、Kazhdan、Rozenblyum、Varshavsky 的合作工作中，我们证明了几何朗兰兹猜想的一个版本给出 Lafforgue 分解的特征空间在处处无分歧情形下的描述。

在本文中，我们概述后一个思想圈。然后我们说明如何用这些方法证明：在\emph{平凡} Langlands 参数处不存在处处无分歧的紧支撑尖形式，从而回答 Michael Harris 的一个问题。

作为一项具有一定独立意义的结果，我们在平凡 Langlands 参数附近以某些显式项计算伪 Eisenstein 级数的谱类比。在合适的坐标下，我们发现它是 Weyl 特征标公式与一个和曲线相关的 zeta 因子的乘积。
\end{minipage}
\end{center}

\vspace{0.8em}

\section{引言}

\subsection{背景与目标}

\pnum{在我们与 Arinkin、Gaitsgory、Kazhdan、Rozenblyum、Varshavsky 的合作工作 [2]、[3]、[4] 中，我们提出了无分歧、紧支撑自守函数的谱分解的一个猜想性表述。这一猜想也由朱歆文在 [32] 中独立发现。}

我们的谱分解受 V.~Lafforgue 的突破 [22] 与 Beilinson--Drinfeld 的几何朗兰兹猜想启发，属于 Langlands 型的分解，但性质不同：它描述的是全部（紧支撑）自守函数，而不仅是特征函数；并且它同时给出互反律与函子性的陈述，而不必把其中任何一个明确纳入表述。

我们工作的一个主要目的是证明：这一谱分解确实可由几何朗兰兹猜想的 $\ell$-进版本推出。

\pnum{在宣传我们的合作工作时——包括我在 IHÉS 的若干报告——我一直试图论证：我们的猜想能给出关于自守函数的、令数论学家感兴趣的新洞察。我们的猜想在最显而易见的情形下最令人满意，即\emph{离散}（亦称\emph{椭圆}）Langlands 参数（参见例 2.4.5.1）。但在其他 Langlands 参数附近，我一直难以给出精确、具体的推论。}

例如，就目前的形式而言，我们的猜想并不能立即重新得到离散系列的 Arthur 重数公式。要理解这类公式，还需要理论的进一步发展。

\pnum{话虽如此，本文的主要新贡献是给出我们工作对自守函数的一个简单、具体的应用，回答 Michael Harris 的一个问题。这个断言所涉及的是平凡 Langlands 参数——它几乎与离散参数相距最远。}

\pnum{此外，为与会议论文集相称，我们在 §2 中对几何朗兰兹纲领以及 [2]、[3]、[4] 中的思想圈作一些介绍。这两部分可以基本独立地阅读。}

最关心这部分综述材料的读者可以跳到 §2；因为那部分材料按其性质是导引性的，而引言余下部分的重点放在本文后半所考虑的问题上。

\subsection{主结果的陈述}

\pnumt{设定。}{}\ 我们固定一个特征为 $p$ 的有限域 $\F_q$，并令 $k = \Fb$ 表示它的代数闭包。我们固定一个素数 $\ell \neq p$，并令 $\e$ 表示 $\overline{\Q}_\ell$；在 [2] 的术语中，这是\emph{系数域}。我们固定一个分裂约化群 $G/\F_q$，并令 $\bG/\e$ 表示它的朗兰兹对偶群。

令 $X_0/\F_q$ 为一条光滑、射影、几何连通的曲线，并令 $X = X_0 \times_{\F_q} k$ 表示它到 $k$ 的基变换。

我们令 $F = \F_q(X_0)$ 表示与 $X_0$ 相伴的整体域。我们令 $\A$ 表示它的 adele 环，并令 $\OO \subseteq \A$ 表示整 adele 的子环。

我们令 $\Aut^{\unr}_{G,c}$ 表示 $F$ 上处处无分歧、紧支撑自守函数的空间。按定义，这意味着 $\Aut^{\unr}_{G,c}$ 是函数
\[
  G(F)\backslash G(\A)/G(\OO) \longrightarrow \e
\]
的向量空间，且这些函数具有有限支撑。我们令 $\Aut^{\unr}_{\cusp} \subseteq \Aut^{\unr}_{G,c}$ 表示\emph{尖}自守形式的子空间。\footnote{我们提醒读者一个众所周知的事实：在函数域上，尖自守形式\emph{先验地}就是紧支撑的。}

\pnum{我们一次性固定 $X$ 的一个 $k$-点作为我们的基本群的基点；我们在记号中省略它。\footnote{更好的做法是按 $X$ 上 lisse 层的范畴来思考，我们经常这样做。这里用 $\pi_1^{\et}$ 只是为了使讨论中的某些点更具体。}我们令 $\pi_1^{\et}(X)$ 表示 $X$ 的 \et{} 基本群，令 $\pi_1^{\arthm}(X) := \pi_1^{\et}(X_0)$ 表示算术基本群，并令 $\mathcal{W}_X := \pi_1^{\arthm}(X) \times_{\ZZ} \ZZ$ 表示 $X$ 的 Weil 群（赋以其标准拓扑，使得 $\pi_1^{\et}(X) \subseteq \mathcal{W}_X$ 是开子群）。}

\begin{notation}
为确定起见，我们始终使用\emph{几何} Frobenius 约定。于是我们通过一个生成元对应于\emph{几何} Frobenius 元素，来等同 $\widehat{\ZZ} \simeq \pi_1^{\et}(\Spec(\F_q))$。
\end{notation}

\pnumt{Lafforgue--Langlands 分解。}{}\ 暂时假定 $G$ 是半单的，以简化讨论。（本文正文对一般约化群工作。）

一个 \emph{Langlands 参数}是连续同态 $\rho: \mathcal{W}_X \to \bG(\e)$。一个 Langlands 参数称为\emph{半单}的，如果对任意抛物子群 $\check{P} \subseteq \bG$，只要 $\rho$ 通过 $\check{P}(\e)$ 分解，就存在 Levi 因子 $\check{M} \subseteq \check{P}$，使得 $\rho$ 进一步通过 $\check{M}(\e)$ 分解（参见 [2] §3.5--3.6）。

我们现在回忆，[22] 构造了分解
\begin{equation}
  \Aut^{\unr}_{\cusp} \cong \bigoplus_{[\sigma]} \Aut^{\unr}_{\cusp,[\sigma]},
\end{equation}
其中 $[\sigma]$ 取遍半单 Langlands 参数的共轭类。

\begin{remark}
上述讨论对分歧自守函数同样适用。我们的主要结果限制在无分歧情形，因此我们选择在全文强调无分歧设定。
\end{remark}

\pnum{}本文的主要结果如下：

\begin{thmA}
令 $G$ 为半单群（且非平凡群）。令 $\triv: \mathcal{W}_X \to \bG(\e)$ 表示平凡 Langlands 参数，即取值于单位元的常值映射。

假定 [2] 的受限几何朗兰兹猜想及其与 Eisenstein 级的相容性。那么直和项 $\Aut^{\unr}_{\cusp,[\triv]} \subseteq \Aut^{\unr}_{\cusp}$ 为零。换言之，不存在具有平凡 Langlands 参数的无分歧尖形式。
\end{thmA}

\begin{remark}
尽管我们在正文中没有强调这一点，但可以少用一些假设。确切地说，根据 [2] 推论 14.3.5，$\Shv_{\Nilp}(\Bun_G)$ 分解为 $X$ 上半单 $\bG$-局部系统上的一个直和。人们只需要平凡局部系统处的受限几何朗兰兹。我们期待后续工作完全处理这个问题。
\end{remark}

\begin{remark}
粗略地说，论证如下。在 §3 中，我们讨论 Eisenstein 级数在算术朗兰兹的谱侧对应什么。然后在 §4 中，我们在平凡局部系统附近给出 $\LS^{\arthm}_{\bG}$ 的局部坐标（见定理 4.3.3.1），并随后在这些坐标中显式计算谱 Eisenstein 级数（见定理 4.7.2.1）。由此，结果基本上是显然的（见引理 3.7.1.1）。
\end{remark}

\begin{remark}
我们指出，定理 4.7.2.1 证明的一个关键点暗示了某些发散级数的操作与范畴迹方法之间的关系。我们在 §4.6.5 中把关于这个主题的想法——就其目前的形式而言——写出来。这部分材料可以基本独立于本文其余部分阅读。
\end{remark}

\subsection{一些评论}

\pnumt{动机 I。}{}\ $\Aut^{\unr}_{\cusp,[\triv]}$ 的消失是 Beuzart-Plessis--Harris--Thorne 即将开展的工作中的一个组成部分，该工作通过迹公式研究函数域的局部 Langlands 对应。上述定理使他们的结果以几何朗兰兹对应为条件，而后者近年来已取得很大进展。

\pnumt{动机 II。}{}\ 对 $G = PGL_n$，任意尖形式都有不可约的 Langlands 参数；也就是说，在此情形 $\Aut^{\unr}_{\cusp,[\sigma]} = 0$，除非 $\sigma$ 不可约；关于最近的处理，我们参考 [22] 引理 16.4（它沿用了 [20] 的部分内容）；特别参见 [20] 定理 VI.9 的陈述。

然而，对一般的 $G$，情形更复杂：尖形式可以有可约的 Langlands 参数；这与这些尖形式不满足 Ramanujan 猜想有关。

尽管如此，Arthur 的猜想对可能出现的 $\sigma$ 给出了一些限制。首先注意，存在一个典范映射 $\mathcal{W}_X \to \ZZ \xrightarrow{1 \mapsto q} \e^\times$，我们记之为 $\gamma \mapsto |\gamma|$；选取 $\sqrt{q} \in \e^\times$，我们于是得到一个典范映射
\begin{equation}
  \begin{aligned}
    \mathcal{W}_X &\longrightarrow \mathcal{W}_X \times SL_2(\e),\\
    \gamma &\longmapsto \Bigl(\gamma, \begin{pmatrix} \sqrt{|\gamma|} & 0 \\ 0 & 1/\sqrt{|\gamma|} \end{pmatrix}\Bigr).
  \end{aligned}
\end{equation}
Arthur 的猜想预测：$\Aut^{\unr}_{\cusp,[\sigma]}$ 将为零，除非 $\sigma$ 沿 (1.2) 延拓为 $\mathcal{W}_X \times SL_2$ 到 $\bG$ 的一个不可约表示。\footnote{形式上，这意味着我们在 $\e$ 上有一个映射 $SL_2 \to \bG$，以及一个连续映射 $\mathcal{W}_X \to \bG(\e)$，其像与 $SL_2(\e) \to \bG(\e)$ 的像交换。}

\pnum{}假设给定这样一个 Arthur 参数 $\sigma^\sharp: \mathcal{W}_X \times SL_2 \to \bG$，并假设它到 $\pi_1^{\et}(X)$ 的限制是平凡的，于是我们有一个映射 $\sigma^\sharp: \ZZ \times SL_2 \to \bG$。令 $f \in \bg(\e)$ 表示 $\sigma^\sharp\bigl(0, \bigl(\begin{smallmatrix} 1 & 0 \\ 0 & 1 \end{smallmatrix}\bigr)\bigr)$ 的对数，并令 $F \in \bG(\e)$ 表示 $\bigl(1, \bigl(\begin{smallmatrix} \sqrt{q} & 0 \\ 0 & 1/\sqrt{q} \end{smallmatrix}\bigr)\bigr)$ 的像。注意 $F$ 编码了 $\sigma^\sharp$ 的底层 Langlands 参数，并注意 $\Ad_F(f) = qf$。我们还注意，$f$ 必须非零，否则 $\sigma^\sharp$ 不会是不可约的。

于是我们得到：

\begin{conj}
假设 $F \in \bG(\e)$ 是半单元素，并令 $\sigma_F: \mathcal{W}_X \to \bG(\e)$ 表示相应的 Langlands 参数 $\mathcal{W}_X \to \ZZ \xrightarrow{1 \mapsto F} \bG(\e)$。那么 $\Aut^{\unr}_{\cusp,[F]}$ 平凡，除非 $q$ 是 $\Ad_F: \bg \to \bg$ 的一个特征值。
\end{conj}

\begin{remark}
当然，在这个猜想中 $q$ 可以替换为 $q^{-1}$（对 Arthur 的 $SL_2$ 作用 Cartan 对合），这部分反映了该猜想在我们规范化方式的各种修改下（例如几何 Frobenius 与算术 Frobenius）的不变性。
\end{remark}

\pnum{}假定几何朗兰兹，我们将证明与这个猜想相近的结果。确切地说，我们将证明 $\Aut^{\unr}_{\cusp,[\sigma_F]} = 0$，除非 $\Ad_F$ 有一个特征值等于 $H^1_{\et}(X) \times H^2_{\et}(X)$ 中出现的某个 Frobenius 特征值；由于 $q$ 是 $H^2_{\et}(X)$ 上的 Frobenius 特征值，这比猜想的假设略强一些。我们指出，$F$ 为单位元的情形给出定理 A。

\subsection{约定与记号}

我们的希望是，本文能成为进入长论文 [2]、[3]、[4] 的一个入口。我们在 §2 中提供了一些背景。我们也请读者参考 [2] 的引言以获取更多背景。我们还力求在每次使用这些论文中的技术结果时给出精确引用，希望这能帮助读者在这些工作中导航。

我们一般遵循 [2] 的约定与记号；特别参考 §0.9。我们使用\emph{几何域} $k = \Fb$ 来处理曲线 $X$ 的几何、$H$-丛的模叠 $\Bun_H$ 等等，并使用特征 $0$ 的\emph{系数域} $\e = \overline{\Q}_\ell$ 来处理局部系统的几何。在 §2 之外，我们取 $k = \Fb$。$k$ 上的几何是经典代数几何，而 $\e$ 上的代数几何是导出的。我们使用高阶范畴方法。我们的 DG 范畴假定是 $\e$-向量空间上的充实范畴。

对一个在 $k$ 上局部有限型的代数叠 $\mathcal{Y}$，我们令 $\Shv(\mathcal{Y})$ 表示 $\mathcal{Y}$ 上 $\e$-层的 DG 范畴（见 [15] A.1.1 (d$'$)）。我们令 $\qLisse(\mathcal{Y}) \subseteq \Shv(\mathcal{Y})$ 表示\emph{拟 lisse} 复形的子范畴，如 [2] 定义 1.2.6 所述；这些对象的（比如说反常）上同调是 lisse 层的余极限（按通常意义）。

我们一般把 $\infty$-范畴简称为\emph{范畴}，以简化术语。

我们令 $\DGCat_{\mathrm{cont}}$ 表示在连续 DG 函子下余完备（且可及）的 DG 范畴的范畴。我们把 $\DGCat_{\mathrm{cont}}$ 视为配有 Lurie 张量积。我们令 $\Vect \in \DGCat_{\mathrm{cont}}$ 表示 $\e$-向量空间的 DG 范畴，它是幺半结构的单位。

对一个 DG 范畴 $\mathcal{C}$，我们令 $\mathcal{C}^c$ 表示其紧对象子范畴。当 $\mathcal{C}$ 有 $t$-结构时，我们令 $\mathcal{C}^\heartsuit$（相应地 $\mathcal{C}^{\leq 0}$、相应地 $\mathcal{C}^{\geq 0}$）表示 $t$-结构的核心（相应地连通对象子范畴、相应地余连通对象子范畴）。

关于范畴迹方法的背景，我们参考 [15]。

对 $\e$ 上的仿射代数群 $H$，我们提醒：$\LS^{\restr}_H = \LS^{\restr}_H(X)$ 是 $X$ 上（具有受限变差的）$H$-局部系统的模叠。回忆一个 $H$-局部系统不过是一个 $t$-正合（等价地：右 $t$-正合）对称幺半函子 $\Rep(H) \to \qLisse(X)$（等价地：一个对称幺半 $\e$-线性函子 $\Rep(H)^{\heartsuit,c} \to \mathrm{Lisse}(X)^{\heartsuit}$）。因此，我们定义叠 $\LS^{\restr}_H$ 来参数化右 $t$-正合对称幺半 DG 函子 $\Rep(H) \to \qLisse(X)$；更确切地说，$\LS^{\restr}_H$ 的 $S = \Spec(A)$ 点是右 $t$-正合对称幺半函子 $\Rep(H) \to A\text{-}\mathrm{mod}(\qLisse(X))$ 的群胚。

沿几何 Frobenius $\Frob_X: X \to X$ 的拉回定义了一个映射 $\LS^{\restr}_H \to \LS^{\restr}_H$，我们也称之为 \emph{Frobenius}。\footnote{注意这个映射是 $\e$ 上叠的映射。$\LS^{\restr}_H$ 的 Frobenius 可以看作 Frobenius 作用在 $X$ 的（$\e$-向量空间）\et{} 上同调上的非交换/非线性版本。}它的 Frobenius 不动点按定义是叠 $\LS^{\arthm}_H$，它（同义反复地）参数化从 $\Rep(H)$ 到 $X$ 上拟 lisse Weil 层的右 $t$-正合对称幺半函子。

最后，我们始终假定 $p = \mathrm{char}(k)$ 满足 [2] §14.4.1 的（温和）假设。

\subsection{致谢}

我感谢 Michael Harris 提出这个问题。我也感谢 Sasha Braverman 与 Dennis Gaitsgory 对这个主题的兴趣以及有益的谈话。我感谢 Dima Arinkin、Dennis Gaitsgory、David Kazhdan、Nick Rozenblyum 与 Yasha Varshavsky 在这个主题上的合作，感谢 Vincent Lafforgue 与 Cong Xue 的相关讨论。

最后，我感谢 2022 年 IHÉS Langlands 纲领暑期学校的组织者邀请我作报告，并在本文写作期间给予我非凡的耐心。

% ============================================================
%  ch02a：§2 引言、2.1 算术与几何朗兰兹（原文 pp.5–9）
% ============================================================
\section{AGKRRV 理论}

我们首先概述 [2] 与 [4] 的工作以及一些背景材料。

这些工作诚然是技术性的。我们不打算在这里概述这些工作所需的每一项技术背景。不过，我们尽量解释某些技术问题为何出现（例如，为何需要导出代数几何）。但在这方面，我们自由地诉诸本主题中一些基础性的想法，它们可能并非所有读者都熟悉：叠、DG 范畴、$D$-模、$\ell$-进层与 $\IndCoh$ 都属此类。尽管这些主题技术性强且并不总是广为人知，但如今已有许多参考文献（以及乐于助人的专家），我们认为有兴趣的读者应当不难找到资源，以探究讨论中出现的背景材料。

本节结构如下。首先，在 §2.1 中，我们稍微解释对自守函数感兴趣的读者应如何看待 de Rham（或 $D$-模）几何朗兰兹猜想，并强调后者相对于前者的若干可喜特征。在 §2.3 中，我们解释受限几何朗兰兹对应；其中一侧涉及 $\ell$-进层的模理论，我们在 §2.2 中解释。在 §2.4 中，我们解释当考虑 Frobenius 时，在有限域上故事如何展开。最后，在 §2.5 中，我们描述我们的主要算术结果（来自 [4]）如何被证明，并强调 Xue 关于 shtuka 上同调层的工作所起的关键作用。

\subsection{算术与几何朗兰兹}

\pnumt{算术朗兰兹。}{} 传统的算术朗兰兹关乎\emph{自守表示}，按定义，它们是出现在某个合适的自守函数空间中的某些\emph{不可约}表示。

朗兰兹哲学的一种粗糙（也许有些俗气）形式预言：$G$ 的自守表示对应于 $G$ 的朗兰兹参数。这里有一些修正，但它们并非我们的重点：应当使用 Arthur 参数，会出现 $L$-包，对一般数域并没有朗兰兹参数（或朗兰兹群）的合适定义，等等。我们的重点是该猜想的\emph{原子性}：基本对象是某个函数空间的不可约子商，而不是函数空间本身。

\pnumt{几何朗兰兹。}{} 与此相对，\emph{几何}朗兰兹猜想的传统形式预言：
\begin{equation}
  \Dmod(\Bun_G) \simeq \IndCoh_{\Nilp^{\mathrm{spec}}}(\LS^{\dR}_{\bG}) \simeq \QCoh(\LS^{\dR}_{\bG}).
\end{equation}
在上式中，$\IndCoh_{\Nilp^{\mathrm{spec}}}$ 是 $\QCoh$ 的一个合适扩张，由 Arinkin--Gaitsgory 定义，我们将在 §2.1.11 中稍作讨论。

由于本次会议主要关心朗兰兹纲领的\emph{算术}方面，我们花一些时间解释\emph{几何}设定的若干起始特征，包括上面使用的记号以及思考那里出现的主要对象的一些方式。

几何朗兰兹猜想的这种形式归功于 Beilinson--Drinfeld 与 Arinkin--Gaitsgory，参见 [1] 与 [13] 中对这一思想圈的导引。我们将把它称为几何朗兰兹猜想的 \emph{de Rham} 形式，因为 $D$-模理论记忆的是 de Rham 上同调群。

\pnum{}几何朗兰兹猜想的输入是一条固定域 $k$ 上的光滑射影曲线 $X$。为简化某些要点，我们假定 $k$ 代数闭，尽管在 de Rham 设定中这并非本质性的。

于是 $\Bun_G$ 是 $X$ 上 $G$-丛的空间。更确切地说，$\Bun_G = \Bun_G(X)$ 是一个\emph{叠}，其点函子由
\[
  \Bun_G(S) := \Hom(X \times S, \mathbf{B}G) = \{X \times S \text{ 上的 } G\text{-丛}\}
\]
给出，其中 $S$ 是仿射概形，$\mathbf{B}G$ 是 $G$ 的分类叠。众所周知，$\Bun_G$ 是局部有限型的光滑代数叠，尽管它不是拟紧的。

Weil 的一个漂亮公式\footnote{这个公式基本上是显然的，只要知道 $G$-丛在光滑射影曲线上是 Zariski（而不只是 étale）局部平凡的。对 $GL_n$，这由下降得到。对其他群，这是 Steinberg 定理。我们在此指出，该定理在有限域上也成立，这一点在主题中常被默认。对单连通 $G$，这是 Harder [19] 的定理。一般地，取满射 $G'\to G$，其中 $G'$ 的导群单连通、$\operatorname{Ker}(G'\to G)$ 是（连通）环面；于是 $X$ 上的 $G$-丛可提升到 $G'$（由类域论的 Brauer 群考虑），从而化归为 Harder 定理。} 说：
\begin{equation}
  \Bun_G(k) = G(F)\backslash G(\A)/G(\OO)
\end{equation}
其中 $\OO = \prod_{x \in X(k)} \OO_x$ 是\emph{整 adele} 环（$\OO_x$ 是以 $x \in X(k)$ 为基点的 Taylor 级数环），$\A = \colim_{S \subseteq X(k)\ \text{有限}} \bigl(\prod_{x \in S} \mathrm{Frac}(\OO_x) \times \prod_{x \notin S} \OO_x\bigr)$ 是类似定义的 adele 环，而 $F = k(X)$ 是 $X$ 上的有理函数域。因此，我们可以把 $\Bun_G$ 看作双重商空间的一个几何化身，无分歧自守函数就生活在那里（如果把 $k$ 换成有限域的话）。
\pnum{}几何表示论中的一个基础类比说：当 $k$ 特征为零时，叠 $\mathcal{Y}$ 上的 $D$-模范畴的行为类似于集合 $\mathcal{Y}(\F_q)$ 上的函数空间——如果这样的东西有意义的话。\footnote{例如，读者可以设想 $\mathcal{Y}$ 定义在 $\ZZ[1/N]$ 上，$N$ 与 $q$ 互素。但我要劝读者不要在这个点上过于字面化。}

这个想法有几个理由。首先，对 $k = \mathbb{C}$，某些 $D$-模通过 Riemann--Hilbert 对应与构造可构造层相关，而后者又通过 Riemann 存在定理与 \et{} 层相关；对 $k = \F_q$，后者又通过 Grothendieck--Deligne 的层—函数对应与函数相关。\footnote{我们参考 [28] 的前几节“Applications de la formule des traces aux sommes trigonométriques”以了解这个概念的背景。}

或者，人们可以设想，$D$-模编码线性微分方程组，其解在 $\mathcal{Y}(\mathbb{C})$ 上定义函数（或分布），而这些函数以不同理由类似于 $\mathcal{Y}(\F_q)$ 上的函数。

在实践中，在这个类比里使用 $\mathcal{Y}$ 上的\emph{全部} $D$-模很重要。例如，此设定中的 \emph{Mellin 变换}是一个等价
\[
  \Dmod(\mathbb{G}_m) \simeq \QCoh(\mathbb{A}^1/\ZZ);
\]
它可以看作几何朗兰兹型等式的简化玩具模型。在 Mellin 变换下，左端的全纯对象与正则全纯对象在右端都没有合理的描述。人们把这视为一个信号：在几何表示论中应当使用“全部”$D$-模的范畴，而不是某种可构造的子范畴。

此外，由 [24]，对 $G = \mathbb{G}_m$，等价 (2.1) 并不来自 Abel 范畴的等价；也就是说，在这个类比中必须使用\emph{导出}范畴。按照现代理解，我们使用 $\infty$-范畴的同伦形式中的 DG 范畴；我们一般遵守这样的约定：我们的 DG 范畴应当具有直和，而它们之间的函子应当是线性的、正合的，并保持直和。同伦形式的好处在于，它通过引入代数工具\emph{减轻}了主题的基础负担——在这种语言中，我们可以最自如地谈论幺半范畴、模范畴、张量积等等。

\begin{example}
按照前面的讨论，人们把 $\Dmod(\Bun_G)$ 看作无分歧自守函数空间 $\Aut^{\unr}_{\mathrm{c}}$ 的类比物。
\end{example}

\pnum{}我们再停留片刻，进一步讨论范畴与向量空间之间的类比。

其起源可以这样理解：对定义在 $\F_q$ 上的叠 $\mathcal{Y}/\F_q$，一个可构造 Weil 层 $\mathcal{F}$ 通过取 Frobenius 在有理点纤维上的迹，给出 $\mathcal{Y}(\F_q)$ 上的一个函数，这给出了从层产生函数的一个相当一般的程序。这就是层与函数之间类比的通常来源。

人们可以说，空间上的函数构成一个向量空间，而空间上的层构成一个范畴，于是（函数的）向量空间被范畴化（为层的范畴）。

\pnum{}我们可以把前面的讨论变得更精确。

从根本上说，前面讨论中函数的来源是：如果有一个（有限维）向量空间 $V$ 连同线性变换 $T: V \to V$，我们可以构造 $\mathrm{tr}_V(T)$ 而得到一个数。

类似地，对一个（可对偶化的 DG）范畴 $\mathcal{C}$ 连同自函子 $T: \mathcal{C} \to \mathcal{C}$，存在一个与这份数据相伴的迹 $\mathrm{tr}_{\mathcal{C}}(T) \in \Vect$；关于这个构造的详细讨论，我们参考 [15]。我们想强调：除了围绕范畴的（也许是严重的）心理障碍之外，这个一般构造相当形式化，并且与通常的迹理论相镜像。

\pnum{}我们现在转而更仔细地解释几何朗兰兹等价。

空间 $\LS^{\dR}_{\bG} = \LS^{\dR}_{\bG}(X)$ 是 $X$ 上 de Rham $\bG$-局部系统的模叠。在本文中，叠 $\LS^{\dR}_{H}$（对 $H/k$ 仿射代数群）按惯例定义为具有如下 $S$-点：
\[
  \LS^{\dR}_{H}(S) := \Hom(X_{\dR} \times S, \mathbf{B}H),
\]
其中 $X_{\dR}$ 是 $X$ 的 de Rham 空间。在这里详细讨论 de Rham 空间过于冗长，但要点是 $\QCoh(X_{\dR}) = \Dmod(X)$。

不那么循惯例地，也可以如下进行。首先，在 $k$-点的层面上：一个 $H$-局部系统应当是什么？我们可以取 $\qLisse^{\dR}(X) \subseteq \Dmod(X)$ 为其每个上同调都是局部系统（即带联络的向量丛）之余极限的对象组成的子范畴；这是 \emph{lisse} $D$-模的一个合适导出范畴，但我们按 [2] 的惯例称它们为\emph{拟 lisse}。

于是 de Rham $H$-局部系统本质上是一个对称幺半函子 $\Rep(H) \to \qLisse^{\dR}(X)$；对 $H = \mathbb{G}_m$ 这还不太对，因为此时这样一份数据是 $\qLisse^{\dR}(X)$ 中的一个张量可逆对象 $\sigma$，即一个\emph{上同调平移过的}带联络线丛；为消除这种歧义，我们细化定义，要求我们的函子是（在作一次上同调平移 $[1]$ 之后）右 $t$-正合的（以照顾到该函子把平凡表示送到常层 $\mathbf{e}_X$ 这一事实，后者位于次数 $\dim(X) = 1$）。\footnote{当然，这个问题之所以出现，只是因为我们坚持使用导出范畴——读者可以对此提出异议。在 §2.1.12 中，我们解释在我们所讲的故事中必须使用导出代数几何，因此我们的仿射测试概形 $S$ 也应当是导出的；导出范畴 $\QCoh(S)$ 对导出几何敏感，而 Abel 范畴 $\QCoh(S)^\heartsuit$ 则不然。

鉴于这些问题，我们作了一个教学上的选择：坚持使用导出范畴与（相差平移的）右 $t$-正合函子。}\footnote{读者可能会问：右 $t$-正合函子 $\Rep(H) \to \qLisse^{\dR}(X)$ 也是 $t$-正合的（两种情形都相差平移），那为什么还要写“右”？原因是，我们很快会在 (2.3) 中推广这一设定。譬如说那里的 $H$ 是平凡群，于是我们谈论的是对称幺半函子 $\Vect \to \QCoh(S)$，它把 $k$ 送到 $\mathbf{e}_X \boxtimes \OO_S$，其中 $\mathbf{e}_X$ 是常层。在平移 $\dim(X)$ 之后，后者总是连通的（即位于上同调次数 $\leq 0$），但只有当 $S$ 是经典概形（即“通常”代数几何的对象，而非导出代数几何的对象）时，它才位于 $t$-结构的核心中。}

我们注意到，这个定义的行为本质上符合预期：一个 de Rham $GL_n$-局部系统是 $X$ 上秩 $n$ 的带联络向量丛；一个 de Rham $SO_n$-局部系统是秩 $n$ 的带联络向量丛 $\mathcal{E}$ 连同保持联络的非退化对称配对 $\mathcal{E} \otimes_{\OO_X} \mathcal{E} \to \OO_X$；de Rham $Sp_{2n}$-局部系统类似，但非退化配对是反对称的；de Rham $G_2$-局部系统是带联络的 octonion 丛；等等。

这种局部系统的 Tannaka 定义——这也许是对一般代数群定义局部系统的最简单方式——经适当调整可给出 $\LS^{\dR}_{H}$ 的 $S$-点：即对称幺半函子
\begin{equation}
  \Rep(H) \to \Dmod(X) \otimes \QCoh(S),
\end{equation}
它们在相差 $\dim(X) = 1$ 的平移的意义下是右 $t$-正合的。

我们建议读者参考 [2] §4.1 以了解进一步的相关讨论。我们简要指出，任何对称幺半函子 $\Rep(H) \to \Dmod(X)$ 都落在 $\qLisse^{\dR}(X)$ 中，尽管 (2.3) 一般不会映入 $\qLisse^{\dR}(X) \otimes \QCoh(S)$。

\pnum{}不必多说：对数论学家而言，$\LS^{\dR}_{\bG}$ 被看作 Galois 参数的模叠。

确实，由 Riemann--Hilbert 对应（的一个很容易的形式），对 $k = \mathbb{C}$，存在 $\bG$-局部系统与同态 $\rho: \pi_1(X(\mathbb{C})) \to \bG(\mathbb{C})$（模共轭）之间的解析等同（尽管这在 $S$-族中不能天真地成立）。

\pnum{}几何朗兰兹猜想中下一块关键结构是\emph{谱作用}。

这是幺半范畴 $\QCoh(\LS^{\dR}_{\bG})$ 在 $\Dmod(\Bun_G)$ 上的作用，由 Drinfeld--Gaitsgory 构造。根据\emph{同上}，这个作用由其与 $\Dmod(\Bun_G)$ 上 Hecke 作用（的一个足够强的版本）的相容性唯一刻画。

在此我们参考 [?] §1.1 中关于唯一性的讨论，以及 [?] §1.5 中与 Hecke 函子的确切联系。% [Tutor Note: 原文两处均作 [?]，据上下文应为 [2]。]

\pnum{}谱作用可以如下可视化。

令 $\mathcal{Y}$ 为任意叠，并设 $\mathcal{C}$ 是 $\QCoh(\mathcal{Y})$ 的一个模范畴。我们把它画成 $\mathcal{Y}$ 上的一个纤维化范畴：
\begin{equation*}
\begin{tikzcd}[column sep=large, row sep=large]
  \mathcal{C}_y \arrow[r] \arrow[d] & \mathcal{C} \arrow[d] \\
  y \arrow[r, hook] & \mathcal{Y}
\end{tikzcd}
\end{equation*}
这里 $y$ 处的纤维 $\mathcal{C}_y$ 定义为：
\[
  \mathcal{C}_y := \mathcal{C} \underset{\QCoh(\mathcal{Y})}{\otimes} \Vect,
\]
如果 $y$ 是 $k$-点；如果它是 $A$-点，则把 $\Vect$ 换成 $A$-$\mathrm{mod}$。启发式地，我们可以写 $\mathcal{C} = \int_{y \in \mathcal{Y}} \mathcal{C}_y \, dy$。这个形式体系在 [14] 中有详细研究。

于是，非形式地说，Drinfeld--Gaitsgory 谱作用说的是：范畴 $\Dmod(\Bun_G)$ 纤维化于 $\LS^{\dR}_{\bG}$ 之上，并且这一结构由 Hecke 函子典范地决定。因此，谱作用的\emph{存在性}可以解释为自守层范畴（这个说法在此 de Rham 语境中指的是 $\Dmod(\Bun_G)$）的（范畴式）\emph{互反律}。

按定义，$\Dmod(\Bun_G)$ 在点 $\sigma \in \LS^{\dR}_{\bG}$ 处的纤维 $\Dmod(\Bun_G)_\sigma$ 是特征值为 $\sigma$ 的 \emph{Hecke 特征层}范畴。

\pnum{}在上面启发式的公式 $\mathcal{C} = \int_{y \in \mathcal{Y}} \mathcal{C}_y \, dy$ 中，我们设想在 $\mathcal{Y}$ 上有一个范畴值测度 $\mathcal{C}_y \, dy$。在我们关心的情况中，我们可能希望计算它。

这正是完整几何朗兰兹猜想的工作。$\QCoh(\LS^{\dR}_{\bG})$-模范畴 $\IndCoh_{\Nilp^{\mathrm{spec}}}(\LS^{\dR}_{\bG})$ 在这个隐喻中编码 Plancherel 测度的类比物。

这里，具有幂零奇异支撑的\emph{非相干}层范畴 $\IndCoh_{\Nilp^{\mathrm{spec}}}(\LS^{\dR}_{\bG})$ 在 [1] 中被定义，此后一直是该领域广泛研究的对象。它是 $\QCoh$ 的一个修改，具有反映 $\LS^{\dR}_{\bG}$ 的奇点某些信息的几何性质。关于这个主题的导引，我们参考 [1]。我们用 $\Nilp^{\mathrm{spec}} \subseteq T^*[-1]\LS^{\dR}_{\bG}$ 表示\emph{谱幂零锥}，并提醒它在许多其他文献中常简记为 $\Nilp$。

由于不可约 $\bG$-局部系统不支持其伴随丛的非零幂零水平截影，我们有：
\[
  \QCoh(\LS^{\dR,\mathrm{irred}}_{\bG}) \underset{\QCoh(\LS^{\dR}_{\bG})}{\otimes} \IndCoh_{\Nilp^{\mathrm{spec}}}(\LS^{\dR}_{\bG}) = \QCoh(\LS^{\dR,\mathrm{irred}}_{\bG})
\]
（详见 [1] 命题 13.3.3）。在我们的类比下，这意味着 Plancherel 测度在 $\LS^{\dR,\mathrm{irred}}_{\bG}$ 上为常数，取值为 $\Vect$。在可约局部系统附近，存在与伴随丛的幂零水平截影相关的修正，它们在这里是 Arthur 的 $SL_2$ 的化身。

\pnum{}我们现在给出一个简单例子，以说明一些基本技术要点。

假设 $X = \mathbb{P}^1$ 且 $G = \mathbb{G}_m$。于是 $\Bun_{\mathbb{G}_m}$ 参数化 $\mathbb{P}^1$ 上的线丛，因此同构于 $\ZZ \times \mathbf{B}\mathbb{G}_m$：$\ZZ$-因子参数化线丛的次数，而 $\mathbf{B}\mathbb{G}_m$-因子编码了 $\mathbb{P}^1$ 上每个线丛都有自同构群 $\mathbb{G}_m$ 这一事实（在 $S$-族中以恰当方式理解）。

因此，
\[
  \Dmod(\Bun_{\mathbb{G}_m}(\mathbb{P}^1)) = \prod_{n \in \ZZ} \Dmod(\mathbf{B}\mathbb{G}_m).
\]
范畴 $\Dmod(\mathbf{B}\mathbb{G}_m)$ 可以相当显式地算出。令 $\pi: \Spec(k) \to \mathbf{B}\mathbb{G}_m$ 为结构映射，我们提醒它是一个光滑覆盖。函子 $\pi^!: \Dmod(\mathbf{B}\mathbb{G}_m) \to \Dmod(\Spec(k)) = \Vect$ 显然保守，并具有左伴随 $\pi_!$。由基变换，$\Vect$ 的 endofunctor $\pi^! \pi_!$ 由与 $C_{\dR}(\mathbb{G}_m)$（$\mathbb{G}_m$ 的 de Rham 上同调）作张量给出。此外，由 Barr--Beck 的一个简单形式，这个 endofunctor $\pi^! \pi_!$ 具有自然的单子结构，对应于 $C_{\dR}(\mathbb{G}_m)$ 上来自群结构 $\mathbb{G}_m$ 的代数结构；并且，诱导函子 $\Dmod(\mathbf{B}\mathbb{G}_m) \to C_{\dR}(\mathbb{G}_m)$-$\mathrm{mod}$ 是一个等价。

最后，当然，$C_{\dR}(\mathbb{G}_m)$ 是 $k$ 的一个平方零扩张，由一个位于上同调次数 $-1$ 的单个生成元 $\eta$ 扩张而成（别名：一个次数 $-1$ 生成元上的对称代数；一个生成元上的同调分次外代数）——这不过反映了圆周的同调在次数 $0$ 与 $1$ 上是 $1$ 维这一初等事实。

于是我们有：
\begin{equation}
  \Dmod(\Bun_{\mathbb{G}_m}(\mathbb{P}^1)) = \prod_{n \in \ZZ} (k \oplus k\eta)\text{-}\mathrm{mod}.
\end{equation}

天真地看，$\mathbb{P}^1$ 是单连通的，因此没有非平凡局部系统，所以人们可能期望 $\LS^{\dR}_{\bG}(\mathbb{P}^1)$ 等于 $\mathbf{B}G$（反映平凡局部系统的非平凡自同构群）。对 $\bG = \mathbb{G}_m$，我们将有 $\QCoh(\mathbf{B}\mathbb{G}_m) = \Rep(\mathbb{G}_m) = \prod_{n \in \ZZ} \Vect$，这与 (2.4) 很接近，但缺少次数 $-1$ 的生成元。

事实上，这是因为我们太天真了。\emph{导出}叠 $\LS^{\dR}_{\bG}(\mathbb{P}^1)$ 等于\footnote{这个公式例如来自把 $\mathbb{P}^1$ 上的局部系统看作两个标准开集 $\mathbb{A}^1 \subset \mathbb{P}^1$ 上的局部系统连同它们在交 $\mathbb{A}^1 \setminus 0$ 上的同构。注意 $\mathbb{A}^1$ 是可缩的，而不只是单连通的。} $\mathbf{B}\bG \times_{\bG} \mathbf{B}\bG$，对 $\bG = \mathbb{G}_m$ 就是 $\mathbf{B}\mathbb{G}_m \times (0 \times_{\mathbb{A}^1} 0)$。这里的纤维积必须在\emph{导出代数几何}意义下取。于是我们得到 $\QCoh(\LS^{\dR}_{\mathbb{G}_m}(\mathbb{P}^1)) = \prod_{n \in \ZZ} \QCoh(0 \times_{\mathbb{A}^1} 0)$。最后，我们注意 $0 \times_{\mathbb{A}^1} 0$ 是 $\Spec$ 的 $k \otimes_{k[t]} k$（张量积是导出的，即包含 $\mathrm{Tor}^{k[t]}_i(k,k)$ 这些群），这恰是我们先前看到的同一个平方零扩张 $C_{\dR}(\mathbb{G}_m)$。

我们注意到，底层的经典叠恢复了我们先前的朴素局部系统概念。

或者，人们可以如下看待导出代数几何的用处。对一般 $X$ 与 $\bG$，标准论证说 $\LS^{\dR}_{\bG}$ 在 $\bG$-局部系统 $\sigma$ 处的切空间是 $H^1_{\dR}(X, \check{\mathfrak{g}}_\sigma)$，即以伴随局部系统 $\sigma$ 为系数的第一个 de Rham 上同调。更一般地，我们应当期望切\emph{复形}是 $C_{\dR}(X, \check{\mathfrak{g}}_\sigma)[1]$。正如上面的例子所示，这个公式只有在把 $\LS^{\dR}_{\bG}$ 解释为\emph{导出}叠时一般才能成立。

总结：我们在几何朗兰兹的谱侧使用导出代数几何，\emph{因为它给出正确的答案}（不像经典代数几何），并且\emph{因为它产生模空间更易处理的无穷小几何}。

\pnumt{结论。}{} 上面我们简要讨论了算术朗兰兹，并对（de Rham）几何朗兰兹作了较长的介绍。

有一个关键区别，§2.1.1 已有所暗示：在算术朗兰兹中，我们研究\emph{原子对象}（不可约表示），而在\emph{几何}朗兰兹中，我们研究分子对象（自守函数空间的一个类比物）。人们可以把这种情形与圆周 $S^1$ 上的 Fourier 理论作比较：\emph{原子}理论说 $S^1$ 的（必然酉的）特征标与 $\ZZ$ 一一对应，但实际的 Fourier 理论说 $L^2(S^1)$ 是 $\ZZ$ 上 $1$ 维 Hilbert 空间（即 $L^2(\ZZ)$）的直接积分。在自守理论中，后者的类比物将是理想的，但我们所能触及的只是前者。

[2] 的一个出发点是试图解决这一差异，至少对函数域，至少在处处无分歧情形。最终，我们得到一个更接近几何朗兰兹猜想的算术视角。

我想强调：我们的工作并非唯一试图弥合这一鸿沟的工作；[32] 与 [11] 是密切相关的努力，而我们都受到 V.~Lafforgue 的突破 [22] 的启发。

% ============================================================
%  ch02b：§2.2 受限变差的局部系统（原文 pp.9–12）
% ============================================================
\subsection{受限变差的局部系统：导引}

\pnumt{问题何在？}{} 现在假设 $k$ 是代数闭域，$X$ 是 $k$ 上的光滑射影曲线。我们令 $\e := \overline{\Q}_\ell$。

我们希望模仿一般的几何朗兰兹故事，但把局部系统理解为 lisse \et{} $\e$-层，而不是带联络的向量丛。会出什么问题？

首先，让我们尽可能乐观：我们希望有一个叠 $\LS^{\et}_{\bG} = \LS^{\et}_{\bG}(X)$，其行为像先前的叠 $\LS^{\dR}_{\bG}$。在恰当的理解下，它的点应当是 $\bG$-局部系统，即右 $t$-正合对称幺半函子 $\Rep(\bG) \to \qLisse(X)$——这里 $\qLisse(X) \subseteq \Shv(X) = \Shv^{\et}(X)$ 按 \et{} 意义理解，如 §1.4 那样。

作为预备步骤，注意 $X$ 上平凡 $\bG$-局部系统的自同构群是 $\bG(\e)$。这暗示 $\LS^{\et}_{\bG}$ 应定义在域 $\e$ 上，并包含一个对应于平凡局部系统的 $\mathbf{B}\bG$ 的拷贝。

\begin{remark}
这里我们看到几何朗兰兹的代数几何中的一个基本分叉：某些对象，如 $X$、$G$、$\Bun_G$ 等，生活在\emph{基域} $k$ 上，而谱对象，如 $\bG$、$\LS^{\et}_{\bG}$ 等，生活在\emph{系数域} $\e$ 上。对 de Rham 理论，系数域就是基域，这一区分可以忽略。
\end{remark}

\pnum{}现在假设 $k$ 特征为零，且 $X$ 的亏格 $g > 0$。

忽略技术问题（叠性、导出结构），我们可能首先猜测 $\LS^{\et}_{\mathbb{G}_m}$ 是域 $\e$ 上类似 $\mathbb{G}_m^{2g}$ 的东西。毕竟，$X$ 的 \et{} 基本群的 Abel 化是 $\widehat{\ZZ}^{2g}$。此外，人们可以看出（忽略同样的技术问题），叠 $\LS^{\dR}_{\mathbb{G}_m}$ 在 $\mathbb{C}$ 上\emph{复解析地}同构于 $\mathbb{G}_m^{2g} \times (\mathbf{B}\mathbb{G}_m \times 0 \times_{\mathbb{A}^1} 0)$。

然而，$\ZZ$ 与 $\widehat{\ZZ}$ 之间的区别在这里是关键。事实上，连续同态 $\widehat{\ZZ}^{2g} \to \e^\times$ 由 $(\OO^{\times}_{\e})^{2g}$ 中的点参数化，其中 $\OO_{\e} \subseteq \e$ 是整元素的通常赋值子环。换言之，我们希望 $\LS^{\et}_{\mathbb{G}_m} = \mathbb{G}_m^{2g} \times (\mathbf{B}\mathbb{G}_m \times 0 \times_{\mathbb{A}^1} 0) \approx \mathbb{G}_m^{2g}$ 太天真了：右端在 $\e$ 上点太多了！

注意，很难找到一个有趣的\footnote{这里有一个并不有趣的例子：$\prod_{\OO^\times_{\e}} \Spec(\e)$。对 $\mathbb{G}_m$ 的情形，$\LS^{\et}_{\mathbb{G}_m}$ 将在这个答案的基础上，于叠性与（可能的）上同调幂零元方面有所不同（当然，还要把 $\OO^\times_{\e}$ 换成 $(\OO^\times_{\e})^{2g}$）。}、以 $\e$ 为系数域且具有 $\e$-点 $\OO^{\times}_{\e}$ 的（譬如说连通且代数的）叠 $\LS^{\et}_{\mathbb{G}_m}$。所以我们放弃。

\pnumt{另一方面，}{} \et{} 局部系统（别名：Galois 表示）的形变理论是一个老故事。通常人们考虑挠系数，但我们在这里不必如此。要点是：对 \et{} 局部系统 $\sigma$，我们有一个 DG 李代数 $C_{\et}(X, \check{\mathfrak{g}}_\sigma)$，因此有一个相伴的形式模问题（参见 [25] 第 13 章与 [18]）。

换言之，尽管我们放弃了 $\LS^{\et}_{\bG}$ 的存在性，我们\emph{确实}知道它的 $\e$-点（毕竟那都是局部系统），并且\emph{确实}知道它在每个这样的点处的形式完备化。

所以目前，为了构造 $\LS^{\et}_{\bG}$ 的\emph{某个}近似，我们可以取“$\LS^{\et}_{\bG}$ 在每个 $\e$-点 $\sigma$ 处的形式完备化”的（通常是不可数的）无交并。这对 $\mathbb{G}_m$ 给出了正确答案，但对其他群我们将构造一个稍好一些的东西（正如 §2.2 余下部分所讨论的）。

\pnumt{定义。}{} 在 [2] 中，我们对任意仿射代数群 $H/\e$ 定义了 $\e$ 上的一个\emph{预叠}（即从连通交换 $\e$-代数到 $\infty$-群胚的函子，也就是模问题），称为 $\LS^{\restr}_{H}$。一般地，它比点的形式邻域记忆得稍多一点，正如我们将看到的。

这个定义是对 (2.3) 的天真模仿；按定义，$\LS^{\restr}_{H}$ 的一个 $S$-点是右 $t$-正合对称幺半函子：
\[
  \Rep(H) \to \qLisse(X) \otimes \QCoh(S).
\]

\begin{remark}
为使这个定义显得更具体，我们解释一下不用 DG 范畴张量积时右端是什么。假设 $\mathcal{C}$ 是一个 DG 范畴，我们提醒它有所有余极限。假设 $S = \Spec(A)$。于是 $\mathcal{C} \otimes \QCoh(S) = \mathcal{C} \otimes A\text{-}\mathrm{mod} = A\text{-}\mathrm{mod}(\mathcal{C})$，即 $\mathcal{C}$ 中带 $A$-作用的对象。所以右端相当具体——其复杂程度与 $A$-$\mathrm{mod}$ 大致相同。
\end{remark}

\pnumt{我们得到了什么？}{} 上面是一个形式定义。仍需论证我们给出了一个\emph{好的概念}，而“好”的含义将随着我们前进而变得更精细。

\pnumt{例：加性群。}{} 首先假设 $H = \mathbb{G}_a$。

我们断言，在这种情形中 $\LS^{\restr}_{\mathbb{G}_a}$ 实际上“看起来与” de Rham 情形相同！更确切地说，我们将证明 $\LS^{\restr}_{\mathbb{G}_a}$ 是对应于复形 $C_{\et}(X)[1]$ 的代数叠，即它（非典范地\footnote{我们实际上为 $C_{\et}(X)[1]$ 选择了一个形式性同构。}）同构于
\[
  \mathbf{B}H^0_{\et}(X) \times H^1_{\et}(X) \times \Omega H^2_{\et}(X)
\]
（其中 $\Omega H^2_{\et}(X)$ 定义为导出概形 $0 \times_{H^2_{\et}(X)} 0$）。人们可以看出（例如通过下面的分析），在 de Rham 设定中同样成立，只是处处用 de Rham 上同调替换 \et{} 上同调。

为看出这一点，令 $\triv \in \Rep(\mathbb{G}_a)$ 为平凡表示，即张量单位。存在一个典范映射 $\triv \to \triv[1] \in \Rep(\mathbb{G}_a)$，对应于 $\mathbb{G}_a$ 的标准 $2$ 维表示 $\begin{pmatrix} 1 & 1 \\ 0 & 1 \end{pmatrix}$（看作平凡表示的非平凡自扩张）。

现在，对对称幺半 DG 范畴 $\mathcal{C}$ 与对称幺半函子 $F: \Rep(\mathbb{G}_a) \to \mathcal{C}$，我们可以把 $F$ 作用到上面的扩张类上，得到一个映射 $\mathbf{1}_{\mathcal{C}} \to \mathbf{1}_{\mathcal{C}}[1]$，即 $(\infty\text{-})\text{群胚}$ $\Hom_{\mathcal{C}}(\mathbf{1}_{\mathcal{C}}, \mathbf{1}_{\mathcal{C}}[1]) = \Omega^{\infty-1}\underline{\mathrm{End}}_{\mathcal{C}}(\mathbf{1}_{\mathcal{C}})$ 中的一个点。\footnote{为向不熟悉这些材料的读者澄清：对谱 $V$（或 $k$-向量空间复形），$\Omega^{\infty}V$ 意为“取底层的 $\infty$-群胚”——至少对连通谱/链复形，这类似于取 Abel 群的底层集合；一般地，可以把它设想为“过渡到连通覆盖，然后取底层的同伦集合”。在显式的集合论模型中，我们可以取 $\ZZ$-模的链复形 $V^\bullet$，截断得到 $\tau^{\leq 0}(V^\bullet)$，然后在 Dold--Kan 下过渡到相应的单纯 Abel 群（从而单纯集合）。

记号 $\Omega^{\infty-1}(V)$ 只是指 $\Omega^{\infty}(V[1])$。} 容易看出\footnote{也就是说，只需利用 $\Rep(\mathbb{G}_a)$ 与交换代数 $\e \times \e[-1]$ 上的模范畴之间存在标准的对称幺半等价。}，这给出群胚的同构：
\[
  \Hom_{\mathrm{ComAlg}(\DGCat_{\mathrm{cont}})}(\Rep(\mathbb{G}_a), \mathcal{C}) \simeq \Hom_{\mathcal{C}}(\mathbf{1}_{\mathcal{C}}, \mathbf{1}_{\mathcal{C}}[1]).
\]
取 $\mathcal{C} = \qLisse(X) \otimes \QCoh(S)$，我们看到 $\LS^{\restr}_{\mathbb{G}_a}$ 的 $S$-点等于
\[
  \Omega^{\infty}\bigl(\underline{\Hom}_{\qLisse(X)}(\mathbf{e}_X, \mathbf{e}_X[1]) \otimes \underline{\Hom}_{\QCoh(S)}(\OO_X, \OO_X)\bigr) = \Omega^{\infty}\bigl(C_{\et}(X)[1] \otimes \Gamma(S, \OO_S)\bigr),
\]
其中 $\mathbf{e}_X$ 是 $X$ 上的常层。把形式体系展开之后，这就证明了断言。

\pnumt{例：乘性群。}{} 这里我们只陈述结果：

空间 $\LS^{\restr}_{\mathbb{G}_m}$ 是一个 ind-代数叠。它是其连通分支的无交并，而每个连通分支（同样非典范地\footnote{这里有一个更典范地构造该分支的方法。首先取 $\LS^{\restr}_{\mathbb{G}_a}$ 并在平凡 $\mathbb{G}_a$-局部系统处形式完备化。所得的叠接受来自 $\mathbf{B}\mathbb{G}_a^{\wedge} = \mathbf{B}\mathbb{G}_m$ 的同态（这里我们用了指数映射）；然后沿到 $\mathbf{B}\mathbb{G}_m$ 的映射作推出。}）同构于 $\mathbf{B}\mathbb{G}_m \times H^1_{\et}(X)^{\wedge}_0 \times \Omega H^2_{\et}(X)$。$\LS^{\restr}_{\mathbb{G}_m}$ 的连通分支与其 $\e$-点一一对应，我们提醒后者恰是 $X$ 上秩 $1$ 的 lisse 层。

\pnumt{工具箱是什么？}{} 这部分材料可以跳过。为方便读者，我们描述证明关于 $\LS^{\restr}_{H}$ 的事情的常用方法。

首先，我们需要探查底层的经典叠，忽略导出代数几何的问题。为此，令 $\Pi_X$ 为附着于 Tannaka 范畴 $\qLisse(X)^{\heartsuit}$ 的 Tannaka 群，于是 $\Pi_X$ 是 $\e$ 上的群概形，并具有对称幺半等价 $\Rep(\Pi_X)^{\heartsuit} \simeq \qLisse(X)^{\heartsuit}$。

对经典概形，$\LS^{\restr}_{H}$ 的 $S$-点典范地双射于群概形映射 $\Pi_X \times S \to H \times S$（模共轭）（其中我们在群胚意义下取商）——参见 [2] 命题 2.5.9（尽管该断言本质上只是 Tannaka 对偶加上簿记）。这允许我们使用代数群论的工具研究 $\LS^{\restr}_{H}$ 的底层经典预叠。

然后我们用形变理论推广到导出概形，这对 $\LS^{\restr}_{H}$ 来说容易计算；参见 [2] §2.2。

\begin{example}
我们用一个例子说明第一种技术。

先前关于 $\LS^{\restr}_{\mathbb{G}_m}$ 是 §2.2.7 中“胖点”的无交并的断言，等价于说：对任意代数闭域扩张 $\e'/\e$，映射 $S = \Spec(\e') \to \LS^{\restr}_{\mathbb{G}_m}$ 通过一个 $\e$-点分解。这变成关于群概形的一般断言：映射 $\Pi_X \times S \to \mathbb{G}_m \times S$ 来自一个定义在 $\e$ 上的映射。由于 $\mathbb{G}_m$ 具有有限型，这归约到同样的断言，但把 $\Pi_X$ 换成仿射代数群 $\Gamma$（即 $\Pi_X$ 的有限型商），我们甚至可以假设它是 Abel 的。这里该断言由交换代数群的表示论显然可见。
\end{example}

\pnumt{$\LS^{\restr}$ 的一般结构。}{} 我们希望下面的结果在此时不会令人惊讶。

首先，$\LS^{\restr}_{H}$ 总是一个\emph{形式}代数叠。更确切地说，若把 $\LS^{\restr}_{H} = \LS^{\restr}_{H}(X) \to \mathbf{B}H = \LS^{\restr}_{H}(\Spec(k))$ 取为在 $X$ 的一个 $k$-点处的纤维，则该映射在 indscheme 中可表，而且更好：在形式概形的无交并这些 indscheme 中可表——参见 [2] 定理 1.4.5。

其次，$\LS^{\restr}_{H}$ 的连通分支与\emph{半单} $H$-局部系统（模等价）一一对应。非形式地说，$\LS^{\restr}_{H}$ 的两个点属于同一连通分支，当且仅当它们的半单化在无穷小意义下接近。确切陈述参见 [2] 命题 3.7.2。

最后，如果我们设想 $\LS^{\et}_{H}$ 存在，那么对每个半单 $\sigma$，将存在局部系统的一个闭子叠 $\LS^{\et}_{H,\sigma}$，其半单化为 $\sigma$；此时 $\LS^{\restr}_{H}$ 在道德上就是 $\LS^{\et}_{H}$ 在每个这样的 $\LS^{\et}_{H,\sigma}$ 处形式完备化的无交并。关于 Betti 与 de Rham 设定中更精确的断言，参见 [2] §4。

% ============================================================
%  ch02c：§2.3 受限几何朗兰兹（原文 p.12）
% ============================================================
\subsection{受限几何朗兰兹}

\pnum{}令 $\Shv_{\Nilp}(\Bun_G) \subseteq \Shv(\Bun_G)$ 表示具有\emph{幂零锥中的奇异支撑}的层子范畴。这里 \et{} 层的奇异支撑由 Beilinson 在 [5] 中定义。

在 Betti 设定中，Ben-Zvi 与 Nadler [7] 说具有幂零奇异支撑的层是应当研究的正确对象（在 [2] 定理 18.1.6 中可以找到一个证明这一想法的精确定理）。我们把这一原则模仿到 \et{} 设定中，猜想：

\begin{conj}[受限几何朗兰兹猜想]
存在一个等价
\[
  \Shv_{\Nilp}(\Bun_G) \simeq \IndCoh_{\Nilp^{\mathrm{spec}}}(\LS^{\restr}_{\bG}).
\]
\end{conj}

这里右端如在 de Rham 情形中那样定义。

\begin{remark}
人们可以在 [23] 猜想 6.3.2 中找到这个猜想的一个简化版本。
\end{remark}

\begin{remark}
当然，猜想 2.3.1 受制于许多相容性。与 Whittaker 系数的相容性、（与 Eisenstein 级数的相容性的一个温和形式）\footnote{关于 Whittaker 相容性与 Eisenstein 相容性的表述，参见 [13]。} 以及 [10] 定理 8.3.0.1 的一个版本，\emph{唯一地}确定了猜想 2.3.1 中的比较函子；在 de Rham 与 Betti 设定中，这一想法是 [16] 的主题。
\end{remark}

\pnumt{证据。}{} 当几何域 $k$ 特征为零时，我们在 [2] 中证明受限 GLC 可由 de Rham 几何朗兰兹猜想推出。

一般地，人们可以直接对 $G = \mathbb{G}_m$ 验证该猜想；并且，用与 [21] 类似的想法，可以把该猜想归约到 $X = \mathbb{P}^1$ 的导出 Satake。


% ============================================================
%  ch02d：§2.4 Frobenius 与 §2.5 Xue 定理（原文 pp.12–16）
% ============================================================
\subsection{Frobenius}

\pnum{}我们现在讨论纳入 Frobenius 之后会发生什么。

现在假设基域 $k$ 是 $\Fb$。假设 $X$ 定义在 $\F_q$ 上；由于 $G$ 是先验定义在 $\ZZ$ 上的，可知 $\Bun_G$ 也自然地定义在 $\F_q$ 上。回忆有理结构可以由 $X$（相应地，$\Bun_G$）的几何 Frobenius 自同态来编码。

因此，存在 Frobenius 自同态（即沿几何 Frobenius 的拉回）作用在 $\Shv(X)$、$\qLisse(X)$ 与 $\Shv(\Bun_G)$ 上。

按 $\LS^{\restr}_{H}$ 的定义，$\qLisse(X)$ 的 Frobenius 自同态诱导 $\LS^{\restr}_{H}$ 的一个“Frobenius”自同态。

\begin{example}
假设 $H = \mathbb{G}_a$。由 §2.2.6，$\LS^{\restr}_{\mathbb{G}_a}$ 是链复形 $C_{\et}(X)[1]$ 的几何化身；这个复形带有其标准的 Frobenius 自同态，而在这本字典下两者自动匹配。

一般地，$\LS^{\restr}_{H}$ 上的 Frobenius 因此可以被看作 $C_{\et}(X)$ 上 Frobenius 的非线性类比。
\end{example}

\pnum{}我们把 $\LS^{\arthm}_{H}$ 定义为 $\LS^{\restr}_{H}$ 的 Frobenius 不动点。更确切地说，我们有形式导出叠的笛卡尔图：
\begin{equation*}
\begin{tikzcd}[row sep=large, column sep=large]
  \LS^{\arthm}_H \arrow[r] \arrow[d] & \LS^{\restr}_H \arrow[d] \\
  \LS^{\restr}_H \arrow[r, "\Delta"] & \LS^{\restr}_H \times \LS^{\restr}_H
\end{tikzcd}
\end{equation*}
其中右端的箭头是 Frobenius 映射的图。

注意 $\LS^{\arthm}_{GL_n}$ 的 $\e$-点是 $X$ 上秩 $n$ 的 lisse Weil 层。更一般地，我们把 $\LS^{\arthm}_{\bG}$ 看作参数化从 Weil 群 $\mathcal{W}_X$ 到代数群 $\bG/\e$ 的连续同态的叠，并把这些同态模共轭。更简洁地说：$\LS^{\arthm}_{\bG}$ 是整体域 $\F_q(X)$ 的无分歧 Langlands 参数的叠。

\begin{remark}
令 $\qLisse^{\phi}(X)$ 表示\emph{拟 lisse Weil 层}的 DG 范畴，按定义它是 $\qLisse(X)$ 上来自 Frobenius 的 $\ZZ$-作用的不动点。同义反复地，$\LS^{\arthm}_{H}$ 参数化对称幺半函子 $\Rep(H) \to \qLisse^{\phi}(X)$，其方式与 $\LS^{\restr}_{H}$ 参数化对称幺半函子 $\Rep(H) \to \qLisse(X)$ 相同（即 $\LS^{\arthm}_{H}$ 的 $S$-点是右 $t$-正合对称幺半函子 $\Rep(H) \to \qLisse(X) \otimes \QCoh(S)$）。

然而，$\qLisse^{\phi}(X)$ 与 $\qLisse(X)$ 具有不同的范畴性质。例如，$\qLisse^{\phi}(X)^{\heartsuit}$ 不是 Tannaka 范畴。这导致两种设定之间的一些形式差异：$\LS^{\arthm}_{H}$ 在某些方面表现得更像 Betti 局部系统的模空间；例如，结果表明（[2] 定理 16.1.4）$\LS^{\arthm}_{H}$ 是一个（非形式的！）代数叠，并且是拟紧的（特别地：只有有限多个连通分支！）。
\end{remark}

\begin{remark}
除了一句评论之外，我们在本节不试图给出更显式的图像：$\LS^{\arthm}_{\bG}$ 的几何比其受限对应物更复杂。但在定理 4.3.3.1 中，我们将在包含平凡表示的一块 $\LS^{\arthm}_{\bG}$ 上给出坐标，从而对其几何提供一些显式分析。
\end{remark}

\pnum{}本质上由注 2.3.1.2，任何受限几何朗兰兹等价\emph{必须}与两边的 Frobenius 自同态相容。

回忆 §2.1.6 中提到的\emph{范畴迹}概念：它取（可对偶化的）DG 范畴连同自函子，产生向量空间。

于是我们可以对受限几何朗兰兹等价的两边取 Frobenius 的迹。如 [2] §16 所概述，$\IndCoh_{\Nilp^{\mathrm{spec}}}(\LS^{\restr}_{\bG})$ 上 Frobenius 的迹与 $\IndCoh(\LS^{\restr}_{\bG})$ 上的相同，即：
\[
  \Gamma(\LS^{\arthm}_{\bG}, \omega),
\]
其中 $\omega$ 是 $\LS^{\arthm}_{\bG}$ 上的\emph{对偶化层}。

另一方面，[4] 的主定理把 $\Shv_{\Nilp}(\Bun_G)$ 上 Frobenius 的迹计算为：
\[
  \Aut^{\unr}_{G,c}.
\]
利用 (2.2)，人们可以把这解释为层—函数对应的一个更高范畴版本（尽管是在特殊情形中，而不是一般的几何现象）。

\pnum{}我们最终得到\emph{算术}猜想：
\[
  \Aut^{\unr}_{G,c} \simeq \Gamma(\LS^{\arthm}_{\bG}, \omega).
\]
正如引言中那样，左端的向量空间是无分歧自守函数的空间，即 $\Bun_G(\F_q)$ 上的紧支撑函数。特别地，无分歧尖形式位于这个空间中。

\pnum{}存在一个典范映射 $\tau: \mathcal{O}_{\LS^{\arthm}_{\bG}} \to \omega_{\LS^{\arthm}_{\bG}}$，编码 $\LS^{\arthm}_{\bG}$ 上的一个“弱 Calabi--Yau”结构——更多讨论参见 §4.6.5。

特别地，存在一个从 $\LS^{\arthm}_{\bG}$ 上的函数到上面右端 $\Gamma(\LS^{\arthm}_{\bG}, \omega)$ 的自然映射。人们应当把 Arthur 的 $SL_2$ 看作在 $\LS^{\arthm}_{\bG}$ 上度量 $\mathcal{O}$ 与 $\omega$ 之间的差异。

\begin{example}
假设 $\sigma \in \LS^{\arthm}_{\bG}$ 是这个叠的一个光滑孤立点。（这样的 $\sigma$ 称为\emph{椭圆}或\emph{离散} Langlands 参数。）于是人们可以看出 $\tau|_\sigma$ 是一个同构。因此，我们的猜想预言存在一个与 $\sigma$ 对应的\footnote{我们对“对应于”的含义有些马虎。更确切地说，我们的猜想结合 [2] §24.2 的讨论，蕴含应当存在 V.~Lafforgue 的 \emph{excursion 代数}作用的一个 $1$ 维无分歧本征形式空间，其特征值由 $\sigma$ 定义。众所周知，对一般 $G$，仅靠 Hecke 算子不足以挑出一个 $1$ 维本征空间。} $1$ 维无分歧自同构形式空间。
\end{example}

\subsection{Xue 的定理与 Frobenius 的迹}

\pnum{}上面我们说：
\begin{equation}
  \mathrm{tr}(\Frob^*_{\Bun_G}, \Shv_{\Nilp}(\Bun_G)) \simeq \Aut^{\unr}_{G,c}
\end{equation}
是 [4] 的主定理。我们简要说明这是如何证明的。关键作用由 Xue 来自 [31] 的定理扮演。

人们也可以转向 [4] 的引言以了解论证概述。我们的概述在这里并无太大不同，只是我们稍微更努力地把 Beilinson 的谱投影扫到地毯下（也许对讨论有害）。

\pnumt{第 1 步。}{} Drinfeld 关于朗兰兹对应的工作\footnote{这里讨论的关于一般约化群的观点来自 [29]。}中的一个教训是：把自守函数 $\Aut^{\unr}_G$ 看作 \emph{shtuka 上同调层}的特殊情形通常很有帮助。

我们简要回顾这个故事。shtuka 构造取一个有限集 $I$、一个表示 $V \in \Rep(\bG^I)$，并产生 $X^I$ 上的一个层 $\Sht_{I,V} \in \Shv(X^I)$。确切地说，附着于 $I$ 与 $V$ 的数据，有一个 \emph{Hecke 函子}：
\[
  H_V: \Shv(\Bun_G) \to \Shv(\Bun_G \times X^I).
\]
这个函子来自一个自然定义的\emph{核} $\mathcal{K}_V \in \Shv(\Bun_G \times \Bun_G \times X^I)$。我们提醒，几何 Satake 在构造中起关键作用，关于构造的更多细节我们参考 [22] 或 [4]。

于是 $\Sht_{I,V}$ 通过沿 Frobenius 的图作 $*$-拉回得到：
\[
  \Bun_G \times X^I \xrightarrow{\mathrm{Graph}_{\Frob} \times \mathrm{id}_{X^I}} \Bun_G \times \Bun_G \times X^I
\]
然后沿 $\Bun_G$ 因子取紧支撑上同调，即推前到 $X^I$。

例如，当 $I = \varnothing$（且 $V$ 是平凡表示的 $1$ 维表示）时，$\mathcal{K}_V = \Delta_!(\mathbf{e}_{\Bun_G})$，所以由基变换，上面的计算给出 $C_{\et,c}(\Bun_G(\F_q)) = \Aut^{\unr}_G$。

shtuka 上同调层之间存在自然的态射。首先，对固定的 $I$，上面的构造给出一个函子 $\Sht_I: \Rep(\bG^I) \to \Shv(X^I)$。但我们也可以变化 $I$；更确切地说，$\Rep(\bG)$ 上的对称幺半结构把指派 $I \mapsto \Rep(\bG^I) = \Rep(\bG)^{\otimes I}$ 变成一个函子 $\mathbf{fSet} \to \DGCat_{\mathrm{cont}}$（其中 $\mathbf{fSet}$ 是有限集的范畴）；沿对角态射的 $*$-拉回使指派 $I \mapsto \Shv(X^I)$ 也变成一个函子 $\mathbf{fSet} \to \DGCat_{\mathrm{cont}}$。于是几何 Satake 的标准函子性说，我们有一个函子之间的自然变换：
\[
  \Sht: (I \mapsto \Rep(\bG^I)) \to (I \mapsto \Shv(X^I)),
\]
即
\[
  \mathbf{fSet} \to \DGCat_{\mathrm{cont}}.
\]
这种函子性是 shtuka 上同调层的一个关键性质，它的存在编码了自守函数的关键对称性：V.~Lafforgue 在 [22] 中正是利用这种函子性来构造 excursion 算子。

\pnumt{第 2 步。}{} 我们现在类似地推广定理的另一边，我们提醒它是 $\mathrm{tr}_{\Shv_{\Nilp}(\Bun_G)}(\Frob)$。答案应当输入 $V \in \Rep(\bG^I)$，并在 $X^I$ 上产生一个层，我们最终将之记为 $\Sht^{\mathrm{tr}}_{I,V}$。当然，构造应当涉及 Hecke 函子与 $\Shv_{\Nilp}(\Bun_G)$，所以我们暂且离题讨论后者。

遵循拓扑设定中的 [27]，我们在 [2] 定理 14.2.4（及其后续讨论）中证明：对 $V \in \Rep(\bG^I)$，Hecke 函子 $H_V$ 把 $\Shv_{\Nilp}(\Bun_G) \subseteq \Shv(\Bun_G)$ 映入 $\Shv_{\Nilp}(\Bun_G) \otimes \qLisse(X^I) \subseteq \Shv(\Bun_G \times X^I)$。

此外，我们在\emph{同上}定理 14.4.3 中证明了一个逆命题：对 $\mathcal{F} \in \Shv(\Bun_G)$，若对所有 $V \in \Rep(\bG)$ 都有 $H_V(\mathcal{F}) \in \Shv(\Bun_G) \otimes \qLisse(X)$，则必有 $\mathcal{F} \in \Shv_{\Nilp}(\Bun_G)$。\footnote{技术上，这个断言中对基域的特征有一些次要限制。回忆 §1.4，我们总是忽略这些小特征，并在讨论中隐含假定我们远离这些特征。}

于是我们用一句格言来总结：$\Shv_{\Nilp}(\Bun_G)$ 可以看作层 $\mathcal{F}$ 的子范畴，使得当我们变化点 $x \in X$ 时，其 \emph{Hecke 变换} $H_{V,x}(\mathcal{F})$ 是局部常值的。

这种关于 $\Shv_{\Nilp}(\Bun_G)$ 的观点实际上对 AGKRRV 系列中几乎\footnote{主要的例外是 [2] 中的 \emph{Künneth 公式}。}每个结果都更好。（从一种观点看，这就是为什么在我们当前的分析中引入 Hecke 函子与一般的 shtuka 层很重要：$\Shv_{\Nilp}(\Bun_G)$ 本身最好用 Hecke 作用来理解。）

\pnumt{第 3 步。}{} 由上，对 $V \in \Rep(\bG)$，我们有一个 Hecke 函子：
\[
  H_V: \Shv_{\Nilp}(\Bun_G) \to \Shv_{\Nilp}(\Bun_G) \otimes \qLisse(X^I).
\]
我们可以把这个函子与 $\Bun_G$ 上的 Frobenius 预复合，得到：
\[
  H_V \circ \Frob_{\Bun_G}: \Shv_{\Nilp}(\Bun_G) \to \Shv_{\Nilp}(\Bun_G) \otimes \qLisse(X^I).
\]
然后我们可以沿\footnote{这类似于说：如果有一个线性变换 $W_1 \to W_1 \otimes W_2$（$W_1$ 有限维），我们有一个相应的向量属于 $W_1^\vee \otimes W_1 \otimes W_2$，并且可以沿前两个因子配对，得到一个“沿 $W_1$ 的迹”，它是 $W_2$ 中的向量。} $\Shv_{\Nilp}(\Bun_G)$ 取迹，得到 $\qLisse(X^I)$ 中的一个对象。这就是所需的对象 $\Sht^{\mathrm{tr}}_{I,V}$。

我们在下面想要证明的是，我们有函子性的等同：
\begin{equation}
  \Sht^{\mathrm{tr}}_{I,V} \simeq \Sht_{I,V}[2|I|],
\end{equation}
其中上同调平移的出现有技术原因，将在下面显现。$I = \varnothing$、$V$ 为 $1$ 维的情形用简洁记号重新得到了 (2.5)。

\pnumt{第 4 步。}{} 观察 $\Sht^{\mathrm{tr}}_{I,V}$ 与 $\Sht_{I,V}$ 之间的一个差异：出于本质的几何原因，$\Sht^{\mathrm{tr}}$ 取值在 $\qLisse(X^I) \subseteq \Shv(X^I)$ 中，而这对 $\Sht_I$ 本身并不显然。

在 [3] 中，我们引入了计算 $\Shv_{\Nilp}(\Bun_G)$ 上迹的方法。细节我们参考\emph{同上}，但简短的答案是：迹可以用\emph{一般的几何要素}（upper-$*$ 与 lower-$!$ 函子）加上\emph{来自（几何）表示论的一个具体要素}来计算。后者是 Beilinson 的谱投影，其工作（就我们的目的而言）是取一个相容的（对 $I$）函子系统 $\mathcal{S}_I: \Rep(\bG^I) \to \Shv(X^I)$，并产生一个相容的系统 $\lambda \mathcal{S}_I: \Rep(\bG^I) \to \qLisse(X^I)$。

在不精确的、高层次的看法下，[3] 的配方给出如下答案：函子 $\Sht^{\mathrm{tr}}$ 的系统是函子 $\Sht_I$ 的系统在取值于 $\qLisse(X^I)$（而不是 $\Shv(X^I)$）中的最佳逼近，即它是 $\lambda \Sht_I$。

于是 Xue 的定理 [31] 说 $\Sht_I$ 本身取值在 $\qLisse(X^I)$ 中，所以 $\lambda \Sht_I$ 与 $\Sht_I$ 本身重合，从而我们得到 (2.6)。

\pnumt{第 5 步。}{} 上面在道德上正确，但我们现在要修正一个谎言。下面的讨论可以与 [4] 注 3.2.6 相比较。

当有限集 $I$ 变化时，函子 $\Sht_I$ 在 upper-$*$ 函子下相容。然而，应用谱投影的程序适用于在 upper-$!$ 函子下相容的函子系统 $\mathcal{S}_I$。所以我们需要 shtuka 函子的一个变体 $\Sht^!$，它在 upper-$!$ 函子下适当相容。

相关的函子 $\Sht^!_{I,-}: \Rep(\bG^I) \to \Shv(X^I)$ 由下面的公式刻画：
\begin{equation}
  C_c(X^I, \Sht^!_{I,V} \boxtimes \mathcal{F}) = C_c(\Bun_G \times X^I, (\Frob \times \mathrm{id})^*(\mathcal{K}_V \boxtimes p_3^*(\mathcal{F})))
\end{equation}
其中 $p_3: \Bun_G \times \Bun_G \times X^I \to X^I$ 是投影。\footnote{一个记号上的说明：我们的记号与 [4] 不一致。我们现在称为 $\Sht^!$ 的函子族被整齐地打包在 [4] 中称为 $\Sht$ 的单个函子里，尽管我们这里称为 $\Sht_I$ 的函子在 [4] 中记号相同。由 Xue 的定理，它们只相差平移。}

先验地，[3] 的结果\emph{实际上}蕴含 $\{\Sht^{\mathrm{tr}}_I\}_{I \in \mathbf{fSet}}$ 是函子 $\{\Sht^!_I\}_{I \in \mathbf{fSet}}$ 在子范畴 $\qLisse \subseteq \Shv$ 中取值的最佳逼近。

逻辑接着通过两次应用 Xue 的定理来进行。首先，该定理告诉我们 $\Sht^!_{I,V} = \Sht_{I,V}[2|I|]$（其中 $2|I|$ 表现为 $2\dim(X^I)$）——我们只需在 (2.7) 的左端用 $\Sht_{I,V}[2|I|]$ 替换 $\Sht^!_{I,V}$，然后利用恒等式 $\mathcal{G} \boxtimes \mathcal{F} = \mathcal{G} \boxtimes \mathcal{F}[-2\dim]$（对 $\mathcal{G}$ lisse）把所得表达式（函子性地）整理成 (2.7) 的右端。特别地，Xue 的定理蕴含 $\Sht^!_{I,V}$ 取值在 $\qLisse$ 中，所以前一段蕴含 $\Sht^{\mathrm{tr}} = \Sht^!$，而我们刚看到这也等于 $\Sht_I[2|I|]$，正如所愿。

% ============================================================
%  ch03a：§3 谱 Eisenstein 级数：引言、3.1 自守 Eisenstein 级数、
%         3.2 一些一般构造（原文 pp.16–19）
% ============================================================
\section{谱 Eisenstein 级数}

我们现在开始走向定理 A。本文余下部分，我们始终假定 $k = \Fb$。

本节的目标是定义并研究如下映射：
\[
  \Eisspec : \Gamma(\LS^{\arthm}_{\bT}, \omega_{\LS^{\arthm}_{\bT}}) \longrightarrow
  \Gamma(\LS^{\arthm}_{\bG}, \omega_{\LS^{\arthm}_{\bG}}).
\]

本节通篇只考虑局部几乎有限型的（预）叠；我们不再一一提及这一假设。

\subsection{自守 Eisenstein 级数}

我们先回顾关于几何 Eisenstein 级数及其函数论对应物——伪 Eisenstein 级数——的若干构造。稍后我们希望在谱侧找到这些构造的对应物。

\pnum{首先，我们有一个典范函子：}
\[
  \Eis_! : \Shv(\Bun_T) \longrightarrow \Shv(\Bun_G),
\]
它由沿 $\Bun_B \to \Bun_T$ 的 $*$-拉回、再沿 $\Bun_B \to \Bun_G$ 的 $!$-推前定义。

由 [6] 中建立的 $\Eis_!$ 的 Hecke 性质，以及 [2] 定理 14.4.3（“Nadler--Yun 定理的逆”，参看 §2.5.3），我们得到：

\begin{prop}
函子 $\Eis_!$ 把 $\qLisse(\Bun_T)$（$= \Shv_{\Nilp_T}(\Bun_T)$）映入 $\Shv_{\Nilp}(\Bun_G)$。
\end{prop}

\pnumt{伪 Eisenstein 级数。}{}\ 令 $\text{ps-Eis} : \Aut^{\unr}_{T,c} \to \Aut^{\unr}_{G,c}$ 为伪 Eisenstein 映射。按定义，它是如下复合：
\[
  \Aut^{\unr}_{T,c} := \mathrm{Func}(\Bun_T(\F_q)) \longrightarrow
  \mathrm{Func}(\Bun_B(\F_q)) \longrightarrow \mathrm{Func}(\Bun_G(\F_q)) := \Aut^{\unr}_{G,c},
\]
即先作限制（注意 $\Bun_B(\F_q) \to \Bun_T(\F_q)$ 的纤维有限），再沿 $\Bun_B(\F_q) \to \Bun_G(\F_q)$ 的纤维求和（这是良定义的，因为我们考虑的是具有有限支撑的函数）。

\pnumt{两者的相容性。}{}\ 函子 $\Eis_!$ 显然缠绕 Frobenius 并保持紧对象，因此我们可以过渡到 Frobenius 的迹，得到映射：
\[
  \tr(\Eis_!) : \tr_{\qLisse(\Bun_T)}(\Frob) \longrightarrow \tr_{\Shv_{\Nilp}(\Bun_G)}(\Frob).
\]
由 [4] 的主定理，我们有同构：
\begin{equation}
  \begin{aligned}
    \tr_{\qLisse(\Bun_T)}(\Frob) &\cong \Aut^{\unr}_{T,c},\\
    \tr_{\Shv_{\Nilp}(\Bun_G)}(\Frob) &\cong \Aut^{\unr}_{G,c}.
  \end{aligned}
\end{equation}

于是 $\tr(\Eis_!)$ 对应于映射 $\Aut^{\unr}_{T,c} \to \Aut^{\unr}_{G,c}$。

由 [4] 定理 5.2.3\footnote{我们指出，前引文献（loc.\ cit.）的陈述（甚至其表述本身）都以 $\Shv_{\Nilp}(\Bun_G)$ 上某个技术性假设为条件；见 [4] §5.1。这一假设最近由作者与 Gaitsgory 验证，将出现在即将发表的工作中。}，同构 (3.1) 由层—函数对应的一个版本给出；由此形式上可得交换图：
\begin{equation}
  \begin{tikzcd}[column sep=large, row sep=large]
    \tr_{\qLisse(\Bun_T)}(\Frob) \arrow[r, "\tr(\Eis_!)"] \arrow[d, "\cong"']
      & \tr_{\Shv_{\Nilp}(\Bun_G)}(\Frob) \arrow[d, "\cong"'] \\
    \Aut^{\unr}_{T,c} \arrow[r, "\text{ps-Eis}"]
      & \Aut^{\unr}_{G,c}.
  \end{tikzcd}
\end{equation}

换言之：\emph{几何 Eisenstein 级数的迹就是伪 Eisenstein 级数。}

\subsection{一些一般构造}

以下除非另有说明，我们只在域 $\e$ 上工作。

下面，我们在一般的叠论语境中给出 $\Eisspec$ 的一个一般构造。

\pnumt{叠的术语。}{}\ 回忆 [8] 中 QCA 叠\zy{QCA 是“quasi-compact with affine diagonal”（拟紧且对角仿射）的缩写；这一概念由 [8] 引入并在其中系统使用。[来源：A/原文 §3.2.1 及其所引 [8]]}这一技术性概念：该术语指拟紧且对角仿射的代数叠。对任意这样的 QCA 叠 $\mathcal{Y}$，[8] 定理 0.4.5 断言 $\IndCoh(\mathcal{Y}) = \mathrm{Ind}(\Coh(\mathcal{Y}))$。此外，由 [8] §3，QCA 叠上的归纳凝聚层有良好的推前理论。

我们也回忆 [2] §5.2 中 ind-代数叠的概念。我们提醒：预叠 $\mathcal{Y}$ 称为\emph{ind-代数的}，如果它是收敛的，且对每个 $n \ge 0$，它的 $n$-截断 ${}^{\le n}\mathcal{Y}$ 可以写成若干 $n$-截断代数叠 $\mathcal{Y}_i$ 沿闭嵌入的滤余极限。如果这些 $\mathcal{Y}_i$ 还可以取为 QCA 的，我们就说 $\mathcal{Y}$ 是\emph{ind-QCA} 的。

我们的主要例子是仿射代数群 $H$ 的 $\LS^{\restr}_H$。由 [2] 推论 5.2.6，$\LS^{\restr}_H$ 是 ind-代数的；此外，该结果的证明表明 $\LS^{\restr}_H$ 实际上是 ind-QCA 的。

由上述 Drinfeld--Gaitsgory 定理，任意 ind-QCA 叠 $\mathcal{Y}$ 的 $\IndCoh(\mathcal{Y})$ 都是紧生成的。同样，ind-QCA 叠之间的态射有良好的 $\IndCoh$-推前理论。

\pnum{下面，我们固定 ind-QCA 代数叠之间的一个 1-可表\footnote{即纤维是代数叠。确切地说，对每个 $S \in \mathrm{AffSch}$ 与 $S \to \mathcal{Z}$，纤维积 $\mathcal{Y} \times_{\mathcal{Z}} S$ 是代数叠。在我们的语境中，这一条件排除了像 $\mathcal{Y} \to \Spec(\e)$ 这样的映射，除非 $\mathcal{Y}$ 是真正的（而非 ind-）代数叠。}映射 $f : \mathcal{Y} \to \mathcal{Z}$。}

此外假设给定自同构 $\phi_{\mathcal{Y}} : \mathcal{Y} \to \mathcal{Y}$ 与 $\phi_{\mathcal{Z}} : \mathcal{Z} \to \mathcal{Z}$，它们被 $f$ 缠绕（即给定同构 $\phi_{\mathcal{Z}} \circ f \simeq f \circ \phi_{\mathcal{Y}}$）。有时我们省略下标，对 $\phi_{\mathcal{Y}}$ 与 $\phi_{\mathcal{Z}}$ 都简写为 $\phi$。

我们构造 $\phi$ 的不动点叠 $\mathcal{Y}^\phi$（相应地 $\mathcal{Z}^\phi$）。确切地说，它是等化子 $\mathrm{Eq}(\mathcal{Y} \rightrightarrows \mathcal{Y})$（两个映射为 $\phi$ 与 $\mathrm{id}$），也可以写成 $\mathcal{Y} \xrightarrow{\mathrm{Graph}_\phi} \mathcal{Y} \times \mathcal{Y} \xleftarrow{\Delta} \mathcal{Y}$ 的纤维积。由假设，我们有诱导映射 $\mathcal{Y}^\phi \to \mathcal{Z}^\phi$，记作 $f^\phi$。

下面，在适当的假设下，我们将构造 $\Gamma(\mathcal{Y}^\phi, \omega_{\mathcal{Y}^\phi})$ 与 $\Gamma(\mathcal{Z}^\phi, \omega_{\mathcal{Z}^\phi})$ 之间的典范映射。

我们把讨论组织如下。在 §3.2.3 与 §3.2.4 中，我们利用 $\IndCoh$ 的标准函子性给出这些映射的“初等”构造。后一个构造尤其有些复杂。之后，我们解释一个使这些构造的存在性变得显然的概念框架（迹的函子性）。作者本人是从后一种观点来思考这些映射的，但担心若缺少具体描述，读者不会相信这套魔法；不需要这种说服的读者可以跳过 §3.2.3--3.2.4。

我们还要指出，ind-QCA 假设是过强的。§3.2.3 并不需要它。它在 §3.2.4 中被轻度使用，用于各种推前函子的存在性，不过更弱的假设已经足够。从根本上说，从更概念化的迹的视角看，这一假设是自然的，它保证 $\IndCoh$ 的紧生成性（从而保证可对偶性）。

\pnumt{推前。}{}\ 首先，假设映射 $f$ 可表且紧合（proper）。此时我们将构造映射：
\begin{equation}
  \Gamma^{\IndCoh}(\mathcal{Y}^\phi, \omega_{\mathcal{Y}^\phi}) \longrightarrow
  \Gamma^{\IndCoh}(\mathcal{Z}^\phi, \omega_{\mathcal{Z}^\phi}).
\end{equation}

事实上，这相当容易。此时映射 $f^\phi : \mathcal{Y}^\phi \to \mathcal{Z}^\phi$ 也是紧合（且可表）的\footnote{确实，由于 $f$ 可表且分离，态射 $\mathcal{Y}^\phi = \mathcal{Y} \times_{\mathcal{Y} \times \mathcal{Y}} \mathcal{Y} \to \mathcal{Y} \times_{\mathcal{Z} \times \mathcal{Z}} \mathcal{Z}$ 是闭嵌入。显然，进一步的投影 $\mathcal{Y} \times_{\mathcal{Z} \times \mathcal{Z}} \mathcal{Z} \to \mathcal{Z} \times_{\mathcal{Z} \times \mathcal{Z}} \mathcal{Z} = \mathcal{Z}^\phi$ 是紧合的。}，而这正是下面所需的全部。

于是 $f^{\phi,\IndCoh}_*$ 左伴随于 $f^{\phi,!}$，故我们得到典范伴随映射 $f^{\phi,\IndCoh}_* f^{\phi,!} \to \mathrm{id}$。这给出映射 $f^{\phi,\IndCoh}_*(\omega_{\mathcal{Y}^\phi}) \to \omega_{\mathcal{Z}^\phi}$；再作用 $\Gamma^{\IndCoh}(\mathcal{Z}^\phi, -)$，即得所需映射 (3.3)。

\begin{remark}
映射 (3.3) 是 $\Gamma(\mathcal{Z}^\phi, \OO_{\mathcal{Z}^\phi})$-线性的。
\end{remark}

\pnumt{拉回。}{}\ 其次，假设 $f$ 最终余连通（且 1-可表）。此时我们将构造映射：
\begin{equation}
  \Gamma^{\IndCoh}(\mathcal{Z}^\phi, \omega_{\mathcal{Z}^\phi}) \longrightarrow
  \Gamma^{\IndCoh}(\mathcal{Y}^\phi, \omega_{\mathcal{Y}^\phi}).
\end{equation}

首先注意\footnote{这一点以及本节中其他类似断言，形式上都可以归约到 [17] 第 4 章 §3 的结果。}，此时 $f^{!,\IndCoh} : \IndCoh(\mathcal{Y}) \to \IndCoh(\mathcal{Z})$ 具有左伴随 $f^{*,\IndCoh}$。我们有 $*$-推前与 $*$-拉回之间的基变换（后者只对最终余连通的态射考虑）。

此时，我们有自然变换
\[
  \begin{aligned}
    &(\mathrm{id} \times f)^{*,\IndCoh}(\mathrm{id} \times \phi_{\mathcal{Z}})^{\IndCoh}_*
      (\mathrm{id} \times f)^{\IndCoh}_* = {}\\
    &(\mathrm{id} \times f)^{*,\IndCoh}(\mathrm{id} \times f)^{\IndCoh}_*
      (\mathrm{id} \times \phi_{\mathcal{Y}})^{\IndCoh}_*\\
    &\qquad\qquad \longrightarrow (\mathrm{id} \times \phi_{\mathcal{Y}})^{\IndCoh}_*
      \in \mathrm{End}(\IndCoh(\mathcal{Y} \times \mathcal{Y})),
  \end{aligned}
\]
它来自伴随。把这一映射作用于 $\Delta^{\IndCoh}_*(\omega_{\mathcal{Y}})$，我们得到典范映射：
\begin{equation}
  \begin{aligned}
    &(\mathrm{id} \times f)^{*,\IndCoh}
      \mathrm{Graph}^{\IndCoh}_{\phi_{\mathcal{Z}} \circ f,*}(\omega_{\mathcal{Y}})\\
    &\qquad\qquad \longrightarrow \mathrm{Graph}_{\phi_{\mathcal{Y}},*}(\omega_{\mathcal{Y}})
      \in \IndCoh(\mathcal{Y} \times \mathcal{Y}).
  \end{aligned}
\end{equation}

这里，对映射 $g : S \to T$，映射 $\mathrm{Graph}_g : S \to S \times T$ 是图像态射，即 $\mathrm{Graph}_g := (\mathrm{id} \times g) \circ \Delta_S$。

令 $\varpi_{\mathcal{Y}}$ 表示典范映射 $\mathcal{Y}^\phi \to \mathcal{Y}$，它把一对 $(y \in \mathcal{Y}, y \simeq \phi_{\mathcal{Z}}(y))$ 送到 $y$；$\varpi_{\mathcal{Z}}$ 类似。% [Tutor Note: 原文此处作 φ_𝒵(y)，但 𝒴^φ 由 φ_𝒴 定义，语义上应为 φ_𝒴(y)；译文保留原文写法，见报告。]

下面，我们将构造典范同构：
\begin{equation}
  \begin{aligned}
    &\Delta^!_{\mathcal{Y}}(\mathrm{id} \times f)^{*,\IndCoh}
      \mathrm{Graph}^{\IndCoh}_{\phi_{\mathcal{Z}} \circ f,*}(\omega_{\mathcal{Y}})\\
    &\qquad\qquad \simeq f^{*,\IndCoh}\varpi^{\IndCoh}_{\mathcal{Z},*}(\omega_{\mathcal{Z}^\phi})
      \in \IndCoh(\mathcal{Y}).
  \end{aligned}
\end{equation}

假设这一构造已经给出，则把 $\Delta^!_{\mathcal{Y}}$ 作用于 (3.5)，我们得到典范映射
\[
  f^{*,\IndCoh}\varpi^{\IndCoh}_{\mathcal{Z},*}(\omega_{\mathcal{Z}^\phi})
  \longrightarrow \Delta^!_{\mathcal{Y}}\mathrm{Graph}_{\phi_{\mathcal{Y}},*}(\omega_{\mathcal{Y}})
  \in \IndCoh(\mathcal{Y}).
\]
由伴随，这又给出典范映射
\[
  \begin{aligned}
    \varpi^{\IndCoh}_{\mathcal{Z},*}(\omega_{\mathcal{Z}^\phi})
    &\longrightarrow f^{\IndCoh}_*\Delta^!_{\mathcal{Y}}
      \mathrm{Graph}_{\phi_{\mathcal{Y}},*}(\omega_{\mathcal{Y}})\\
    &\simeq f^{\IndCoh}_*\varpi^{\IndCoh}_{\mathcal{Y},*}(\omega_{\mathcal{Y}^\phi})
      \in \IndCoh(\mathcal{Z}).
  \end{aligned}
\]
这里我们用了基变换同构 $\Delta^!_{\mathcal{Y}}\mathrm{Graph}^{\IndCoh}_{\phi_{\mathcal{Y}},*} \simeq \varpi^{\IndCoh}_{\mathcal{Y},*}\varpi^!_{\mathcal{Y}}$。现在把整体截面函子 $\Gamma^{\IndCoh}(\mathcal{Z}, -)$ 作用于上式两端，即得所需映射 (3.4)。

还剩下同构 (3.6) 的给出。首先，纤维积图表
\[
  \begin{tikzcd}[column sep=large, row sep=large]
    \mathcal{Y} \arrow[r, "\mathrm{Graph}_{\phi_{\mathcal{Z}} \circ f}"] \arrow[d, "f"]
      & \mathcal{Y} \times \mathcal{Z} \arrow[d, "f \times \mathrm{id}"] \\
    \mathcal{Z} \arrow[r, "\mathrm{Graph}_{\phi_{\mathcal{Z}}}"]
      & \mathcal{Z} \times \mathcal{Z}
  \end{tikzcd}
\]
给出基变换同构
\[
  \mathrm{Graph}^{\IndCoh}_{\phi_{\mathcal{Z}} \circ f,*}(\omega_{\mathcal{Y}})
  \simeq (f \times \mathrm{id})^!\mathrm{Graph}^{\IndCoh}_{\phi_{\mathcal{Z}},*}(\omega_{\mathcal{Z}}).
\]
这又给出
\[
  \begin{aligned}
    &\Delta^!_{\mathcal{Y}}(\mathrm{id} \times f)^{*,\IndCoh}
      \mathrm{Graph}^{\IndCoh}_{\phi_{\mathcal{Z}} \circ f,*}(\omega_{\mathcal{Y}})\\
    &\qquad\qquad \simeq \Delta^!_{\mathcal{Y}}(\mathrm{id} \times f)^{*,\IndCoh}
      (f \times \mathrm{id})^!\mathrm{Graph}^{\IndCoh}_{\phi_{\mathcal{Z}},*}(\omega_{\mathcal{Z}}).
  \end{aligned}
\]

我们有\footnote{这类交换性是一般事实：见 [12] 命题 7.1.6。不过，在目前的设定中它尤其容易：由 Künneth 公式，我们可以把 $(f \times \mathrm{id})^!$ 写成 $f^! \otimes \mathrm{id}$，把 $(\mathrm{id} \times f)^{*,\IndCoh}$ 写成 $\mathrm{id} \otimes f^{*,\IndCoh}$。}
$(\mathrm{id} \times f)^{*,\IndCoh}(f \times \mathrm{id})^! \simeq (f \times \mathrm{id})^!(\mathrm{id} \times f)^{*,\IndCoh}$，于是可把上式右端改写为
\[
  \begin{aligned}
    &\Delta^!_{\mathcal{Y}}(f \times \mathrm{id})^!(\mathrm{id} \times f)^{*,\IndCoh}
      \mathrm{Graph}^{\IndCoh}_{\phi_{\mathcal{Z}},*}(\omega_{\mathcal{Z}})\\
    &\qquad\qquad = \mathrm{Graph}^{\sigma,!}_f(\mathrm{id} \times f)^{*,\IndCoh}
      \mathrm{Graph}^{\IndCoh}_{\phi_{\mathcal{Z}},*}(\omega_{\mathcal{Z}}).
  \end{aligned}
\]
这里我们用 $\mathrm{Graph}^\sigma_g$ 表示 $(g \times \mathrm{id}) \circ \Delta$，即交换两个笛卡尔因子之后的图像映射。

现在，对纤维积图表
\[
  \begin{tikzcd}[column sep=large, row sep=large]
    \mathcal{Y} \arrow[r, "\mathrm{Graph}^\sigma_f"] \arrow[d, "f"]
      & \mathcal{Z} \times \mathcal{Y} \arrow[d, "\mathrm{id} \times f"] \\
    \mathcal{Z} \arrow[r, "\Delta_{\mathcal{Z}}"]
      & \mathcal{Z} \times \mathcal{Z}
  \end{tikzcd}
  \qquad
  \begin{tikzcd}[column sep=large, row sep=large]
    \mathcal{Z}^\phi \arrow[r, "\varpi_{\mathcal{Z}}"] \arrow[d, "\varpi_{\mathcal{Z}}"]
      & \mathcal{Z} \arrow[d, "\mathrm{Graph}_{\phi_{\mathcal{Z}}}"] \\
    \mathcal{Z} \arrow[r, "\Delta_{\mathcal{Z}}"]
      & \mathcal{Z} \times \mathcal{Z}
  \end{tikzcd}
\]
作基变换，得到同构
\[
  \begin{aligned}
    \mathrm{Graph}^{\sigma,!}_f(\mathrm{id} \times f)^{*,\IndCoh}
      \mathrm{Graph}^{\IndCoh}_{\phi_{\mathcal{Z}},*}(\omega_{\mathcal{Z}})
    &\simeq f^{*,\IndCoh}\Delta^!_{\mathcal{Z}}
      \mathrm{Graph}^{\IndCoh}_{\phi_{\mathcal{Z}},*}(\omega_{\mathcal{Z}})\\
    &\simeq f^{*,\IndCoh}\varpi^{\IndCoh}_{\mathcal{Z},*}(\omega_{\mathcal{Z}^\phi}),
  \end{aligned}
\]
正如所愿。

\begin{remark}
映射 (3.4) 是 $\Gamma(\mathcal{Z}^\phi, \OO_{\mathcal{Z}^\phi})$-线性的。
\end{remark}

\pnumt{范畴框架。}{}\ 我们现在给出处理上述这类构造的一个更概念化的途径。

假设 $F : \mathcal{C} \to \mathcal{D} \in \DGCat_{\mathrm{cont}}$ 是可对偶 DG 范畴之间的映射。假设 $\mathcal{C}$（相应地 $\mathcal{D}$）配有自函子 $\phi_{\mathcal{C}}$（相应地 $\phi_{\mathcal{D}}$），并且：
\begin{itemize}
  \item $F$ 具有连续右伴随 $F^R$。
  \item $F$ 以 lax 方式缠绕 $\phi$，即给定映射（常常是同构）$F \circ \phi_{\mathcal{C}} \to \phi_{\mathcal{D}} \circ F$。
\end{itemize}
于是迹的标准函子性给出与这份数据相伴的典范映射
\[
  \tr_{\mathcal{C}}(\phi_{\mathcal{C}}) \longrightarrow \tr_{\mathcal{D}}(\phi_{\mathcal{D}}) \in \Vect.
\]
确切地说，我们有
\[
  \tr_{\mathcal{C}}(\phi_{\mathcal{C}}) \longrightarrow \tr_{\mathcal{C}}(\phi_{\mathcal{C}} F^R F)
  \simeq \tr_{\mathcal{D}}(F \phi_{\mathcal{C}} F^R) \longrightarrow \tr_{\mathcal{D}}(\phi_{\mathcal{D}} F F^R)
  \longrightarrow \tr_{\mathcal{D}}(\phi_{\mathcal{D}}).
\]
这里我们用了迹的循环性与标准的伴随映射。

\pnumt{推前/拉回再论。}{}\ 首先假设给定如上可表且紧合的 $f : \mathcal{Y} \to \mathcal{Z}$。我们还假设 $\mathcal{Y}$ 与 $\mathcal{Z}$ 是 ind-QCA 的（以确保 $\IndCoh$ 可对偶）。

在 §3.2.5 的设定中，取 $\mathcal{C} = \IndCoh(\mathcal{Y})$、$\mathcal{D} = \IndCoh(\mathcal{Z})$、$F = f^{\IndCoh}_*$、$\phi_{\mathcal{C}} = \phi^{\IndCoh}_{\mathcal{Y},*}$，以及 $\phi_{\mathcal{D}} = \phi^{\IndCoh}_{\mathcal{Z},*}$。

于是范畴形式体系给出典范映射\footnote{这里等号都是标准的同构；例如见 [15] §3.5.3。}：
\[
  \Gamma^{\IndCoh}(\mathcal{Y}^\phi, \omega_{\mathcal{Y}^\phi})
  = \tr_{\IndCoh(\mathcal{Y})}(\phi^{\IndCoh}_{\mathcal{Y},*})
  \longrightarrow \tr_{\IndCoh(\mathcal{Z})}(\phi^{\IndCoh}_{\mathcal{Z},*})
  = \Gamma^{\IndCoh}(\mathcal{Z}^\phi, \omega_{\mathcal{Z}^\phi}).
\]
一次直接的图表追踪表明，这一映射还原出 (3.3)。我们指出，紧合性是 $f^{\IndCoh}_*$ 具有连续右伴随所需要的。

其次，取 $f : \mathcal{Y} \to \mathcal{Z}$ 1-可表且最终余连通。现在取 $\mathcal{C} = \IndCoh(\mathcal{Z})$、$\mathcal{D} = \IndCoh(\mathcal{Y})$、$F = f^{\IndCoh,*}$、$\phi_{\mathcal{C}} = \phi^!_{\mathcal{Z}}$、$\phi_{\mathcal{D}} = \phi^!_{\mathcal{Y}}$。于是范畴形式体系给出典范映射\footnote{注意 $\tr_{\mathcal{C}}(F) = \tr_{\mathcal{C}^\vee}(F^\vee)$，故 $\tr(\phi^!) = \tr(\phi^{\IndCoh}_*)$。}：
\[
  \Gamma^{\IndCoh}(\mathcal{Z}^\phi, \omega_{\mathcal{Z}^\phi})
  = \tr_{\IndCoh(\mathcal{Z})}(\phi^!_{\mathcal{Z}})
  \longrightarrow \tr_{\IndCoh(\mathcal{Y})}(\phi^!_{\mathcal{Y}})
  = \Gamma^{\IndCoh}(\mathcal{Y}^\phi, \omega_{\mathcal{Y}^\phi}).
\]
一次图表追踪表明，这一映射与 (3.4) 一致。

\begin{remark}
第二次图表追踪是例行公事，但不出所料，比第一次要复杂一些。我们在此省略验证。实际上，就我们的目的而言，读者可以把这些范畴构造当作 (3.3) 与 (3.4) 的\emph{定义}，完全忽略 §3.2.3--3.2.4 的材料。我们之所以纳入这些显式构造，只是为了让大家看到构造的具体样子。
\end{remark}

% ============================================================
%  ch03b：§3 谱 Eisenstein 级数：3.3 谱 Eisenstein 级数、
%         3.4 谱 vs. 函数论 Eisenstein 级数（原文 pp.19–21）
% ============================================================
\subsection{谱 Eisenstein 级数}

\pnumt{设定。}{}\ 我们有标准的对应：
\begin{equation*}
  \begin{tikzcd}[column sep=large, row sep=large]
    & \LS^{\restr}_{\bB} \arrow[dl, "q"'] \arrow[dr, "p"] & \\
    \LS^{\restr}_{\bT} & & \LS^{\restr}_{\bG}.
  \end{tikzcd}
\end{equation*}
映射 $p$ 可表且紧合，而映射 $q$ 拟光滑且 1-可表。此外，这些空间各自带有一个 Frobenius 自映射，且映射 $p$ 与 $q$ 缠绕这些 Frobenius。因此，由 (3.3) 与 (3.4)，我们得到典范映射：
\[
  \Gamma(\LS^{\arthm}_{\bT}, \omega_{\LS^{\arthm}_{\bT}}) \longrightarrow
  \Gamma(\LS^{\arthm}_{\bB}, \omega_{\LS^{\arthm}_{\bB}}) \longrightarrow
  \Gamma(\LS^{\arthm}_{\bG}, \omega_{\LS^{\arthm}_{\bG}}).
\]

\begin{defn}
上述映射的复合称为\emph{谱 Eisenstein 级数}
$\Eisspec : \Gamma(\LS^{\arthm}_{\bT}, \omega_{\LS^{\arthm}_{\bT}}) \to
\Gamma(\LS^{\arthm}_{\bG}, \omega_{\LS^{\arthm}_{\bG}})$。
\end{defn}

\subsection{谱 vs.\ 函数论 Eisenstein 级数}

\pnum{回忆：\emph{受限几何朗兰兹}（restricted geometric Langlands，见猜想 2.3.1）预测范畴之间存在如下等价：}
\begin{equation}
  \Shv_{\Nilp}(\Bun_G) \simeq \IndCoh_{\Nilp^{\mathrm{spec}}}(\LS^{\restr}_{\bG}).
\end{equation}

这一等价应当满足各种相容性。我们在此强调其中两个特别值得注意的：
\begin{itemize}
  \item（Hecke 相容性）等价 (3.7) 是 $\QCoh(\LS^{\restr}_{\bG})$-模范畴之间的等价；这里右端具有显然的作用，而左端带有 [2] 定理 14.3.2 的作用。
  \item（Eisenstein 相容性，$P = B$ 情形\zy{$P = B$ 指把抛物子群 $P$ 取成 Borel 子群 $B$ 的情形，即极小抛物情形。[来源：A/原文 §3.4.1 括注「$P = B$ case」]}）图表
  \begin{equation*}
    \begin{tikzcd}[column sep=large, row sep=large]
      \qLisse(\Bun_T) \arrow[r, "\Eis_!"] \arrow[d]
        & \Shv_{\Nilp}(\Bun_G) \arrow[d, "\cong"] \\
      \QCoh(\LS^{\restr}_{\bT}) \arrow[r, "\Eisspec"]
        & \IndCoh_{\Nilp^{\mathrm{spec}}}(\LS^{\restr}_{\bG})
    \end{tikzcd}
  \end{equation*}
  交换；这里左侧的竖直箭头是 [2] 例 21.2.9 中无条件构造出的等价。
\end{itemize}

\pnum{我们现在回忆如下结果：}

\begin{prop}
设 $\mathcal{Y}$ 是拟光滑的 ind-代数叠，带有自映射 $\phi_{\mathcal{Y}}$。令
$\mathcal{N} \subseteq T^*[-1]\mathcal{Y}$ 是闭锥形子叠（位于 $\mathcal{Y}$ 的 $-1$-平移余切丛中），使得对每个点
$y \in \mathcal{Y}^\phi$，映射 $d\phi[-1] : \mathcal{N}_y \to \mathcal{N}_y$ 都收缩到 $0 \in T^*_y[-1]\mathcal{Y}$。则映射
\[
  \tr_{\IndCoh_{\mathcal{N}}(\mathcal{Y})}(\phi^!)
  \longrightarrow \tr_{\IndCoh(\mathcal{Y})}(\phi^!)
  = \Gamma(\mathcal{Y}^\phi, \omega_{\mathcal{Y}^\phi})
\]
是同构。
\end{prop}

见 [2] §24.6.8。\footnote{在前引文献（loc.\ cit.）中，这一结果是以猜想的形式表述的。但它实际上可以由 [1] 的形式体系直接证明。}

特别地，我们发现受限几何朗兰兹给出同构：
\begin{equation}
  \Aut^{\unr}_{G,c} \simeq \Gamma(\LS^{\arthm}_{\bG}, \omega_{\LS^{\arthm}_{\bG}}).
\end{equation}

在受限几何朗兰兹的 Hecke 相容性之下，这是
$\Exc_{\bG} := \Gamma(\LS^{\arthm}_{\bG}, \OO_{\LS^{\arthm}_{\bG}})$-模之间的等价，其中左端从 [4] 继承其
$\Exc_{\bG}$-模结构。我们提醒（见 [2] §24.2）：$\Aut^{\unr}_c$ 上的 $\Exc_{\bG}$-模结构细化了
巡游算子在它上面的 Lafforgue--Xue 作用（见 §1.2.3）。

\pnum{我们现在得到：}

\begin{prop}
假设受限几何朗兰兹对 $G$ 成立，并带有其 Eisenstein 相容性。则如下图表交换：
\begin{equation*}
  \begin{tikzcd}[column sep=large, row sep=large]
    \Aut^{\unr}_{T,c} \arrow[r, "\text{ps-Eis}"] \arrow[d, "\cong"]
      & \Aut^{\unr}_{G,c} \arrow[d, "\cong"] \\
    \Gamma(\LS^{\arthm}_{\bT}, \omega_{\LS^{\arthm}_{\bT}}) \arrow[r, "\Eisspec"]
      & \Gamma(\LS^{\arthm}_{\bG}, \omega_{\LS^{\arthm}_{\bG}})
  \end{tikzcd}
\end{equation*}
其中竖直同构来自 (3.8)。
\end{prop}

事实上，这来自 ps-Eis 与 $\Eisspec$ 通过迹的实现，见 (3.2) 与 §3.2.6。

% ============================================================
%  ch03c：§3.5 主结果的表述（局部化；定理 3.5.2.1）（原文 pp.21–22）
%  边界：本块起于 §3.5 标题（p021 中部）。p021 顶部的 3.4.3 与命题 3.4.3.1
%  属 §3.4，归 ch03b。本块止于定理 3.5.2.1（p022 中部）；§3.6 起于 p022
%  下半页，归 ch03d。
% ============================================================
\subsection{主要结果的表述}

\pnum{我们将考察映射 $\Eisspec$ 的\emph{局部化}版本。下面简要讨论相关的形式体系。}

设给定交换图：
\[
\begin{tikzcd}[column sep=huge, row sep=large]
  & \LS^{\arthm}_{\bB} \arrow[dl] \arrow[dr] & \\
  \LS^{\arthm}_{\bG} \arrow[dr, "f"{below}] & & \LS^{\arthm}_{\bT} \arrow[dl, "g"{below}] \\
  & \A^1 &
\end{tikzcd}
\]

令 $\mathring{\A}^1 := \A^1\setminus 0$。注意：
\[
  \Gamma(\LS^{\arthm}_{\bG},\omega)[f^{-1}]
  := \colim_{f\cdot-}\Gamma\bigl(\LS^{\arthm}_{\bG},\omega_{\LS^{\arthm}_{\bG}}\bigr)
  \simeq \Gamma\bigl(\LS^{\arthm}_{\bG}\times_{\A^1}\mathring{\A}^1,\omega\bigr)
\]
类似地，可以用 $(\check T,g)$ 或 $(\check B, f|_{\LS^{\arthm}_{\check B}} = g|_{\LS^{\arthm}_{\check B}})$ 代替 $(\bG,f)$。

现在注意到我们有一个对应：
\begin{equation}
\begin{tikzcd}[column sep=huge, row sep=large]
  & \LS^{\arthm}_{\bB}\times_{\A^1}\mathring{\A}^1 \arrow[dl] \arrow[dr] & \\
  \LS^{\arthm}_{\bG}\times_{\A^1}\mathring{\A}^1 & & \LS^{\arthm}_{\bT}\times_{\A^1}\mathring{\A}^1
\end{tikzcd}
\end{equation}
其中左箭头紧合且可表，右箭头最终余连通且 1-可表。% [Tutor Note: proper 首次括注按 glossary 归 ch03a，此处只用「紧合」。]如同在谱 Eisenstein 级数的定义中那样，我们得到典范映射：
\begin{equation}
  \Gamma\bigl(\LS^{\arthm}_{\bT},\omega_{\LS^{\arthm}_{\bT}}\bigr)[g^{-1}]
  \longrightarrow
  \Gamma\bigl(\LS^{\arthm}_{\bG},\omega_{\LS^{\arthm}_{\bG}}\bigr)[f^{-1}].
\end{equation}
我们显然有：

\begin{lem}
\begin{enumerate}[label=(\arabic*), leftmargin=2.2em]
\item 映射 $\Eisspec:\Gamma(\LS^{\arthm}_{\bT},\omega_{\LS^{\arthm}_{\bT}})\to\Gamma(\LS^{\arthm}_{\bG},\omega_{\LS^{\arthm}_{\bG}})$
把乘以 $g$ 的算子与乘以 $f$ 的算子衔接起来，即该映射自然地是 $\e[t]$-模的态射。
\item 映射 (3.10) 由反转 $t$ 的作用得到，即在 $\e[t]$ 上张量 $\e[t,t^{-1}]$。
\item 图表
\[
\begin{tikzcd}[column sep=huge, row sep=large]
  \Gamma\bigl(\LS^{\arthm}_{\bT},\omega_{\LS^{\arthm}_{\bT}}\bigr)
    \arrow[r, "\Eisspec"] \arrow[d]
  & \Gamma\bigl(\LS^{\arthm}_{\bG},\omega_{\LS^{\arthm}_{\bG}}\bigr) \arrow[d] \\
  \Gamma\bigl(\LS^{\arthm}_{\bT},\omega_{\LS^{\arthm}_{\bT}}\bigr)[g^{-1}]
    \arrow[r, "(3.10)"]
  & \Gamma\bigl(\LS^{\arthm}_{\bG},\omega_{\LS^{\arthm}_{\bG}}\bigr)[f^{-1}]
\end{tikzcd}
\]
交换。
\end{enumerate}
\end{lem}

据此，当含义无歧义时，我们滥用记号，把映射 (3.10) 也记作 $\Eisspec$。

\pnum{现在我们能够陈述关于 $\Eisspec$ 的主要结果。}

\begin{thm}
存在函数 $\delta_{\bG}\in\Om^{\infty}\Gamma(\LS^{\arthm}_{\bG},\OO)$ 与
$\delta_{\bT}\in\Om^{\infty}\Gamma(\LS^{\arthm}_{\bT},\OO)$，它们填入交换图
\begin{equation}
\begin{tikzcd}[column sep=huge, row sep=large]
  & \LS^{\arthm}_{\bB} \arrow[dl] \arrow[dr] & \\
  \LS^{\arthm}_{\bG} \arrow[dr, "\delta_{\bG}"'] & & \LS^{\arthm}_{\bT} \arrow[dl, "\delta_{\bT}"] \\
  & \A^1 &
\end{tikzcd}
\end{equation}
并且满足：
\begin{enumerate}[label=(\arabic*), leftmargin=2.2em]
\item 映射 $\delta_{\bG}$ 在平凡 Weil 局部系统处取非零值。
\item 对象
\[
  \Gamma\bigl(\LS^{\arthm}_{\bG},\omega_{\LS^{\arthm}_{\bG}}\bigr)[\delta_{\bG}^{-1}]\in\Vect
\]
位于上同调次数 $0$，\footnote{注意，(3.8) 预言
$\Gamma(\LS^{\arthm}_{\bG},\omega_{\LS^{\arthm}_{\bG}})$ 位于上同调次数 $0$。尽管我们最终会假设受限几何朗兰兹，
我们在此力求陈述一个不依赖于它的定理，因此保留了这个论断。}以 $\bT$ 代替 $\bG$ 亦类似。
\item 映射
\begin{equation}
  \Eisspec:\Gamma\bigl(\LS^{\arthm}_{\bT},\omega_{\LS^{\arthm}_{\bT}}\bigr)[\delta_{\bT}^{-1}]
  \longrightarrow
  \Gamma\bigl(\LS^{\arthm}_{\bG},\omega_{\LS^{\arthm}_{\bG}}\bigr)[\delta_{\bG}^{-1}]
  \in\Vect^{\heartsuit}
\end{equation}
是满射。
\end{enumerate}
\end{thm}

该结果将在 §4 中证明。

% ============================================================
%  ch03d：§3.6 定理 A 的证明、§3.7 玩具模型（原文 pp.22–25）
%  边界：本块起于 §3.6 标题（p022 下半页），止于推论 3.7.2.1（p025 中部）；
%  其后的 §4 归 ch04a。p022 上半页的定理 3.5.2.1 归 ch03c。
% ============================================================
\subsection{定理 A 的证明}

我们现在由定理 3.5.2.1 与先前的观察推出本文的主定理。此处提醒，我们已经假设 $G$ 是半单的。\footnote{这个假设使讨论略有简化。适当表述后，本文的结果对一般约化群同样适用。}% [Tutor Note: 原文句末在脚注标记之后残留 "here."（"我们已经假设 G 是半单的。\footnote{...} here."），语序不通，疑为作者编辑残留；译文把 here 的含义并入「此处」。]

\pnum{首先，我们按照 [22] 与 [30] 回忆 Langlands 参数的显式含义。}

设 $\sigma$ 是 $\LS^{\arthm}_{\bG}$ 的一个 $\e$-点，即 $X$ 上的一个 Weil $\bG$-局部系统。我们得到一个映射
\[
  \mathrm{ev}_\sigma : \Exc_{\bG} \to \e
\]
它把函数 $f\in\Exc=\Gamma(\LS^{\arthm}_{\bG},\OO)$ 送到其在 $\sigma$ 处的值。我们滥用记号，也用 $\mathrm{ev}_\sigma$ 记诱导映射（通过取 $H^0$ 得到）$H^0(\Exc)\to\e$，它是经典交换代数之间的映射。我们用 $\mathfrak m_\sigma\subseteq H^0(\Exc)$ 记相应的极大理想。

现在回忆（来自 [4]，它建立在 [22] 与 [30] 之上）：$\Exc$——从而 $H^0(\Exc)$——作用在 $\Aut^{\unr}_{G,c}$ 上。

于是我们定义
\[
  \Aut^{\unr}_{G,c,[\sigma]} \subseteq \Aut^{\unr}_{G,c}
\]
为右端的 $\mathfrak m_\sigma$-挠部分，即 $\psi\in\Aut^{\unr}_{G,c,[\sigma]}$ 当且仅当 $\mathfrak m_\sigma^n\cdot\psi=0$ 对 $n\gg 0$ 成立。

我们令 $\Aut^{\unr}_{G,\cusp,[\sigma]} := \Aut^{\unr}_{G,c,[\sigma]}\cap\Aut^{\unr}_{G,\cusp}$。

\begin{warning}
由于 $\Aut^{\unr}_{G,\cusp}$ 是有限维的\footnote{我们提醒（这里用到 $G$ 半单），存在定义在 $\F_q$ 上的拟紧开子集 $\mathcal U\subseteq\Bun_G$，使得任何 $\psi\in\Aut^{\unr}_{G,\cusp}$ 在 $\mathcal U(\F_q)$ 之外为零（参见 [9] 命题 1.4.6 的层论版本）；由于 $\mathcal U(\F_q)$ 有限，断言显然成立。}，$\Aut^{\unr}_{G,\cusp}$ 分解为直和：
\begin{equation}
  \Aut^{\unr}_{G,\cusp} \simeq \bigoplus_{\sigma/\sim}\Aut^{\unr}_{G,\cusp,[\sigma]}.
\end{equation}
（这里隐含的等价关系 $\sim$ 把 $\mathfrak m_{\sigma_1}=\mathfrak m_{\sigma_2}$ 的 $\sigma_1$ 与 $\sigma_2$ 等同起来；根据 [22] 命题 0.38，这恰好发生在 $\sigma_1$ 与 $\sigma_2$ 具有等价半单化之时。）
\end{warning}

然而，对紧支撑自守函数，我们并没有类似的分解 (3.13)；存在这样的函数，它们不属于任何 $\Aut^{\unr}_{G,c,[\sigma]}$。

\pnum{回忆\footnote{参见 [26] II.2.4 中一个强得多的论断。} $\mathrm{ps}\text{-}\Eis(\Aut^{\unr}_{T,c})\cap\Aut^{\unr}_{G,\cusp}=0$。}

因此，只需证明任何 $\psi\in\Aut^{\unr}_{G,c,[\triv]}\subseteq\Aut^{\unr}_{G,c}$ 都落在映射 $\mathrm{ps}\text{-}\Eis:\Aut^{\unr}_{T,c}\to\Aut^{\unr}_{G,c}$ 的像中。这将是我们的目标。

\pnum{回忆 $\delta_{\bG}\in\Exc_{\bG}:=\Gamma(\LS^{\arthm}_{\bG},\OO)$。由于 $\Aut^{\unr}_{G,c}$ 被 $\Exc_{\bG}$ 作用，我们可以反转 $\delta_{\bG}$ 的作用：}
\[
  \Aut^{\unr}_{G,c}[\delta_{\bG}^{-1}] := \colim_{\delta_{\bG}\cdot-}\Aut^{\unr}_{G,c}.
\]

现在我们借助受限几何朗兰兹（以及命题 3.4.3.1）从谱 Eisenstein 级数转译过来。由引理 3.5.1.1 (1)，映射 $\mathrm{ps}\text{-}\Eis$ 把 $\delta_{\bT}$ 与 $\delta_{\bG}$ 分别在 $\Aut^{\unr}_{T,c}$ 与 $\Aut^{\unr}_{G,c}$ 上的作用衔接起来。此外，诱导映射
\[
  \mathrm{ps}\text{-}\Eis : \Aut^{\unr}_{T,c}[\delta_{\bT}^{-1}]
  \longrightarrow \Aut^{\unr}_{G,c}[\delta_{\bG}^{-1}]
\]
由定理 3.5.2.1 是满射。

这意味着，对我们的给定\footnote{需要说明：这对任何紧支撑自守函数都成立，但若该形式另有 Langlands 参数，则该论断可能基本上是空的。}自守函数 $\psi$，存在整数 $n\ge 0$，使得 $\delta_{\bG}^n\cdot\psi=\mathrm{ps}\text{-}\Eis(\psi_0)$，其中 $\psi_0\in\Aut^{\unr}_{T,c}$。

令 $\lambda\in\e$ 为 $\delta_{\bG}$ 在平凡局部系统 $\triv\in\LS^{\arthm}_{\bG}$ 处的值。注意 $(\delta_{\bG}-\lambda)\in\mathfrak m_{\triv}\subseteq H^0(\Exc)$，故对 $m\gg 0$ 有
\[
  (\delta_{\bG}-\lambda)^m\cdot\psi=0.
\]

依假设，$\lambda\ne 0$。因此可以找到多项式 $q(t)\in\e[t]$ 使
\[
  q(t)\cdot t^n \equiv 1 \pmod{(t-\lambda)^m}.
\]

于是我们显然得到
\[
  \mathrm{ps}\text{-}\Eis(q(\delta_{\bT})\cdot\psi_0)
  = q(\delta_{\bG})\cdot\mathrm{ps}\text{-}\Eis(\psi_0)
  = q(\delta_{\bG})\cdot\delta_{\bG}^n\cdot\psi
  = \psi.
\]

这就完成了论证。

\subsection{定理 3.5.2.1 的玩具模型}

我们现在给出一个更简单的设定，在其中定理 3.5.2.1 的某种形式成立。我们最终会把定理 3.5.2.1 的证明归约到这个特殊情形。我们所考虑的特殊情形，是关于 Grothendieck--Springer 消解的一个标准结果。

\pnumt{类比。}{}作为类比，我们把图表
\[
\begin{tikzcd}[column sep=huge, row sep=large]
  & \LS^{\restr}_{\bB} \arrow[dl, "q"{above}] \arrow[dr, "p"{above}] & \\
  \LS^{\restr}_{\bT} & & \LS^{\restr}_{\bG}
\end{tikzcd}
\]
换成图表：
\begin{equation}
\begin{tikzcd}[column sep=huge, row sep=large]
  & \mathbf{B}\check B \arrow[dl] \arrow[dr] & \\
  \mathbf{B}\check T & & \mathbf{B}\check G.
\end{tikzcd}
\end{equation}

代替 Frobenius，我们考虑 (3.14) 中的每一项连同其恒等自同态。取该映射之下的不动点，我们得到图表：
\[
\begin{tikzcd}[column sep=huge, row sep=large]
  & \check B^{\mathrm{ad}}/\check B \arrow[dl] \arrow[dr] & \\
  \check T^{\mathrm{ad}}/\check T & & \check G^{\mathrm{ad}}/\check G
\end{tikzcd}
\]

这里对代数群 $H$，$H^{\mathrm{ad}}/H$ 表示 $H$ 通过共轭作用在自身上所得的（叠）商；我们提醒\footnote{这里的同构存在一些符号含糊。通常这无关紧要，就当前的讨论而言亦然。但它在本文后面会有影响。我们在 §4.4.2 中澄清隐含的符号约定。}，$H^{\mathrm{ad}}/H\simeq(\mathbf{B}H)^{S^1}=(\mathbf{B}H)^{\id=\id}$。我们注意到，自然映射 $\OO_{H^{\mathrm{ad}}/H}\to\omega_{H^{\mathrm{ad}}/H}$ 是同构。

因此，§3.2 的形式体系给出典范映射：
\begin{equation}
  \Eissoyt : \Gamma\bigl(\check T^{\mathrm{ad}}/\check T,\OO_{\check T^{\mathrm{ad}}/\check T}\bigr)
  \longrightarrow
  \Gamma\bigl(\check G^{\mathrm{ad}}/\check G,\OO_{\check G^{\mathrm{ad}}/\check G}\bigr).
\end{equation}

\begin{lem}
(3.15) 的左端与右端都集中在上同调次数 $0$，且映射 $\Eissoyt$ 是满射。
\end{lem}

\begin{proof}
最直接的证明如下。我们把 $\Rep(\check T)$ 与 $\bigoplus_{\lambda\in\bar\Lambda}\Vect$ 等同，把 $\Rep(\check G)$ 与 $\bigoplus_{\lambda\in\bar\Lambda^+}\Vect$ 等同；这里 $\oplus$ 表示 $\DGCat_{\mathrm{cont}}$ 上的余积，并且我们已经默认选取了各不可约表示的同构类的代表。在此等同下，我们得到典范同构：
\begin{equation}
\begin{aligned}
  \bigoplus_{\lambda\in\bar\Lambda}k
  &\simeq \tr_{\Rep(\check T)}(\id)\,
    \bigl(\simeq\Gamma(\check T^{\mathrm{ad}}/\check T,\OO_{\check T^{\mathrm{ad}}/\check T})\bigr),\\
  \bigoplus_{\lambda\in\bar\Lambda^+}k
  &\simeq \tr_{\Rep(\check G)}(\id)\,
    \bigl(\simeq\Gamma(\check G^{\mathrm{ad}}/\check G,\OO_{\check G^{\mathrm{ad}}/\check G})\bigr).
\end{aligned}
\end{equation}

这些恒等式显然蕴涵 (3.15) 的两端都集中在次数 $0$。我们用 $e_{\check T,\bar\lambda}$（$\bar\lambda\in\bar\Lambda$）与 $e_{\check G,\bar\lambda}$（$\bar\lambda\in\bar\Lambda^+$）记由所显示的这些同构给出的向量空间的基向量。由 Borel--Weil--Bott，对支配的 $\bar\lambda$，映射 $\Rep(\check T)\to\Rep(\check G)$ 把 $\ell^{w_0(\bar\lambda)}\in\Rep(\check T)^{\vee}$（对应于 $w_0(\bar\lambda)$ 的一维表示）送到 $V^{\lambda}\in\Rep(\check G)$（最高权为 $\lambda$ 的表示）% [Tutor Note: 原文此处写作 V^λ、最高权写作 λ，与本段其余地方的 λ̄∈Λ̄ 记号不一致；译文照录。]
，从而把 $e_{\check T,w_0(\bar\lambda)}$ 送到 $e_{\check G,\bar\lambda}$；由此得到满射性。
\end{proof}

\begin{remark}
我们注意到（尽管我们并不需要这一点），在 (3.16) 中，同构 $\bigoplus_{\lambda\in\bar\Lambda^+}k\simeq\Gamma(\check G^{\mathrm{ad}}/\check G,\OO_{\check G^{\mathrm{ad}}/\check G})$ 把 $e_{\check G,\bar\lambda}$（记号同前）送到与 $\check G$ 的表示 $V^{\lambda}$ 相应的迹函数，对 $\check T$ 类似。因此，复合
\[
  \bigoplus_{\lambda\in\bar\Lambda}k
  \simeq \Gamma(\check T^{\mathrm{ad}}/\check T,\OO_{\check T^{\mathrm{ad}}/\check T})
  \xrightarrow{\ \Eissoyt\ }
  \Gamma(\check G^{\mathrm{ad}}/\check G,\OO_{\check G^{\mathrm{ad}}/\check G})
  \simeq \mathrm{Fun}(\check T)^{W}
  = \Bigl(\bigoplus_{\lambda\in\bar\Lambda}k\Bigr)^{W}
\]
可由 Weyl 特征标公式（以及 Borel--Weil--Bott）显式算出。
\end{remark}

\pnum{我们最后记录一个变体。}

注意我们有交换图：
\[
\begin{tikzcd}[column sep=huge, row sep=large]
  & \check B^{\mathrm{ad}}/\check B \arrow[dl] \arrow[dr] & \\
  \check T^{\mathrm{ad}}/\check T \arrow[dr] & & \check G^{\mathrm{ad}}/\check G \arrow[dl] \\
  & \check T/\!/W &
\end{tikzcd}
\]
其中 $\check T/\!/W=\Spec(\e[\bar\Lambda]^W)$ 是 GIT 商，左下方的映射是同义反复的（tautological），右下方的映射是标准的特征多项式映射（由该图表唯一刻画）。

因此，$\Eissoyt$ 是 $\mathrm{Fun}(\check T/\!/W)$-模的映射。引理 3.7.1.1 于是说明 $\Eissoyt$ 是 $\mathrm{Fun}(\check T/\!/W)$-模的满态射，故我们得到：

\begin{cor}
对任何 $g\in\Gamma(\check T/\!/W,\OO_{\check T/\!/W})$，映射
\[
  \Eissoyt : \Gamma\bigl(\check T^{\mathrm{ad}}/\check T,\OO_{\check T^{\mathrm{ad}}/\check T}\bigr)[g^{-1}]
  \longrightarrow
  \Gamma\bigl(\check G^{\mathrm{ad}}/\check G,\OO_{\check G^{\mathrm{ad}}/\check G}\bigr)[g^{-1}]
\]
是满射。
\end{cor}

% ============================================================
%  ch04a：§4 开头、4.1 基点与 Weil 群、4.2 伴随商（原文 pp.25–27）
%  起止：自 “4. Grothendieck–Springer theory for LS^arthm” 一节开头，
%        至 “4.3. Non-resonance” 标题之前（§4.3 起属 ch04b）。
% ============================================================
\section{$\LS^{\arthm}$ 的 Grothendieck--Springer 理论}

本节的目的是证明定理 3.5.2.1。由于该定理完全落在谱侧，本节通篇默认在域 $\e$ 上工作。

\subsection{基点与 Weil 群记号}

下面我们取定一个几何点 $x_0 \in X(k)$ 作为基点，它将在我们的基本群中充当基点的角色；这里提醒一句：$k = \Fb$。

我们建议读者对自己宽容一些，假定 $x_0$ 定义在 $\F_q$ 上；此时 §4.1 余下的材料基本上可以一概略去不看。

话虽如此，为使 $X_0$ 没有有理点的情形也能覆盖，我们在这里收录了一些技术性材料。

\pnum{}令 $\widetilde{X}$ 表示 $X$ 以 $x_0$ 为基点的万有覆盖；按定义，$\widetilde{X}$ 连通、在 $X$ 上 pro-有限 étale，配有 $x_0$ 的一个提升 $\widetilde{x}_0$，并在所有这类数据中是初始的。注意 $\widetilde{X}$ 同时也是 $X_0$ 的万有覆盖，因此存在 $\pi_1^{\et}(X_0)$（$= \pi_1^{\et}(X_0,x_0)$）在 $\widetilde{X}$ 上的一个同义反复的作用（它把 $\widetilde{X}$ 实现为 $X$ 上的 $\pi_1^{\et}(X)$-torsor、以及 $X_0$ 上的 $\pi_1^{\et}(X_0)$-torsor，其中每个 torsor 都理解为对 pro-étale 拓扑局部平凡）。

令 $\Frob_X : X \to X$ 表示几何 Frobenius 映射。我们一次性选定点 $\Frob_X(x_0)$ 到 $\widetilde{X}$ 的一个提升。容易看出，存在唯一的映射 $\widetilde{\Frob}_X$ 使下图交换
\[
\begin{tikzcd}[column sep=large, row sep=large]
  \widetilde{X} \arrow[r, "\widetilde{\Frob}_X"] \arrow[d, "\pi"'] & \widetilde{X} \arrow[d, "\pi"] \\
  X \arrow[r, "\Frob_X"] & X
\end{tikzcd}
\]
并把 $\widetilde{x}_0$ 送到我们选定的 $\Frob_X(x_0)$ 的提升（今后记作 $\widetilde{\Frob}_X(\widetilde{x}_0)$）。

这一选择还（与此相关地）定义了 $\widehat{\ZZ}$ 在 $\mathrm{Sch}_{/X_0}$ 中的 $\widetilde{X}$ 上的一个作用；生成元 $-1 \in \widehat{\ZZ}$ 的逆\footnote{我们指出，按本文的约定，当我们把 $\widehat{\ZZ} \simeq \mathrm{Gal}(\F_q)$ 等同起来时，生成元 $-1 \in \widehat{\ZZ}$ 的逆对应的是\emph{算术} Frobenius。}通过一个映射 $\gamma : \widetilde{X} \to \widetilde{X}$ 作用，该映射的特征是：在 $k$ 上是 Frobenius 半线性的，并且 $\gamma \circ \widetilde{\Frob}_X$ 是 $\widetilde{X}$ 的绝对 Frobenius。

这份数据定义了映射 $\mathcal{W}_X \to \ZZ$ 的一个分裂。我们令 $F \in \mathcal{W}_X$ 表示 $1 \in \ZZ$ 在该分裂下的像。按定义，$F^{-1} \in \mathcal{W}_X \subseteq \pi_1^{\et}(X_0)$ 通过上面记作 $\gamma$ 的映射作用在 $\widetilde{X}$ 上。

对 $g \in \pi_1^{\et}(X)$，我们令 ${}^F g := \Ad_F(g)$。注意点 $\widetilde{\Frob}_X(\widetilde{x}_0)$ 的选取给出了同构 $\pi_1^{\et}(X,\Frob_X(x_0)) \simeq \pi_1^{\et}(X,x_0)$，而复合
\[
  \pi_1^{\et}(X,x_0) \xrightarrow{\ \pi_1^{\et}(\Frob_X)\ } \pi_1^{\et}(X,\Frob_X(x_0)) \simeq \pi_1^{\et}(X,x_0)
\]
就是映射 $g \mapsto {}^F g$\footnote{确实，按 $\pi_1^{\et}(\Frob_X)(g)$ 的定义，有 $\widetilde{\Frob}_X(g\widetilde{x}_0) = \pi_1^{\et}(\Frob_X)(g) \cdot \widetilde{\Frob}_X(\widetilde{x}_0)$。
  又按 $F = \gamma^{-1}$ 的定义，有 $F^{-1}\cdot\widetilde{\Frob}_X(g\widetilde{x}_0) = \Phi_{\widetilde{X}}(g\widetilde{x})$，其中 $\Phi_{\widetilde{X}}$ 是绝对 Frobenius。由函子性，绝对 Frobenius 是带 $\pi_1^{\et}(X)$-作用的空间之间的映射，故 $\Phi_{\widetilde{X}}(g\widetilde{x}_0) = g \cdot \Phi_{\widetilde{X}}(\widetilde{x}_0)$。
  与我们先前的等式比较，可见该表达式等于 $F^{-1}\pi_1^{\et}(\Frob_X)(g) \cdot \widetilde{\Frob}_X(\widetilde{x}_0) = F^{-1}\pi_1^{\et}(\Frob_X)(g)F \cdot F^{-1}\widetilde{\Frob}_X(\widetilde{x}_0) = F^{-1}\pi_1^{\et}(\Frob_X)(g)F \cdot \Phi_{\widetilde{X}}(\widetilde{x}_0)$。因此 $g = \pi_1^{\et}(\Frob_X)(g)F$，断言得证。% [Tutor Note: 原文第二句作 “Φ_X̃(gx̃)”，据上下文应为 “Φ_X̃(gx̃_0)”]
  }，也就是说，我们有
\begin{equation}
  \pi_1^{\et}(\Frob_X)(g) = {}^F g.
\end{equation}

\pnum{点 $\widetilde{\Frob}_X(\widetilde{x}_0)$ 的选取还定义了一个同构
\[
  x_0^* \simeq \Frob_X(x_0)^* : \mathrm{Lisse}(X)^\heartsuit \to \Vect^\heartsuit
\]
（$\e$-线性对称幺半函子之间）。事实上，我们断言它来自对称幺半 DG 函子的同构
\begin{equation}
  x_0^* \simeq \Frob_X(x_0)^* : \qLisse(X) \to \Vect.
\end{equation}
确实，只要 $\qLisse(X)$ 是其核心的导出范畴，这就形式地成立。由 [2] 定理 E.2.8，$X \neq \mathbb{P}^1$ 时正是如此。稍初等一些（并且允许亏格 $0$），我们取 $U \subseteq X$ 为包含 $x_0$ 与 $\Frob_X(x_0)$ 的仿射开集，此时由（更简单的）定理 E.2.8 (a)，$\qLisse(U)$ 是其核心的导出范畴。}

\begin{remark}
下面的注不会被用到。眼下暂且让 $k$ 是\emph{任意}代数闭域（而不只是 $\overline{\F}_q$）。令 $Y/k$ 为有限型的连通概形，$y_1, y_2 \in Y(k)$ 为两个点。此时自然要问：是否存在对称幺半 DG 函子的同构
\[
  y_1^* \simeq y_2^* : \qLisse(Y) \to \Vect ?
\]
我们断言这是对的。确实，光滑连通曲线的情形我们上面已经有效地处理过了。任意连通曲线的情形可以按显然的方式归结为正规化（把奇异曲线各不可约分支之间的交点当作路标来引路）。最后，一般情形由如下初等论证得到：存在一条连通曲线 $C$ 与一个映射 $C \to Y$，使 $y_1$ 与 $y_2$ 都落在其像中。（人们不禁要问：在这个导出设定中是否存在一个纯 Tannaka 的论证？）
\end{remark}

\subsection{伴随商}

以下令 $H$ 为仿射代数群。

\pnum{令 $\LS^{\restr,\circ}_H$ 表示 $\LS^{\restr}_H$ 的\emph{单位}连通分支，即包含 $X$ 上平凡 $H$-局部系统的那个连通分支。我们提醒（见 [2] 命题 3.7.2）：$\LS^{\restr}_H$ 参数化（在某种确切的意义下）$X$ 上半单化平凡的 $H$-局部系统；下文中我们把这些称为\emph{幂幺} $H$-局部系统。

于是我们置
\[
  \LS^{\arthm,\circ}_H := \LS^{\arthm}_H \times_{\LS^{\restr}_H} \LS^{\restr,\circ}_H .
\]
换言之，$\LS^{\arthm,\circ}_H$ 是 Frobenius 作用在 $\LS^{\restr,\circ}_H$ 上的不动点；可以把它设想为参数化\emph{几何幂幺}的 Weil $H$-局部系统。}

\pnum{存在一个典范映射 $\mathbf{B}H \to \LS^{\restr,\circ}_H$，对应于平凡 $H$-局部系统。该映射是 Frobenius 等变的，其中 Frobenius 在 $\mathbf{B}H$ 上平凡地作用。

过渡到 Frobenius 不动点，我们得到映射
\begin{equation}
  \tau_H : H^{ad}/H \to \LS^{\arthm,\circ}_H .
\end{equation}}

\begin{remark}
粗略地说，复合 $H \to H^{ad}/H \to \LS^{\arthm,\circ}_H$ 把 $h \in H$ 送到由
\[
  \mathcal{W}_X \twoheadrightarrow \mathcal{W}_X/\pi_1^{\et}(X) \xrightarrow{\ \mathrm{Fr}\ } \ZZ \xrightarrow{\ 1 \mapsto h\ } H
\]
定义的 Weil 表示 $\mathcal{W}_X \to H$。由于我们只考虑 Weil 表示的\emph{群胚}，该映射如所愿地通过伴随商分解。
\end{remark}

\begin{remark}
下文的一个主要想法是：$\tau_H$ 离同构并不太远。为给后文提供动机，我们指出 $\tau_H$ 成为同构的如下障碍。粗略地说（例如在域值点的层面上），敌人显然是那些几何幂幺、且几何非平凡的 Weil 表示。

设 $H = GL_2$。令 $\lambda \in \e^\times$，并令 $\sigma_\lambda$ 表示 Frobenius 以乘 $\lambda$ 作用（而 $\pi_1^{\et}(X)$ 平凡作用）的一维 Weil 群表示。扩张 $0 \to \sigma_\lambda \to \sigma \to \sigma_1 \to 0$ 由 $\ZZ$ 的适当群上同调分类，即由复形
\[
  \underline{\Hom}_{\e[\ZZ]\text{-}\mathrm{mod}}\bigl(C_{\et,\bullet}(X), \sigma_\lambda\bigr)
\]
的 $H^1$ 分类。这里 $C_{\et,\bullet}(X)$ 是 $X$ 的\emph{\et{} 同调}的复形，$\ZZ$-作用的生成元在 \et{} 同调上按几何 Frobenius 作用。若 $\lambda \neq 1$，容易看出有正合列
\[
  0 \to H^1\bigl(\underline{\Hom}_{\e[\ZZ]\text{-}\mathrm{mod}}(C_{\et,\bullet}(X),\sigma_\lambda)\bigr) \to H^1_{\et}(X,\sigma_\lambda) \xrightarrow{\ \phi_X^{-1}-\id\ } H^1_{\et}(X,\sigma_\lambda),
\]
其中 $\phi_X$ 是作用在 $H^1_{\et}$ 上的几何 Frobenius。\footnote{注意，同调与上同调上的几何 Frobenius 是互为转置（即对偶）的态射。然而，如果我们把例如同调 $H_{\et,1}(X) = (\pi_1^{\et}(X)^{\mathrm{ab}})^{\wedge}_\ell \otimes \e$ 连同其几何 Frobenius（由 (4.1)，它对应于 $\Ad_F : \pi_1^{\et}(X) \to \pi_1^{\et}(X)$）看作 $\ZZ$-表示，那么上同调上的\emph{对偶} $\ZZ$-作用的生成元按\emph{算术} Frobenius 作用。这就解释了上式中出现的那个逆号。} 这里 $H^1_{\et}(X,\sigma_\lambda) = H^1_{\et}(X,\e)$，只是 Frobenius 作用取为标准作用的 $\lambda$ 倍。因此，若 $\lambda$ 是出现在 $H^1_{\et}(X,\e)$ 中的 Frobenius 特征值，我们就得到所需类型的\emph{几何}非平凡扩张。
\end{remark}

\pnumt{分裂。}{}通过限制到 $x_0$，我们得到映射 $\LS^{\restr}_H \to \mathbf{B}H$。由 §4.1.2\footnote{确切地说，当 $x_0$ 不是 $\F_q$-有理点时，这种等变性依赖于辅助选择。}，该映射把 Frobenius 与恒等映射缠结在一起，于是在不动点上我们得到映射
\[
  \LS^{\arthm}_H \to H^{ad}/H .
\]
我们把这个映射记作 $\chi_H = \chi_{H,x_0}$，并类似地记它在 $\LS^{\arthm,\circ}_H$ 上的限制。

由构造，复合
\[
  H^{ad}/H \xrightarrow{\ \tau_H\ } \LS^{\arthm}_H \xrightarrow{\ \chi_H\ } H^{ad}/H
\]
是恒等映射。

% ============================================================
%  ch04b：§4.3 非共振、§4.4 定理 4.3.3.1 的证明（原文 pp.27 末–p.31 中段）
%  起止：自 “4.3. Non-resonance” 标题起，至 §4.4.6 末尾 “as was desired” 止；
%        其后 “4.5. Setup for the proof of Theorem 3.5.2.1” 起属 ch04c。
% ============================================================
\subsection{非共振}

\pnum{定义集合 $\mathcal{R}_X \subseteq \e^\times$ 为（几何）Frobenius 作用在 $H^1_{\et}(X,\e) \times H^2_{\et}(X,\e)$ 上的特征值所构成的集合。}

\begin{remark}
由曲线的 Weil 猜想，$1 \notin \mathcal{R}_X$。此外，$q$ 总是落在 $\mathcal{R}_X$ 之中（不过这一点与我们眼下的关切关系不大）。
\end{remark}

\pnum{}令 $H$ 为仿射代数群。

令 $V$ 为 $H$ 的有限维表示，并令 $\rho_V : H \to GL(V)$ 为相应的同态。令 $\ch_V : H^{ad}/H \times \mathbf{A}^1 \to \mathbf{A}^1$ 为使下图交换的映射：
\[
\begin{tikzcd}[column sep=large, row sep=large]
  H \times \mathbf{A}^1 \arrow[r, "\rho_V \times \id"] \arrow[d]
    & GL(V) \times \mathbf{A}^1 \arrow[d, "{(g,\lambda) \mapsto \det(g-\lambda\cdot\id_V)}"] \\
  H^{ad}/H \times \mathbf{A}^1 \arrow[r, "\ch_V"]
    & \mathbf{A}^1.
\end{tikzcd}
\]
显式地，对 $h \in H$ 与 $\lambda \in \mathbf{A}^1$，$\ch_V([h],\lambda)$ 就是 $\rho_V(h)$ 的特征多项式在 $\lambda$ 处的取值。

我们于是令 $(H^{ad}/H)^{\nonres} \subseteq H^{ad}/H$ 表示由满足 $\prod_{\lambda \in \mathcal{R}_X} \ch_{\mathfrak{h}}([h],\lambda) \neq 0$ 的共轭类 $[h]$ 构成的开集，其中 $\mathfrak{h}$ 是 $H$ 的伴随表示。显式地，$(H^{ad}/H)^{\nonres}$ 是那些使矩阵 $\rho_{\mathfrak{h}}(h) \in GL(\mathfrak{h})$ 在 $\mathcal{R}_X$ 中没有特征值的共轭类 $[h]$ 的集合。我们指出，开浸入 $(H^{ad}/H)^{\nonres} \hookrightarrow H^{ad}/H$ 显然是仿射的。

\begin{remark}
由注 4.3.1.1，$[1] \in (H^{ad}/H)^{\nonres}$。
\end{remark}

\begin{notation}
对任何配有到 $(H^{ad}/H)^{\nonres}$ 的结构映射的叠 $\mathcal{Y}$，我们令 $\mathcal{Y}^{\nonres} := \mathcal{Y} \times_{H^{ad}/H} (H^{ad}/H)^{\nonres}$。我们特别在 $\mathcal{Y} = \LS^{\arthm,\circ}_H$、配以结构映射 $\chi_H$ 的情形使用这个记号。
\end{notation}

\pnumt{主要几何结果。}{}\ 下面的结果把算术局部系统与伴随商作比较：

\begin{thm}
来自 (4.3) 的映射
\[
  \tau_H : (H^{ad}/H)^{\nonres} \to \LS^{\arthm,\circ,\nonres}_H
\]
是同构。
\end{thm}

该结果的证明是 §4.4 的内容。

\subsection{定理 4.3.3.1 的证明}

\pnumt{映射是同构的判别准则。}{}\ 我们首先作如下观察：

\begin{lem}
设 $f : \mathcal{Y}_1 \to \mathcal{Y}_2$ 是局部几乎有限型的代数叠（在代数闭域 $\e$ 上）之间的态射。

则 $f$ 是同构，当且仅当：
\begin{enumerate}[label=(\roman*)]
  \item $f$ 是形式 étale 的，即其余切复形 $\Omega^1_{\mathcal{Y}_1/\mathcal{Y}_2} \in \QCoh(\mathcal{Y}_1)$ 为零。
  \item 映射 $\mathcal{Y}_1(\e) \to \mathcal{Y}_2(\e)$ 是 (1-) 群胚的同构。
\end{enumerate}
这里我们要明确指出，条件 (ii) 可以拆成如下两个独立的条件：
\begin{enumerate}[label=(\text{ii}$_{\arabic*}$)]
  \item 对每个 $y_1 \in \mathcal{Y}_1(\e)$，自同构群的映射 $\Aut_{\mathcal{Y}_1(\e)}(y_1) \to \Aut_{\mathcal{Y}_2(\e)}(f(y_1))$ 是同构\footnote{我们强调，这里没有任何导出（derived）的余地：这是一对集合之间的映射。}。
  \item 对每个 $y_2 \in \mathcal{Y}_2(\e)$，存在 $y_1 \in \mathcal{Y}_1(\e)$ 以及同构 $f(y_1) \simeq y_2 \in \mathcal{Y}_2(\e)$。
\end{enumerate}
\end{lem}

\begin{proof}
只需证明：对每个局部几乎有限型、并配有映射 $S \to \mathcal{Y}_2$ 的仿射 $S$，映射 $S \times_{\mathcal{Y}_2} \mathcal{Y}_1 \to S$ 是同构。性质 (i) 与 (ii) 显然在这样的基变换下保持，于是我们可以假定 $\mathcal{Y}_2$ 是仿射概形。此外，众所周知 $\mathcal{Y}_1 \to \mathcal{Y}_2$ 是同构当且仅当 $\mathcal{Y}_1 \times_{\mathcal{Y}_2} \mathcal{Y}_2^{\mathrm{cl}} \to \mathcal{Y}_2^{\mathrm{cl}}$ 是同构；因此我们可以进一步假定 $\mathcal{Y}_2$ 是经典的。

此时 $\mathcal{Y}_1$ 是 $\e$-点上自同构群平凡的代数叠，从而是代数空间。又因 $\mathcal{Y}_1 \to \mathcal{Y}_2$ 是 étale 的，$\mathcal{Y}_1$ 也是经典的。现在 $f$ 是 radicial 态射（因为它局部有限型、且在 $\e$-点上单射），又是 étale 的，从而是开浸入。最后，由于 $f$ 在 $\e$-点上满射，它必定是同构。
\end{proof}

下面我们将对定理 4.3.3.1 中所考虑的映射 $\tau_H$ 验证上述这些假设。

\pnumt{约定、公式与符号。}{}\ 在继续之前，我们先确立某些重要的符号约定。粗略地说，要在概念上区分 $\tau_H([h])$ 与 $\tau_H([h^{-1}])$ 是困难的，但定理 4.3.3.1 确实能把它们区分开来，因此我们必须稍微更显式地解释该如何理解映射 $\tau_H$。（读者不妨跳过这段题外话，需要时再回头看。）

下面我们令 $\Frob_X : X \to X$ 表示几何 Frobenius 映射。

我们把 $X$ 上的 lisse Weil 层约定为如下的一对 $(\sigma,\alpha)$：其中 $\sigma \in \mathrm{Lisse}(X)^\heartsuit = \qLisse(X)^\heartsuit$，并配有同构 $\alpha : \sigma \xrightarrow{\ \sim\ } \Frob_X^*(\sigma)$。映射 $\alpha$ 的方向正是这里所说的“符号”。我们首先解释：为什么（以及在什么意义上）这个符号对于我们已有的约定而言是正确的。

注意 $V := x_0^*(\sigma)$ 是 $\pi_1^{\et}(X)$ 的一个表示 $\rho^{\mathrm{geom}}$。我们还得到一个同构
\begin{equation}
  V = x_0^*(\sigma) \xrightarrow{\ \alpha_{x_0}\ } \Frob_X(x_0)^*(\sigma) \xrightarrow{\ (4.2)\ } x_0^*(\sigma) = V,
\end{equation}
我们把它记作 $\rho(F)$。按 §4.1.1 的记号，可以同义反复地得到 $\rho(F) \circ \rho^{\mathrm{geom}}(g) = \rho^{\mathrm{geom}}(\pi_1^{\et}(\Frob_X)(g)) \circ \rho(F)$。由 (4.1)，我们可以把这个等式改写为
\[
  \rho(F)\rho^{\mathrm{geom}}(g)\rho(F)^{-1} = \rho^{\mathrm{geom}}(FgF^{-1}),
\]
于是我们得到 $\mathcal{W}_X$ 在 $V$ 上的表示，其中 $F$ 按 (4.4) 作用——倘若 $\alpha$ 的方向反了过来，我们就需要把 (4.4) 取逆。

类似地，对仿射代数群 $H$，一个 Weil $H$-局部系统指的是 $X$ 上的 $H$-局部系统 $\sigma_H$，连同同构 $\alpha : \sigma_H \xrightarrow{\ \sim\ } \Frob_X^*(\sigma_H)$（$H$-局部系统之间的）。作为推论，对 $h \in H$，$\tau_H([h])$ 的 $\sigma_H$ 是平凡的，而 $\alpha$ 由乘以 $h$ 给出。这就保证相应的 Weil 群表示 $\mathcal{W}_X \to H(\e)$ 通过 $\ZZ = \mathcal{W}_X/\pi_1^{\et}(X)$ 分解，并把生成元送到 $h$（正如它本来应该的那样）。

最后，对 lisse Weil 层 $(\sigma,\alpha)$，其上同调上自然的“几何” Frobenius 作用 $\phi_\sigma$ 由如下算子给出：
\begin{equation}
  C_{\et}(X,\sigma) \to C_{\et}(X,\Frob_X^*(\sigma)) \xrightarrow{\ \alpha^{-1}\ } C_{\et}(X,\sigma),
\end{equation}
其中第一个映射是同义反复的那个。

\begin{remark}
我们希望把上述材料的逻辑地位讲清楚。首先，我们论证了映射 $\alpha$ 应当按某个确定的方向来考虑。但这在某种程度上只是直觉性的论证，而非数学论证。毋宁说，我们只是把先前隐含的某个\footnote{说得更清楚些：$H^{ad}/H = (\mathbf{B}H)^{S^1}$ 有一个“回路反转”（loop reversal）的自同构，而我们需要消除这个自同构所带来的含糊。}约定明确化了（并说明它如何导向注 4.2.2.1 中非正式地暗示的那个定向）。从逻辑上讲，确立这一约定对于定理 4.3.3.1 的陈述是绝对必要的。
\end{remark}

\pnumt{$\tau_H$ 的形式 étale 性。}{}\ 我们将证明 $\tau_H|_{(H^{ad}/H)^{\nonres}}$ 是形式 étale 的。

首先注意，我们归结为验证切复形为零（例如，两侧的余切复形都是完美的）。此外，由于两侧都局部几乎有限型，我们可以在所有 $\e$-点的纤维上作验证。

一般地，对 $\sigma \in \LS^{\arthm}_H(\e)$，切复形可以如下计算\footnote{见 [2] 命题 2.2.2、§24.5.1。}：
\[
  T_{\LS^{\arthm}_H,\sigma} = C_{\et}(X,\mathfrak{h}_\sigma[1])^{\ZZ} = \mathrm{Ker}\bigl(\id - \phi_\sigma : C_{\et}(X,\mathfrak{h}_\sigma[1]) \to C_{\et}(X,\mathfrak{h}_\sigma[1])\bigr).
\]
这里 $\mathfrak{h}_\sigma$ 是由 $\sigma$ 诱导出的 $X$ 上的伴随 Weil 局部系统，$C_{\et}(X,\mathfrak{h}_\sigma[1])$ 是它的 étale 上同调复形（差一个平移），而我们取的是关于 Frobenius $\phi_\sigma$（来自 $\sigma$ 上的 Weil 结构）的作用的 $\ZZ$-不变量。

对 $h \in H(\e)$，切复形 $T_{H^{ad}/H,[h]} \in \Vect$ 是同伦核（即平移后的锥）：
\[
  T_{H^{ad}/H,[h]} = \mathrm{Ker}(\id - \Ad_h : \mathfrak{h}[1] \to \mathfrak{h}[1]),
\]
其中 $\Ad_h$ 是 $h$ 在 $\mathfrak{h}$ 上的伴随作用；更自然地说，把 $H^{ad}/H = \mathrm{Maps}(\mathbf{B}\ZZ,\mathbf{B}H)$ 写出来，该公式给出的是 $\ZZ$ 在 $\mathfrak{h}[1] = T_{\mathbf{B}H,\Spec(k)}$ 上作用的 $\ZZ$-不变量，其中生成元按 $\Ad_h$ 作用。（一个推论：我们看到，在上述对 $T_{H^{ad}/H,[h]}$ 的描述中本可以用 $\Ad_{h^{-1}}$ 代替 $\Ad_h$。）

另一方面，我们有
\[
  C_{\et}(X,\mathfrak{h}_{\tau_H([h])}[1]) = C_{\et}(X,\e) \otimes \mathfrak{h}[1],
\]
其 Frobenius 为
\[
  \phi_{\tau_H([h])} = \phi_X \otimes \Ad_{h^{-1}},
\]
其中逆号的出现是因为 (4.5) 中出现了 $\alpha^{-1}$。

Frobenius 等变的映射 $\e = H^0(X,\e) \to C_{\et}(X,\e)$（Frobenius 在来源上平凡地作用）诱导出交换图
\[
\begin{tikzcd}[column sep=3.4cm, row sep=large]
  \mathfrak{h}[1] = \e \otimes \mathfrak{h}[1] \arrow[r, "{\id - \Ad_{h^{-1}}}" description] \arrow[d]
    & \mathfrak{h}[1] = \e \otimes \mathfrak{h}[1] \arrow[d] \\
  C_{\et}(X,\e) \otimes \mathfrak{h}[1] \arrow[r, "{\id - \phi_X \otimes \Ad_{h^{-1}}}" description]
    & C_{\et}(X,\e) \otimes \mathfrak{h}[1].
\end{tikzcd}
\]
沿着各行过渡到（同伦）核，就得到 $\tau_H$ 在 $[h]$ 处的微分\footnote{要看出这一点，把 $H^{ad}/H$ 看作 $\Spec(k)$ 上算术局部系统的模叠，再把上面关于 $\LS^{\arthm}_H$ 的讨论相应地用上去即可。}。

因此我们看到：$\tau_H$ 在 $[h]$ 处是形式 étale 的，当且仅当
\[
  \id - \phi_X \otimes \Ad_{h^{-1}} : \tau^{\geq 1}C_{\et}(X,\e) \otimes \mathfrak{h} \to \tau^{\geq 1}C_{\et}(X,\e) \otimes \mathfrak{h}
\]
是同构；等价地，当且仅当上同调上的诱导映射
\begin{align*}
  \id - \phi_X \otimes \Ad_{h^{-1}} &: H^1_{\et}(X,\e) \otimes \mathfrak{h} \to H^1_{\et}(X,\e) \otimes \mathfrak{h},\\
  \id - \phi_X \otimes \Ad_{h^{-1}} &: H^2_{\et}(X,\e) \otimes \mathfrak{h} \to H^2_{\et}(X,\e) \otimes \mathfrak{h}
\end{align*}
都是同构。显然这恰好在 $1$ 不是 $\phi_X \otimes \Ad_{h^{-1}}$ 的特征值时发生，而这又恰好在 $1$ 不能写成 $\lambda \cdot \mu$（其中 $\lambda$ 是 $\phi_X$ 的特征值、$\mu$ 是 $\Ad_{h^{-1}} = \Ad_h^{-1}$ 的特征值）时发生，也就是说，当 $\Ad_h$ 没有特征值落在 $\mathcal{R}_X$ 中时发生。这正是 $[h]$ 落在 $(H^{ad}/H)^{\nonres}$ 中的定义条件，于是我们得到断言。

\pnumt{稳定化子。}{}\ 接下来我们验证引理 4.4.1.1 中的条件 (ii$_1$)。事实上，这一步是显然的，我们在此也绝不会用到非共振的微妙之处。我们在这里把论证明确写出来：

设 $\sigma \in \LS^{\arthm}_H(\e)$。按定义，$\sigma$ 提升为连续的 Weil 群表示 $\rho : \mathcal{W}_X \to H(\e)$，该提升在共轭意义下良定义。此时 $\Aut_{\LS^{\arthm}_H(\e)}(\sigma)$ 是 $\rho$ 的像在 $H(\e)$ 中的稳定化子。

类似地，$H^{ad}/H$ 中的一个 $\e$-点提升为某个 $h \in H(\e)$，而 $\Aut_{H^{ad}/H(\e)}([h])$ 是 $h$ 的稳定化子。

现在对 $h \in H$，$\tau_H([h])$ 就是 Weil 群表示 $\mathcal{W}_X \to \ZZ \xrightarrow{1 \mapsto h} H(\e)$，其稳定化子显然与 $h$ 的稳定化子一致。

\pnumt{提升同构类：设定。}{}\ 最后我们验证条件 (ii$_2$)。设 $\sigma \in \LS^{\arthm,\circ,\nonres}_H(\e)$。我们把 $\sigma$ 提升为连续表示 $\rho : \mathcal{W}_X \to H(\e)$。令 $\rho_0 : \pi_1^{\et}(X) \to H(\e)$ 表示 $\rho$ 在几何基本群上的限制；我们的任务是证明 $\rho_0$ 是平凡的。

这是一个具体的线性代数问题；我们在下面把细节写清楚。我们沿用 §4.1.1 的记号（特别地，$F \in \mathcal{W}_X$ 以及 $g \mapsto {}^F g$）。此外，我们引入更多记号：
\begin{itemize}
  \item 令 $H_\circ \subseteq H$ 表示 $\mathrm{Image}(\rho_0)$ 的 Zariski 闭包\zy{原文这里用小圆圈作下标（$H_\circ$），全文一贯如此（如 $H_\circ^{\mathrm{ab}}$），并非上标零；译文按原样排作下标小圆圈。[来源：A/原文 pp.30--31 版面]}。
  \item 令 $\theta : H_\circ \xrightarrow{\ \sim\ } H_\circ$ 表示 $\rho(F)$ 的伴随作用，即 $\theta(h) = \Ad_{\rho(F)}(h)$。
\end{itemize}

在上述记号下，注意我们有
\begin{equation}
  \rho_0({}^F g) = \theta(\rho_0(g)) \qquad \text{对一切 } g \in \pi_1^{\et}(X).
\end{equation}
现在我们对假设再作如下几点观察。

第一，$\rho_0$ 定义了 $\LS^{\restr,\circ}_H \subseteq \LS^{\restr}_H$ 中的一个点，这意味着由 [2] 命题 3.7.2，$\rho_0$ 通过 $H$ 的一个幂幺子群分解。等价地说，$H_\circ$ 是幂幺的。

第二，注意非共振条件意味着 $\Ad_{\rho(F)} = \mathrm{Lie}(\theta) : \mathfrak{h} \to \mathfrak{h}$ 在 $\mathcal{R}_X$ 中没有特征值。

\pnumt{提升同构类：证明。}{}\ 在上述记号下，我们的任务是证明 $H_\circ$ 是平凡的。由幂幺性，只需证明其交换化 $H_\circ^{\mathrm{ab}}$ 是平凡的。令 $V := \mathrm{Lie}(H_\circ^{\mathrm{ab}})$；由于 $H_\circ^{\mathrm{ab}}$ 是交换幂幺群，我们滥用记号，把它与它的 Lie 代数（即与之相伴的 $\e$-概形）等同起来。

由函子性，$\theta$ 诱导出 $V$ 的一个自同构，我们仍记作 $\theta$。设 $V \neq 0$；则对转置 $\theta^\vee : V^\vee \to V^\vee$，存在特征向量 $\mu \in V^\vee$；我们令 $\lambda \in \e^\times$ 表示其特征值。注意由非共振假设，$\lambda \notin \mathcal{R}_X$。

于是我们得到一个连续同态
\[
\begin{tikzcd}[column sep=large, row sep=large]
  \pi_1^{\et}(X) \arrow[r, "\rho_0"] \arrow[rrr, bend right=22, "\bar{\rho}_0"']
    & H_\circ(\e) \arrow[r]
    & V \arrow[r, "\mu"]
    & \e
\end{tikzcd}
\]
由 (4.6)，它满足
\begin{equation}
  \bar{\rho}_0({}^F g) = \lambda \cdot \bar{\rho}_0(g).
\end{equation}
我们还指出，$\mathrm{Ker}\bigl(H_\circ(\e) \to V \to \e\bigr)$ 是 $H_\circ$ 的某个代数子群的 $\e$-点集，于是按后者的定义，同态 $\bar{\rho}_0$ 必须是非平凡的。

现在 $\bar{\rho}_0$ 延拓为 $\e$-线性映射 $H_{\et,1}(X,\e) \to \e$，即它来自一个非零的上同调类 $\eta \in H^1_{\et}(X,\e)$。由于 $g \mapsto {}^F g$ 在 $H_{\et,1}(X,\e)$ 上诱导出（几何）Frobenius（见 (4.1)），(4.7) 意味着
\[
  \phi_X(\eta) = \lambda \cdot \eta.
\]
这与非共振假设矛盾，于是我们断定 $V = 0$，此即所欲证者。

% ============================================================
%  ch04c：§4.5 定理 3.5.2.1 证明的设定、§4.6 L-值与迹
%  （原文 pp.31 中段–39，起于 §4.5，止于引理 4.6.6.4 的证明）
% ============================================================
\subsection{定理 3.5.2.1 证明的设定}

我们开始证明定理 3.5.2.1。

\pnum{回忆我们的目标是定义交换图 (3.11) 并验证它的若干性质。}

我们首先定义函数 $\delta\colon\bT\to\A^1$ 如下：
\[
  \delta(t) := \prod_{\lambda\in\mathcal{R}_X}\ch_{\bg}(t,\lambda),\qquad t\in\bT,
\]
其中我们使用 §4.3.2 的记号，并把 $\bg$ 通过伴随作用看作 $\bT$ 的表示。更显式地，我们有：
\[
  \delta(t) = \prod_{\lambda\in\mathcal{R}_X}(1-\lambda)^{\dim(\bT)}\prod_{\check\alpha\in\check\Delta}(\check\alpha(t)-\lambda),
\]
其中我们把 $\check\alpha$ 看作映射 $\bT\to\mathbb{G}_m\subset\A^1$。

显然 $\delta$ 是 $W$-不变态射，因此诱导一个映射 $\bT/\!/W\to\A^1$；我们仍用 $\delta$ 记这个函数。

\pnum{现在我们构造一个交换图：}
\begin{equation}
\begin{tikzcd}[column sep=0.9em, row sep=1.5em]
& \LS^{\arthm,\circ}_{\bB} \arrow[dl] \arrow[d, "\chi_{\bB}"] \arrow[dr] & \\
\LS^{\arthm,\circ}_{\bT} \arrow[d, "\chi_{\bT}"]
  & \bB^{\ad}/\bB \arrow[dl] \arrow[dr]
  & \LS^{\arthm,\circ}_{\bG} \arrow[d, "\chi_{\bG}"] \\
\bT^{\ad}/\bT \arrow[dr] & & \bG^{\ad}/\bG \arrow[dl] \\
& \bT/\!/W \arrow[d, "\delta"] & \\
& \A^1 &
\end{tikzcd}
\end{equation}

我们现在这样定义 $\delta_{\bG}$：它在 $\LS^{\arthm,\circ}_{\bG}$ 上的限制由 (4.8) 中出现的唯一映射 $\LS^{\arthm,\circ}_{\bG}\to\A^1$ 给出，而它在 $\LS^{\arthm}_{\bG}\setminus\LS^{\arthm,\circ}_{\bG}$ 上的限制\footnote{我们提醒，
$\LS^{\restr,\circ}_{\bG}\subseteq\LS^{\restr}_{\bG}$ 是一个连通分支，因此先在
$\LS^{\arthm,\circ}_{\bG}:=\LS^{\arthm}_{\bG}\times_{\LS^{\restr}_{\bG}}\LS^{\restr,\circ}_{\bG}$
上定义这个函数、并在别处令其为零，是合法的。}恒为零。我们用完全同样的方式定义 $\delta_{\bT}$，只需把上一句中的 $\bG$ 处处换成 $\bT$。

下面我们验证映射对 $(\delta_{\bG},\delta_{\bT})$ 满足定理 3.5.2.1 的结论。

\pnum{首先，由 (4.8) 显然可知交换图 (3.11) 存在。}

\pnum{其次，我们需要验证定理 3.5.2.1 的 (1)，即 $\delta_{\bG}$ 在平凡 Weil 局部系统处取非零值。由构造，只需证明 $\delta(1)\neq 0$。显然
$\delta(1)=\prod_{\lambda\in\mathcal{R}_X}(1-\lambda)^{\dim(G)}$，而我们记得 $1\notin\mathcal{R}_X$（见注 4.3.1.1）。}

\pnum{接下来我们注意到，$\delta_{\bG}$ 非零的轨迹恰好是 $\LS^{\arthm,\circ,\nonres}_{\bG}$（按定义）。因此，由定理 4.3.3.1，我们有：}
\begin{equation}
  \Gamma(\LS^{\arthm}_{\bG},\omega)[\delta_{\bG}^{-1}]
  = \Gamma(\LS^{\arthm,\circ,\nonres}_{\bG},\omega)
  \cong \Gamma((\bG^{\ad}/\bG)^{\nonres},\omega),
\end{equation}
它在次数 $0$ 集中，因为它是 $\Gamma(\bG^{\ad}/\bG,\omega)\cong\Gamma(\bG^{\ad}/\bG,\OO)$ 在函数
$\bG^{\ad}/\bG\to\A^1$ 处的局部化，从而验证了定理 3.5.2.1 的假设 (2)。

\pnum{剩下要验证满射性（即定理 3.5.2.1 的 (3)）。我们将在本节余下部分做这件事；这里先作一些预备性的、起定向作用的评注。}

回忆引理 3.7.1.1 的设定。我们注意到有两个映射：
\[
\begin{gathered}
\Gamma(\LS^{\arthm}_{\bT},\omega)[\delta_{\bT}^{-1}]
\xrightarrow{\Eisspec}
\Gamma(\LS^{\arthm}_{\bG},\omega)[\delta_{\bG}^{-1}]
= \Gamma(\LS^{\arthm,\circ,\nonres}_{\bG},\omega)\\
\stackrel{(4.9)}{\cong}
\Gamma(\bG^{\ad}/\bG,\omega)[(\delta_{\bG}\circ\tau_{\bG})^{-1}]
= \Gamma((\bG^{\ad}/\bG)^{\nonres},\omega),
\end{gathered}
\]
以及：
\[
\begin{gathered}
\Gamma(\LS^{\arthm}_{\bT},\omega)[\delta_{\bT}^{-1}]
\stackrel{\text{定理 4.3.3.1}}{=}
\Gamma(\bT^{\ad}/\bT,\omega)[(\delta_{\bT}\circ\tau_{\bT})^{-1}]
\xrightarrow{\mathrm{Eis}^{\mathrm{spec},\mathrm{toy}}}
\Gamma(\bG^{\ad}/\bG,\omega)[(\delta_{\bG}\circ\tau_{\bG})^{-1}]\\
= \Gamma((\bG^{\ad}/\bG)^{\nonres},\omega).
\end{gathered}
\]
由推论 3.7.2.1，如果这两个映射重合，我们就完成了。

这一期望有点太天真：我们转而证明它们相差可逆的 $L$-值，这对我们的目的来说已经足够。

\subsection{$L$-值与迹}

在 §4.5.6 中，我们对 $L$-值作了一个晦涩的评注。在本小节中，我们将在范畴迹与 $L$-值之间建立精确的联系；这是我们所需的主要计算输入。

在下面，我们令 $H$ 表示 $\e$ 上的一个幂幺代数群。（实践中 $H=\bN$。）

\pnumt{类与迹。}{假设 $\mathcal{C}\in\DGCat_{\mathrm{cont}}$ 是一个可对偶的 DG 范畴，$T\colon\mathcal{C}\to\mathcal{C}$ 是一个自函子。}

令 $\mathcal{C}^c\subseteq\mathcal{C}$ 表示紧对象子范畴，并令 $\mathcal{C}^{c,T,\mathrm{lax}}$ 表示由对 $(F,\alpha)$ 构成的范畴，其中 $F\in\mathcal{C}^c$ 且 $\alpha\colon F\to T(F)$ 是 $\mathcal{C}$ 中的态射。

给定如上某个 $(F,\alpha)$，存在一个典范点 $\cl(F,\alpha)\in\Om^\infty\tr_{\mathcal{C}}(T)$。事实上，这由 §3.2.5 中迹的函子性推出：给 $\Vect$ 配备恒等自映射，把 $F$ 看作函子 $\Vect\to\mathcal{C}$，把 $\alpha$ 看作 lax 缠绕映射，于是函子性给出映射 $\e=\tr_{\Vect}(\id_{\Vect})\to\tr_{\mathcal{C}}(T)\in\Vect$，即一个点 $\cl(F,\alpha)\in\Om^\infty\tr_{\mathcal{C}}(T)$。

更一般地，我们回忆存在一个映射 $\cl(-)\colon K(\mathcal{C}^{c,T,\mathrm{lax}})\to\tr_{\mathcal{C}}(T)\in\Spectra$，从 $\mathcal{C}^{c,T,\mathrm{lax}}$ 的 $K$-理论谱到 $T$ 的迹（后者经由遗忘函子 $\Vect\to\Spectra$ 视为谱）。

\begin{notation}
注意 $\Vect^{c,\id,\mathrm{lax}}=\{(W,\phi):W\in\Vect^c,\ \phi\colon W\to W\}$ 是一个对称幺半范畴，因此在上述设定中典范地作用在 $\mathcal{C}^{c,T,\mathrm{lax}}$ 上。显式地，对 $(W,\phi)\in\Vect^{c,\id,\mathrm{lax}}$ 与 $(F,\alpha)\in\mathcal{C}^{c,T,\mathrm{lax}}$，$W\otimes F$ 配备自同态 $\phi\otimes\id+\id\otimes\alpha$。

在类映射下，有：
\begin{equation}
  \cl(W\otimes F,\phi\otimes\id+\id\otimes\alpha)=\tr_W(\phi)\cdot\cl(F,\alpha).
\end{equation}
（我们并不需要这一点，但这个等式可以很容易地提升为谱层面上的适当陈述。）
\end{notation}

\pnumt{问题的陈述。}{回忆 $H$ 是幂幺的。由 [2] 命题 3.3.2，$\LS^{\restr}_{H}$ 是拟紧代数叠；特别地，其结构层 $\OO_{\LS^{\restr}_{H}}\in\QCoh(\LS^{\restr}_{H})$ 是紧的（与非幂幺群的情形不同）。按滥用记号，我们仍用 $\OO_{\LS^{\restr}_{H}}$ 表示在完全忠实嵌入
$\QCoh(\LS^{\restr}_{H})\hookrightarrow\IndCoh(\LS^{\restr}_{H})$ 下 $\IndCoh(\LS^{\restr}_{H})$ 中的“同一个”对象（在 $\IndCoh$ 文献中通常记作“$\Xi$”）。}

我们有一个映射 $\tau_0\colon\mathbf{B}H\to\LS^{\restr}_{H}$，对应于平凡局部系统。于是可以构造
\[
  \tau^{\IndCoh}_{0,*}(\OO_{\mathbf{B}H})\in\IndCoh(\LS^{\restr}_{H}).
\]
注意 $\OO_{\LS^{\restr}_{H}}$ 与 $\tau^{\IndCoh}_{0,*}(\OO_{\mathbf{B}H})$ 这两个对象都是凝聚的，并带有显然的典范 Frobenius 等变结构。因此，我们可以构造它们的类：\footnote{由 $H$ 的幂幺性，注意 $\LS^{\arthm}_{H}$ 的每个点都是非共振的，即映射 $H^{\ad}/H\to\LS^{\arthm}_{H}$ 是同构。}
\[
  \cl(\OO_{\LS^{\restr}_{H}},\alpha),\quad
  \cl(\tau^{\IndCoh}_{0,*}(\OO_{\mathbf{B}H}),\alpha)
  \in\Om^\infty\Gamma(\LS^{\arthm}_{H},\omega).
\]
我们的目标是比较这两个类。

\pnumt{一个 $L$-值。}{令 $\zeta_X(t)=\sum_{n\ge 0}|X_0(\mathbf{F}_{q^n})|t^n\in\Q(t)\subseteq\Q((t))$ 表示曲线 $X$ 的 $\zeta$ 函数。}

我们提醒 $\zeta$ 函数具有形式：
\[
  \zeta_X(t)=\frac{p_X(t)}{(1-t)(1-qt)},\qquad p_X(t)\in\Q[t].
\]
令 $\zeta^{\star}_X(t)=(1-t)\cdot\zeta_X(t)$。由 Weil 猜想，$\zeta^{\star}_X(1)$ 非零，因而等于 $\zeta_X(t)$ 在 $t=1$ 处 Taylor 展开\footnote{注意——与数论中不同——我们是在变量 $t=q^{-s}$ 中展开，而不是在 $s$ 本身中展开。}的首项。

\pnumt{主引理。}{我们将证明：}

\begin{lem}
存在一个等价
\[
  \cl(\tau^{\IndCoh}_{0,*}(\OO_{\mathbf{B}H}),\alpha)
  \simeq \zeta^{\star}_X(1)^{\dim H}\cdot\cl(\OO_{\LS^{\restr}_{H}},\alpha)
  \in\Om^\infty\Gamma(\LS^{\arthm}_{H},\omega).
\]
\end{lem}

撇开同伦的信息，这个结果只是说上述两个点在集合
$\pi_0(\Om^\infty\Gamma(\LS^{\arthm}_{H},\omega))=H^0\Gamma(\LS^{\arthm}_{H},\omega)$ 中的像相等。

\begin{proof}
\textbf{第 1 步。}我们从一个玩具模型开始。

令 $V$ 与 $W$ 是配备自同态 $\phi_V$ 与 $\phi_W$ 的有限维向量空间。令 $\mathcal{Y}$ 表示叠（在 $\e$ 上）$V\times\Om_0W$，我们把它配备自映射 $\phi=\phi_V\times\Om_0\phi_W$；这里 $\Om_0W$ 是导出环路空间 $0\times_W 0$。

例如我们令 $\OO_0\in\Coh(\mathcal{Y})$ 表示原点的结构层，$\OO_V\in\Coh(\mathcal{Y})$ 表示 $V\subseteq\mathcal{Y}$ 的结构层，等等。

我们使用 $\mathcal{Y}$ 上 lax $\phi_*$-等变凝聚层范畴 $\Coh(\mathcal{Y})^{\phi_*,\mathrm{lax}}$，即
$\Coh(\mathcal{Y})^{\phi_*,\mathrm{lax}}=\{F\in\Coh(\mathcal{Y}),\ \alpha\colon F\to\phi_*(F)\}$（见 §4.6.1）。

Koszul 分解给出恒等式：
\begin{equation}
\begin{gathered}
  \sum(-1)^i[\Lambda^iW^{\vee}\otimes\OO_V]\simeq[\OO_{\mathcal{Y}}],\\
  \sum(-1)^j[\Lambda^jV^{\vee}\otimes\OO_V]\simeq[\OO_0].
\end{gathered}
\end{equation}
这里记号的含义如下。首先，$\Lambda^iW^{\vee}\otimes\OO_{\mathcal{Y}}\in\Coh(\mathcal{Y})^{\phi_*,\mathrm{lax}}$ 配备来自记号 4.6.1.1 的 lax 等变结构，其中 $\Lambda^iW^{\vee}$ 配备自同态 $\Lambda^i\phi_W^{\vee}$；对 $\Lambda^jV^{\vee}\otimes\OO_V$ 有类似的记号。记号 $[-]$ 用于表示 $K$-理论谱\footnote{对我们的目的而言，也可以简单地使用 Grothendieck 群。}$K(\Coh(\mathcal{Y})^{\phi_*,\mathrm{lax}})$ 中的类。

由 (4.10)，我们得到：
\[
\begin{gathered}
  \cl(\OO_0,\alpha)=\sum(-1)^j\tr(\Lambda^j\phi_V^{\vee})\cdot\cl(\OO_V,\alpha)
  =\sum(-1)^j\tr(\Lambda^j\phi_V)\cdot\cl(\OO_V,\alpha)\\
  =\det(\id_V-\phi_V)\cdot\cl(\OO_V,\alpha)\in\Om^\infty\Gamma(\mathcal{Y}^{\phi},\omega),
\end{gathered}
\]
其中 $\mathcal{Y}^{\phi}$ 表示导出不动点。

类似地，我们有：
\[
  \cl(\OO_{\mathcal{Y}},\alpha)=\det(\id_W-\phi_W)\cdot\cl(\OO_V,\alpha)
  \in\Om^\infty\Gamma(\mathcal{Y}^{\phi},\omega).
\]
比较这两个恒等式，我们得到：
\[
  \det(\id_W-\phi_W)\cdot\cl(\OO_0,\alpha)
  =\det(\id_V-\phi_V)\cdot\cl(\OO_{\mathcal{Y}},\alpha).
\]
现在假设 $\det(\id_W-\phi_W)$ 非零，从而在域 $\e$ 中可逆。我们得到：
\[
  \cl(\OO_0,\alpha)=\frac{\det(\id_V-\phi_V)}{\det(\id_W-\phi_W)}\cdot\cl(\OO_{\mathcal{Y}},\alpha)
  \in\Om^\infty\Gamma(\mathcal{Y}^{\phi},\omega).
\]
取 $V=H^1_{\et}(X)$ 与 $W=H^2_{\et}(X)$，配备它们各自的 Frobenius 自同态，我们注意到：
\[
  \frac{\det(\id_V-\phi_V)}{\det(\id_W-\phi_W)}=\zeta^{\star}_X(1),
\]
这正是 Grothendieck 迹公式。

\textbf{第 2 步。}其次，假设我们处于以下更一般的设定中。

我们假设 $\mathcal{Z}$ 是一个 QCA 叠，配备自同态 $\phi=\phi_{\mathcal{Z}}$，并且配备一个与映射 $\phi$ 缠绕的拟光滑映射 $\mathcal{Z}\to\mathcal{Y}=V\times\Om_0W$。

我们令 $\mathcal{Z}_0$ 表示 $\mathcal{Z}$ 在 $0\in\mathcal{Y}$ 上的纤维。我们注意到 $\mathcal{Z}_0$ 是最终余连通的，因此 $\OO_{\mathcal{Z}_0}$ 是 $\mathcal{Z}$ 上的凝聚层。

前面的分析于是表明：
\begin{equation}
  \cl(\OO_{\mathcal{Z}_0},\alpha)
  =\frac{\det(\id_V-\phi_V)}{\det(\id_W-\phi_W)}\cdot\cl(\OO_{\mathcal{Z}},\alpha)
  \in\Om^\infty\Gamma(\mathcal{Z}^{\phi},\omega)
\end{equation}
（假设 $1$ 不是 $\phi_W$ 的特征值）。

\textbf{第 3 步。}现在我们希望把上述形式体系应用于推导我们的断言。

选取一列嵌套子群 $\{1\}=H_0\subseteq H_1\subseteq\cdots\subseteq H_r=H$，其中每个 $H_i$ 在 $H$ 中正规，且 $H_{i+1}/H_i\simeq\mathbb{G}_a$。我们指出 $r=\dim(H)$。

然后定义代数叠：
\[
  \mathcal{Z}_i:=\LS^{\restr}_{H}\times_{\LS^{\restr}_{H/H_i}}\mathbf{B}(H/H_i),
\]
其中 $\mathbf{B}(H/H_i)\to\LS^{\restr}_{H/H_i}$ 是映射 $\tau_0$，即它对应于 $X$ 上的平凡 $H/H_i$-局部系统。注意 $\mathcal{Z}_r=\LS^{\restr}_{H}$，$\mathcal{Z}_0=\mathbf{B}H$，并且我们有闭嵌入：
\[
  \mathcal{Z}_0\hookrightarrow\mathcal{Z}_1\hookrightarrow\cdots\hookrightarrow\mathcal{Z}_r.
\]
我们可以把 $\mathcal{Z}_i$ 的定义改写如下。

注意 $H/H_i$ 作用在分类叠 $\mathbf{B}H_i$ 上；形式地说，这由纤维序列 $\mathbf{B}H_i\to\mathbf{B}H\to\mathbf{B}(H/H_i)$ 编码。展开定义，这诱导了 $H/H_i$ 在 $\LS^{\restr}_{H_i}$ 上的一个作用。于是我们有：
\[
  \mathcal{Z}_i\simeq\LS^{\restr}_{H_i}/(H/H_i).
\]
现在对每个 $i$，我们有一个图：
\[
\begin{tikzcd}[column sep=1.6em, row sep=1.4em]
  H/H_i \arrow[r] \arrow[d, Rightarrow]
    & H/H_{i+1} \arrow[r] \arrow[d, Rightarrow]
    & \mathbf{B}G_a \arrow[d, Rightarrow]\\
  \LS^{\restr}_{H_i} \arrow[r]
    & \LS^{\restr}_{H_{i+1}} \arrow[r]
    & \LS^{\restr}_{G_a}
\end{tikzcd}
\]
其中各行是纤维序列，且顶行是群的纤维序列。这里 $\mathbf{B}G_a$ 在 $\LS^{\restr}_{G_a}$ 上的作用由群叠的同态 $\mathbf{B}G_a\to\LS^{\restr}_{G_a}$ 诱导，后者对应于沿 $X\to\Spec(k)$ 拉回局部系统（即 $G_a$ 的映射 $\tau_0$）。在这个图中取商，并利用\footnote{我们指出，由纯度，存在一个与 Frobenius 相容的典范这样的分裂。}
\[
  \LS^{\restr}_{G_a}\simeq\mathbf{B}G_a\times H^1_{\et}(X)\times\Om_0H^2_{\et}(X),
\]
我们得到一个纤维方块：
\[
\begin{tikzcd}[column sep=2.2em, row sep=1.6em]
  \mathcal{Z}_i \arrow[r] \arrow[d] & \mathcal{Z}_{i+1} \arrow[d] & \\
  \Spec(k) \arrow[r] & \LS^{\restr}_{G_a}/\mathbf{B}G_a \arrow[r, equal]
    & H^1_{\et}(X)\times\Om_0H^2_{\et}(X)
\end{tikzcd}
\]
现在由前一步的归纳，我们得到结果。\footnote{形式地说，归纳应当在 $K$-理论类上进行，并推广 (4.11)。我们只在最后才映射到 Hochschild 同调。}
\end{proof}

\pnumt{一段扩展的题外话：经由范畴迹处理发散级数。}{我们解释思考上述引理 4.6.4.1 证明的一个一般格式。这些材料是非正式的，可以跳过。不过，我们相信这是一个重要的哲学要点，值得强调。}

粗略地说，想法是：所谓的\emph{范畴泛函分析}（例如关于 $\Perf$ 与 $\Coh$ 之间区别的精细考虑）通过范畴迹与实际分析（例如对无穷级数求和）相关联。引人注目的是，我们将看到 Hochschild 同调使我们有时能够“正确地”计算无穷和，而完全不必提及这些和所在的域 $\e$ 上的任何拓扑。

我们考虑以下几何设定。令 $\mathcal{Y}$ 是（$\e$ 上的）代数叠，我们假设它拟光滑且 QCA。假设 $\mathcal{Y}$ 配备自映射 $\phi\colon\mathcal{Y}\to\mathcal{Y}$。函子
\[
\begin{gathered}
  \Upsilon_{\mathcal{Y}}\colon\QCoh(\mathcal{Y})\to\IndCoh(\mathcal{Y}),\\
  F\mapsto F\otimes_{\OO_{\mathcal{Y}}}\omega_{\mathcal{Y}}
\end{gathered}
\]
保持紧对象，并缠绕源与靶的自映射 $\phi^*$ 与 $\phi^!$。此外，这个函子是 $\QCoh(\mathcal{Y})$-模范畴的态射。

现在回忆（例如 [15] §3.8.8），对于一个可对偶的 $\QCoh(\mathcal{Y})$-模范畴 $\mathcal{C}$ 与一个与 $\phi^*$ 适当相容的自函子 $T\colon\mathcal{C}\to\mathcal{C}$，存在一个典范对象
\[
  \tr^{\mathrm{enh}}(T)\in\QCoh(\mathcal{Y}^{\phi}),
\]
其基本性质是 $\Gamma(\mathcal{Y},\tr^{\mathrm{enh}}(T))=\tr(T)\in\Vect$。这一构造满足迹的通常函子性。我们有\footnote{在这个公式中，我们把 $\omega_{\mathcal{Y}^{\phi}}$ 看作 $\QCoh$ 的对象，而不是 $\IndCoh$ 的对象。换言之，我们隐含地取 $\IndCoh$ 中的“真正”对偶化层，并对它应用遗忘函子 $\Psi\colon\IndCoh(\mathcal{Y}^{\phi})\to\QCoh(\mathcal{Y}^{\phi})$。由于 $\omega_{\mathcal{Y}^{\phi}}$ 是线丛，这是一件相当温和的事情，所以我们不专门在记号中标明它。}如下公式：
\[
  \tr^{\mathrm{enh}}(\phi^*)=\OO_{\mathcal{Y}^{\phi}},\qquad
  \tr^{\mathrm{enh}}(\phi^!)=\omega_{\mathcal{Y}^{\phi}},
\]
于是 $\Upsilon_{\mathcal{Y}}$ 给出一个典范映射
\[
  \tau\colon\OO_{\mathcal{Y}^{\phi}}=\tr^{\mathrm{enh}}(\phi^*)\to\tr^{\mathrm{enh}}(\phi^!)=\omega_{\mathcal{Y}}
  \in\QCoh(\mathcal{Y}^{\phi}).
\]
（这个映射 $\tau$ 可以看作导出不动点上的一个弱 Calabi--Yau 结构。）

我们令 $\mathcal{Y}^{\phi,\mathrm{good}}\subseteq\mathcal{Y}^{\phi}$ 表示 $\tau$ 为同构的点的轨迹。注意 $\mathcal{Y}^{\phi,\mathrm{good}}$ 包含 $(\mathcal{Y}^{\mathrm{sm}})^{\phi}$（$\mathcal{Y}$ 的光滑轨迹的不动点），但一般更大：事实上可以验证 $\mathcal{Y}^{\phi,\mathrm{good}}$ 恰是 $\mathcal{Y}^{\phi}$ 的拟光滑轨迹。

现在，任何带自映射 $\alpha\colon F\to\phi^*(F)$ 的完美对象 $F\in\Perf(\mathcal{Y})$ 给出一个类
\[
  \cl(F,\alpha)_{\QCoh(\mathcal{Y})}\in\Om^\infty\Gamma(\mathcal{Y}^{\phi},\OO),
\]
即 $\phi$ 的不动点 $\mathcal{Y}^{\phi}$ 上的一个函数。在这个记号中，我们用下标 $\cl(-,-)_{\QCoh}$ 强调我们把 $F$ 看作 $\QCoh$ 的对象（这很快就会成为一个重要的区分）。这个函数可以相当显式地理解：在点 $y\in\mathcal{Y}^{\phi}$ 处，我们取所得映射的迹：
\begin{equation}
  \beta_y\colon y^*(F)\xrightarrow{\alpha}y^*(\phi^*(F))=\phi(y)^*(F)\cong y^*(F),
\end{equation}
其中最后一个同构使用了 $y$ 是不动点所隐含的等同 $y\simeq\phi(y)$。换言之，我们有：
\[
  \cl(F,\alpha)_{\QCoh(\mathcal{Y})}=(y\mapsto\tr(\beta_y)).
\]
现在改为假设 $F\in\Coh(\mathcal{Y})$，但仍配备映射 $\alpha\colon F\to\phi^*(F)$。因为 $F$ 在 $\QCoh(\mathcal{Y})$ 中可能不是紧的，我们不能再构造它在 $\Gamma(\mathcal{Y}^{\phi},\OO)$ 中的类。不过，我们可以扭一下并构造 $F\otimes\omega_{\mathcal{Y}}$，它属于 $\Coh$，因为 $\mathcal{Y}$ 拟光滑（从而 Gorenstein）。于是我们得到映射
\[
  \widetilde\alpha=\alpha\otimes\id\colon
  F\otimes\omega_{\mathcal{Y}}\to\phi^*(F)\otimes\omega_{\mathcal{Y}}
  =\phi^!(F\otimes\omega_{\mathcal{Y}}).
\]
因此，我们可以构造类
\[
  \cl(F\otimes\omega_{\mathcal{Y}},\widetilde\alpha)_{\IndCoh}
  \in\Gamma(\mathcal{Y}^{\phi},\omega_{\phi}).
\]
显然地，在 $F\in\Perf(\mathcal{Y})\subseteq\Coh(\mathcal{Y})$ 的特殊情形，我们有：
\begin{equation}
  \cl(F\otimes\omega_{\mathcal{Y}},\widetilde\alpha)_{\IndCoh}
  =\tau(\cl(F,\alpha)_{\QCoh}).
\end{equation}
遵循这个等式，我们定义\emph{正则化类}
\[
  \cl^{\mathrm{reg}}(F,\alpha)_{\QCoh}\in\Om^\infty\Gamma(\mathcal{Y}^{\phi,\mathrm{good}},\OO),
\]
为 $\cl(F\otimes\omega_{\mathcal{Y}},\widetilde\alpha)_{\IndCoh}$ 在如下复合下的像：
\[
  \Gamma(\mathcal{Y}^{\phi},\omega)\to
  \Gamma(\mathcal{Y}^{\phi,\mathrm{good}},\omega)
  \xrightarrow{\tau^{-1}}\Gamma(\mathcal{Y}^{\phi,\mathrm{good}},\OO).
\]
由构造，当 $F$ 完美时，这个正则化类与（$\cl(F,\alpha)_{\QCoh}$ 到 $\mathcal{Y}^{\phi,\mathrm{good}}$ 的限制）重合。

假设 $y\in\mathcal{Y}^{\phi}$。映射 (4.13) 仍有意义。然而，如果 $F$ 在 $y$ 附近不完美，那么虽然复形 (4.13) 在每个次数都是有限维的，它却下方无界，因此 $\beta_y$ 的迹没有良好定义。我们定义\emph{正则化迹}
\[
  \tr^{\mathrm{reg}}(\beta_y):=(\cl^{\mathrm{reg}}(F,\alpha)_{\QCoh})(y)
\]
为正则化类在 $y\in\mathcal{Y}^{\phi}$ 处的值。

启发式地，正则化迹可以看作给无穷和赋予一个实际值：
\begin{equation}
  \sum_{i\in\Z}(-1)^i\tr(H^i(\beta_y))\colon H^i(y^*(F))\to H^i(y^*(F))\in\e,
\end{equation}
其中我们重申每个求和项都有良好定义，求和项对 $i\gg 0$ 消失，但一般而言会出现无穷多项。

\begin{example}
我们解释在最简单的情形中这是如何运作的。假设 $\mathcal{Y}=\Om_0\A^1$，$\phi$ 是乘以数 $\lambda\in\e$。取层 $F$ 为 $\OO_0$，即点 $0\in\Om_0\A^1$ 的结构层。注意 $0$ 典范地是 $\phi$ 的不动点，所以我们可以把 $0$ 看作 $\mathcal{Y}^{\phi}$ 的点。由标准计算，$0^*(\OO_0)$ 在偶的非正上同调次数有 $1$ 维上同调，在这些次数之外上同调消失；此外，映射 $\beta_0$ 在 $H^{-2i}(0^*(\OO_0))$ 上作为乘以 $\lambda^i$ 作用。因此，(4.15) 中的和是几何级数 $\sum_{i\ge 0}\lambda^i$。我们强调这是一个形式表达式；目前，$\lambda$ 是域 $\lambda$ 的任意元素，绝不是“小”的。% [Tutor Note: 原文此处作 the field λ，应为系数域 e；疑为笔误，照录并登记。]

现在假设 $\lambda\neq 1$。则 $0\in\mathcal{Y}^{\phi,\mathrm{good}}$（事实上 $\mathcal{Y}^{\phi,\mathrm{good}}=\mathcal{Y}^{\phi}=\Spec(\e)=\{0\}$）。于是正则化迹 $\tr^{\mathrm{reg}}(\beta_0)$ 有良好定义，而引理 4.6.4.1 证明中第 1 步的（完全初等的）论证算出
\[
  \tr^{\mathrm{reg}}(\beta_0)=\frac{1}{1-\lambda}.
\]
换言之，我们给出了几何级数公式 $\sum\lambda^i\text{ “=” }\frac{1}{1-\lambda}$ 的直接的、纯粹代数的意义，而通常这需要我们知道 $\lambda^i$ 在某种适当的解析意义下 $i\to\infty$ 时趋于 $0$。
\end{example}

\begin{remark}[正则化迹与函数方程]
现在假设 $\phi\colon\mathcal{Y}\to\mathcal{Y}$ 事实上是同构。则 $\phi^*=\phi^!$（比如，作为限制在 $\Perf$ 或 $\Coh$ 上的函子）。因此，对于带 $\alpha\colon F\to\phi^*(F)$ 的 $F\in\Perf(\mathcal{Y})$，我们也得到映射 $\alpha'\colon F\to\phi^!(F)$。对 $y\in\mathcal{Y}^{\phi}$，我们得到典范映射
\[
  \gamma_y\colon y^!(F)\to y^!(F),
\]
以与 $\beta_y$ 相同的方式定义。由于 $F$ 完美，我们有
\[
  y^!(F)=y^!(F\otimes\omega_{\mathcal{Y}}\otimes\omega_{\mathcal{Y}}^{-1})
  =y^!(F\otimes\omega_{\mathcal{Y}})\otimes y^*(\omega_{\mathcal{Y}}^{-1})
  =y^*(F)\otimes y^*(\omega_{\mathcal{Y}})^{\otimes-1}.
\]
这个映射把 $\gamma_y$（对 $F$）与 $\beta_y$（对 $F$ 与 $\omega_{\mathcal{Y}}$ 两者）缠绕起来。如果我们令 $\epsilon_y$ 为映射
\[
  y^*(\omega_{\mathcal{Y}})\to y^*(\omega_{\mathcal{Y}})
\]
的迹（用 $\gamma_y$ 与显然的同构 $\omega_{\mathcal{Y}}\simeq\phi^!(\omega_{\mathcal{Y}})=\phi^*(\omega_{\mathcal{Y}})$ 构造），则我们发现
\[
  \epsilon_y\cdot\tr(\gamma_y)=\tr(\beta_y).
\]
现在我们可以完全像对凝聚 $F$ 那样，当 $y\in\mathcal{Y}^{\phi,\mathrm{good}}$ 时定义 $\tr^{\mathrm{reg}}(\gamma_y)$。我们得到一个同义反复的“函数方程”：
\[
  \epsilon_y\cdot\tr^{\mathrm{reg}}(\gamma_y)=\tr^{\mathrm{reg}}(\beta_y)
\]
（其中 $\epsilon_y$ 被看作一个 $\epsilon$-因子）。

让我们看看这个逻辑在例 4.6.5.1 的设定中如何展开。我们应当有 $\lambda\neq 0$，使得 $\phi$ 是同构。我们注意到 $0^!(\OO_0)$ 在偶的非负次数有上同调，而 $\gamma_0$ 在 $H^{2i}(0^!(\OO_0))$ 上的作用是乘以 $\lambda^{-i}$。此外，$0^*(\omega)=\e[-1]$，其中“$\gamma$”算子是乘以 $\lambda^{-1}$。因此，正则化迹 $\tr^{\mathrm{reg}}(\gamma_0)$ 启发式地使和
\[
  \tr^{\mathrm{reg}}(\gamma_0)\text{ “=” }\sum_{i\ge 0}\lambda^{-i}
\]
有意义，而 $\epsilon_0=-\lambda^{-1}$。

因此，在这个情形，我们的函数方程启发式地给出：
\[
  -\lambda^{-1}\sum_{i\ge 0}\lambda^{-i}\text{ “=” }
  \epsilon_0\tr^{\mathrm{reg}}(\gamma_0)=\tr^{\mathrm{reg}}(\beta_0)
  \text{ “=” }\sum_{i\ge 0}\lambda^i.
\]
\end{remark}

这个所得方程 $-\lambda^{-1}\sum_{i\ge 0}\lambda^{-i}\text{ “=” }\sum_{i\ge 0}\lambda^i$ 是发散级数研究中的一个经典式子；在 $\mathbb{C}$ 上，左边对 $|\lambda|>1$ 有定义，而右边对 $|\lambda|<1$ 有定义，但当然，这两个函数的解析延拓在其定义域上重合。

我们注意到，这类对发散级数的操作与代数曲线 $\zeta$ 函数的函数方程密切相关。

\begin{remark}
用上述思想的某种版本，在范畴上解释函数域上自守形式解析理论的某些解析方面，将是非常有意义的。
\end{remark}

\pnumt{变体。}{在实践中，我们需要 §4.6.4 讨论的一个轻微推广。}

像之前一样令 $H$ 是一个幂幺群，现在令 $S$ 是一个通过自同构作用在 $H$ 上的环面。（实践中，$S=\bT$ 作用在 $H=\bN$ 上。）为简洁起见，我们令 $Q$ 表示半直积 $S\ltimes H$。

在这个情形，我们定义有理映射
\[
  \zeta^{\star}_{X,H,S}\colon S\dashrightarrow\A^1
\]
通过公式
\[
  (s\in S)\mapsto
  \frac{\det\big(\id-\Ad_{s^{-1}}\otimes\phi_X\curvearrowright\mathfrak{h}\otimes H^1_{\et}(X)\big)}
       {\det\big(\id-\Ad_{s^{-1}}\otimes\phi_X\curvearrowright\mathfrak{h}\otimes H^2_{\et}(X)\big)}.
\]
这里 $\phi_X$ 是作用在 $X$ 的 \et{} 上同调上的 Frobenius，而我们在让 $\Ad_{-}$ 表示由 $S$ 在 $H$ 上的作用导出的 $S$ 在 $\mathfrak{h}$ 上的作用时，稍微滥用了记号。

\begin{remark}
注意 $\zeta^{\star}_{X,H,S}$ 在 $1\in S$ 处有定义，并在该处取值 $\zeta^{\star}_X(1)^{\dim H}$（见引理 4.6.4.1）。
\end{remark}

\begin{remark}
假设 $\mu_1,\dots,\mu_r\colon S\to\mathbb{G}_m$ 是 $S$ 在其表示 $\mathfrak{h}$ 中出现的特征标，按重数计算（故 $r=\dim(H)$）。则我们有：
\[
  \zeta^{\star}_{X,H,S}(s)=\prod_{i=1}^{r}\zeta^{\star}_X(\mu_i(s^{-1})).
\]
特别地，$\zeta^{\star}_{X,H,S}$ 的定义域是 $\bigcap_i\{s\in S\mid\mu_i(s)\neq q\}$，而 $\zeta^{\star}_{X,H,S}$ 在 $\bigcap_i\{s\in S\mid\mu_i(s)\notin\mathcal{R}_X\}$ 上（有定义且）可逆。
\end{remark}

同时引入以下记号会很方便：
\[
\begin{gathered}
  p_{X,H,S}(s)=\det\big(\id-\Ad_{s^{-1}}\otimes\phi_X\curvearrowright\mathfrak{h}\otimes H^1_{\et}(X)\big)
  =\prod_{i=1}^{r}p_X(\mu_i(s^{-1})),\\
  q_{X,H,S}(s)=\det\big(\id-\Ad_{s^{-1}}\otimes\phi_X\curvearrowright\mathfrak{h}\otimes H^2_{\et}(X)\big)
  =\prod_{i=1}^{r}(1-q\cdot\mu_i(s^{-1})),
\end{gathered}
\]
于是 $p_{X,H,S}$ 与 $q_{X,H,S}$ 是 $S$ 上的（正则）函数，且
$\zeta^{\star}_{X,H,S}=\dfrac{p_{X,H,S}}{q_{X,H,S}}$。

我们引入记号：
\[
  \LS^{\restr}_{Q;S}:=\LS^{\restr}_{Q}\times_{\LS^{\restr}_{S}}\mathbf{B}S
  =(\LS^{\restr}_{H})/S.
\]
我们观察到 $\LS^{\restr}_{Q;S}$ 是拟紧代数叠（由 $H$ 的幂幺性）。我们令 $\LS^{\arthm}_{Q;S}$ 表示 $\LS^{\restr}_{Q;S}$ 的 Frobenius 不动点。显式地，我们有：
\[
  \LS^{\arthm}_{Q;S}=\LS^{\arthm}_{Q}\times_{\LS^{\arthm}_{S}}S^{\ad}/S.
\]
我们令 $\tau_1\colon\mathbf{B}Q\to\LS^{\restr}_{Q;S}$ 表示显然的映射。

\begin{remark}
更显式地说，我们提醒 $S^{\ad}/S\xrightarrow{\sim}\LS^{\arthm,\circ}_{S}$（由定理 4.3.3.1），并回忆 $S$ 是环面。因此，$\LS^{\arthm}_{Q;S}$ 是 $\LS^{\arthm}_{Q}$ 中单位元的连通分支。
\end{remark}

\begin{lem}
存在一个等价
\[
  q_{X,H,S}\cdot\cl(\tau^{\IndCoh}_{1,*}(\OO_{\mathbf{B}Q}),\alpha)
  \simeq p_{X,H,S}\cdot\cl(\OO_{\LS^{\restr}_{Q;S}},\alpha)
  \in\Om^\infty\Gamma(\LS^{\arthm}_{Q;S},\omega).
\]
\end{lem}

\begin{remark}
非正式地，这个引理应被理解为说：
\[
  \cl(\tau^{\IndCoh}_{1,*}(\OO_{\mathbf{B}Q}),\alpha)
  \simeq\zeta^{\star}_{X,H,S}\cdot\cl(\OO_{\LS^{\restr}_{Q;S}},\alpha).
\]
\end{remark}

\begin{proof}[引理 4.6.6.4 的证明]
其证明本质上与引理 4.6.4.1 的证明相同。差别如下。

第一，在第 1 步中，应当假设 $V$ 与 $W$ 是 $S$-表示，并应当在 (4.11) 中计入 $S$-作用。对偶表示 $V^{\vee}$ 与 $W^{\vee}$ 出现在 (4.11) 中，这解释了上文 $\zeta^{\star}_{X,H,S}$ 的定义中出现 $s^{-1}$ 而非 $s$。

第二，应当注意引理 4.6.4.1 证明中的子群 $H_i$ 可以取为在 $S$-作用下不变的（证明：把 $S$ 在 $\mathfrak{h}/[\mathfrak{h},\mathfrak{h}]$ 上的作用对角化，然后归纳）。

除此之外，论证逐字相同。
\end{proof}

% ============================================================
%  ch04d：§4.7 结论（含定理 4.7.2.1 及证明）与参考文献
%  （原文 pp.39 中段–42；起于 §4.7，止于参考文献表与作者地址）
% ============================================================
\subsection{结论}

我们现在回到 §4.5 的设定。

\pnum{令 $\LS^{\arthm}_{\bT,\delta_{\bT}\neq0}\subseteq\LS^{\arthm}_{\bT}$ 表示 $\delta_{\bT}$ 的非零轨迹。}

我们有有理函数：
\[
  \LS^{\arthm,\circ}_{\bT}
  \stackrel{\text{定理 4.3.3.1}}{=}
  \bT^{\ad}/\bT=\bT\times^{\mathbf{B}}\bT
  \xrightarrow{p_1}\bT
  \xrightarrow{\zeta_{X,\bN,\bT}}\A^1,
  % [Tutor Note: 原文此处排作 \times^{\mathbf{B}}（粗体分类叠记号），疑为 \times_{\mathbf{B}\check T} 的排版变体；照录并登记。]
\]
它在 $\LS^{\arthm}_{\bT,\delta_{\bT}\neq0}$ 上显然有定义且可逆。按滥用记号，我们也用 $\zeta_{X,\bN,\bT}$ 记所得映射
\[
  \zeta_{X,\bN,\bT}\colon\LS^{\arthm}_{\bT,\delta_{\bT}\neq0}\to\A^1\setminus0.
\]

\pnum{我们现在证明以下结果：}

\begin{thm}
存在一个交换图：
{\scriptsize
\[
\begin{tikzcd}[column sep=3.5em, row sep=3.2em]
  \Gamma(\LS^{\arthm}_{\bT,\delta_{\bT}\neq0},\omega) \arrow[r, equal]
    & \Gamma(\LS^{\arthm}_{\bT},\omega)[\delta_{\bT}^{-1}] \arrow[rr, "\Eisspec"]
    & & \Gamma((\bG^{\ad}/\bG)^{\nonres},\omega)\\
  \Gamma(\LS^{\arthm}_{\bT,\delta_{\bT}\neq0},\omega) \arrow[r, equal]
    \arrow[u, "\cong"', "\zeta_{X,\bN,\bT}\cdot-"]
    & \Gamma(\LS^{\arthm}_{\bT},\omega)[\delta_{\bT}^{-1}] \arrow[r, "\mathrm{Eis}^{\mathrm{spec},\mathrm{toy}}"']
    & \Gamma(\bG^{\ad}/\bG,\omega)[(\delta_{\bG}\circ\tau_{\bG})^{-1}]
      \arrow[r, equal] \arrow[ur, "id"']
    & \Gamma((\bG^{\ad}/\bG)^{\nonres},\omega)
\end{tikzcd}
\]
}
\end{thm}

由推论 3.7.2.1，这显然给出 §4.5.6 中所需的满射性。因此，只剩下证明这个定理。

\begin{proof}[定理 4.7.2.1 的证明]
\textbf{第 1 步。}由紧合态射连接的叠的如下交换图
\[
\begin{tikzcd}[column sep=2.6em, row sep=1.8em]
  \mathbf{B}\bB \arrow[r, "\tau_{\bB}"] \arrow[d, "p^{\mathrm{toy}}"]
    & \LS^{\restr}_{\bB} \arrow[d, "p"]\\
  \mathbf{B}\bG \arrow[r, "\tau_{\bG}"]
    & \LS^{\restr}_{\bG}
\end{tikzcd}
\]
给出所得两个函子
\[
  \Rep(\bB)\to\IndCoh(\LS^{\restr}_{\bG})
\]
的一个等同，且保持紧对象。

对 Frobenius 取迹，这给出交换图
\begin{equation}
\begin{tikzcd}[column sep=3.2em, row sep=1.8em]
  \Gamma(\bB^{\ad}/\bB,\omega) \arrow[r] \arrow[d]
    & \Gamma(\LS^{\arthm}_{\bB},\omega) \arrow[d]\\
  \Gamma(\bG^{\ad}/\bG,\omega) \arrow[r]
    & \Gamma(\LS^{\arthm}_{\bG},\omega).
\end{tikzcd}
\end{equation}

\textbf{第 2 步。}其次，构造交换方块：
\[
\begin{tikzcd}[column sep=2.6em, row sep=1.8em]
  \mathbf{B}\bB \arrow[r, "\tau_{\bB}"] \arrow[d, "q^{\mathrm{toy}}"]
    & \LS^{\restr}_{\bB} \arrow[d, "q"]\\
  \mathbf{B}\bT \arrow[r, "\tau_{\bT}"]
    & \LS^{\restr}_{\bT}.
\end{tikzcd}
\]
我们得到自然变换
\[
  q^{*,\IndCoh}\tau_{\bT,*}^{\IndCoh}
  \to \tau_{\bB,*}^{\IndCoh}q^{\mathrm{toy},*,\IndCoh},
\]
它是函子
\[
  \Rep(\bT)\to\IndCoh(\LS^{\restr}_{\bB})
\]
之间的自然变换。但这个函子\emph{不是}同构，因此仅凭迹的函子性对它几乎无话可说。

尽管如此，我们断言有交换图：
\begin{equation}
\begin{tikzcd}[column sep=3.95em, row sep=4.4em, cells={font=\small}]
  \Gamma(\bT^{\ad}/\bT,\omega) \arrow[r, "p_{X,\bN,\bT}\cdot-"] \arrow[d, equal]
    & \Gamma(\bT^{\ad}/\bT,\omega) \arrow[r, "\tr(q^{*,\IndCoh}\tau_{\bT,*}^{\IndCoh})"]
    & \Gamma(\LS^{\arthm}_{\bB},\omega)\\
  \Gamma(\bT^{\ad}/\bT,\omega) \arrow[r, "q_{X,\bN,\bT}\cdot-"]
    & \Gamma(\bT^{\ad}/\bT,\omega) \arrow[ur, "\tr(\tau_{\bB,*}^{\IndCoh}q^{\mathrm{toy},*,\IndCoh})"']
    & 
\end{tikzcd}
\end{equation}

为构造这个图，注意这些映射自然是 $\Gamma(\LS^{\arthm}_{\bT},\OO)$-模的态射。这个代数显然通过其因子
\[
  \Gamma(\LS^{\arthm,\circ}_{\bT},\OO)=\Gamma(\bT^{\ad}/\bT,\OO)
\]
作用在 $\Gamma(\bT^{\ad}/\bT,\omega)$（我们想构造的图的源）上。因此，只需给出交换图
\[
\begin{tikzcd}[column sep=3.95em, row sep=4.4em, cells={font=\small}]
  \Gamma(\bT^{\ad}/\bT,\omega) \arrow[r, "p_{X,\bN,\bT}\cdot-"] \arrow[d, equal]
    & \Gamma(\bT^{\ad}/\bT,\omega) \arrow[r, "\tr(q^{*,\IndCoh}\tau_{\bT,*}^{\IndCoh})"]
    & \Gamma(\LS^{\arthm}_{\bB},\omega) \arrow[r, "\mathrm{proj}"] \arrow[d, equal]
    & \Gamma(\LS^{\arthm,\circ}_{\bB},\omega) \arrow[d, equal]\\
  \Gamma(\bT^{\ad}/\bT,\omega) \arrow[r, "q_{X,\bN,\bT}\cdot-"]
    & \Gamma(\bT^{\ad}/\bT,\omega) \arrow[r, "\tr(\tau_{\bB,*}^{\IndCoh}q^{\mathrm{toy},*,\IndCoh})"]
    & \Gamma(\LS^{\arthm}_{\bB},\omega) \arrow[r, "\mathrm{proj}"]
    & \Gamma(\LS^{\arthm,\circ}_{\bB},\omega),
\end{tikzcd}
\]
它是 $\Gamma(\bT^{\ad}/\bT,\OO)$-模的交换图（因为
$\LS^{\arthm,\circ}_{\bB}=\LS^{\arthm}_{\bB}\times_{\LS^{\arthm}_{\bT}}\LS^{\arthm,\circ}_{\bT}$）。

回忆 $H^{\ad}/H$ 上的对偶化层典范平凡。这个平凡化可以如下构造：等价 $\QCoh(\mathbf{B}H)\simeq\IndCoh(\mathbf{B}H)$ 在 Hochschild 同调上给出同构
$\Gamma(H^{\ad}/H,\OO)\simeq\Gamma(H^{\ad}/H,\omega)$。令
$\vol_H\in\Om^\infty\Gamma(H^{\ad}/H,\omega)$ 表示所得的生成元——显式地，它是 $\Rep(H)$ 的单位对象的类。因此，为得到上面的图，只需给出同构\footnote{技术上，为使 $p_{X,\bN,\bT}$ 与 $q_{X,\bN,\bT}$ 作用在这个等式的外侧（而不是作用在 $\vol_{\bT}$ 本身上），我们把所见的对象都看作上述意义上的 $\Gamma(\bT^{\ad}/\bT,\OO)$-模。}
\[
\begin{gathered}
  p_{X,\bN,\bT}\cdot\mathrm{proj}
  \Big(\tr\big(q^{*,\IndCoh}\tau_{\bT,*}^{\IndCoh}\big)(\vol_{\bT})\Big)\\
  =q_{X,\bN,\bT}\cdot\mathrm{proj}
  \Big(\tr\big(\tau_{\bB,*}^{\IndCoh}q^{\mathrm{toy},*,\IndCoh}\big)(\vol_{\bT})\Big)
  \in\Om^\infty\Gamma(\LS^{\arthm,\circ}_{\bB},\omega).
\end{gathered}
\]
由构造，我们有：
\[
  \tr\big(\tau_{\bB,*}^{\IndCoh}q^{\mathrm{toy},*,\IndCoh}\big)(\vol_{\bT})
  =\cl\big(\tau_{\bB,*}^{\IndCoh}\OO_{\mathbf{B}\bB},\alpha\big)
  \in\Om^\infty\Gamma(\LS^{\arthm}_{\bB},\omega).
\]
类似地，由基变换，我们有：
\[
  \tr\big(q^{*,\IndCoh}\tau_{\bT,*}^{\IndCoh}\big)(\vol_{\bT})
  =\cl\big(\OO_{\LS^{\restr}_{\bB}\times_{\LS^{\restr}_{\bT}}\mathbf{B}\bT},\alpha\big)
  =\cl\big(\OO_{\LS^{\restr}_{\bB;\bT}},\alpha\big)
  \in\Om^\infty\Gamma(\LS^{\arthm}_{\bB},\omega).
\]
于是该恒等式由引理 4.6.6.4 推出，注意
$\LS^{\arthm}_{\bB;\bT}=\LS^{\arthm,\circ}_{\bB}$（见注 4.6.6.3）。

\textbf{第 3 步。}拼接图 (4.16) 与 (4.17)，我们得到交换图：
\[
\begin{tikzcd}[column sep=1.1em, row sep=3.4em]
  \Gamma(\bT^{\ad}/\bT,\omega) \arrow[r, "p_{X,\bN,\bT}\cdot-"] \arrow[d, equal]
    & \Gamma(\bT^{\ad}/\bT,\omega) \arrow[r, "\tr(\tau_{\bT,*}^{\IndCoh})"]
    & \Gamma(\LS^{\arthm}_{\bT},\omega) \arrow[r, "\Eisspec"]
    & \Gamma(\LS^{\arthm}_{\bG},\omega) \arrow[d, equal]\\
  \Gamma(\bT^{\ad}/\bT,\omega) \arrow[r, "q_{X,\bN,\bT}\cdot-"]
    & \Gamma(\bT^{\ad}/\bT,\omega) \arrow[r, "\mathrm{Eis}^{\mathrm{spec},\mathrm{toy}}"']
    & \Gamma(\bG^{\ad}/\bG,\omega) \arrow[r, "\tr(\tau_{\bG,*}^{\IndCoh})"]
    & \Gamma(\LS^{\arthm}_{\bG},\omega).
\end{tikzcd}
\]
这个图细化了我们本该证明的定理。
\end{proof}

\begin{thebibliography}{99}
\bibitem{arinkin-gaitsgory} D. Arinkin and D. Gaitsgory, \emph{Singular support of coherent sheaves and the geometric Langlands conjecture}, Selecta Math. \textbf{21} (2015), no. 1, 1--199. Springer.
\bibitem{AGKRV1} D. Arinkin, D. Gaitsgory, D. Kazhdan, S. Raskin, N. Rozenblyum, and Y. Varshavsky, \emph{The stack of local systems with restricted variation and geometric Langlands theory with nilpotent singular support}, preprint, 2020. arXiv:2010.01906.
\bibitem{AGKRV2} D. Arinkin, D. Gaitsgory, D. Kazhdan, S. Raskin, N. Rozenblyum, and Y. Varshavsky, \emph{Duality for automorphic sheaves with nilpotent singular support}, preprint, 2020. arXiv:2012.07665.
\bibitem{AGKRV3} D. Arinkin, D. Gaitsgory, D. Kazhdan, S. Raskin, N. Rozenblyum, and Y. Varshavsky, \emph{Automorphic functions as the trace of Frobenius}, preprint, 2021. arXiv:2102.07906.
\bibitem{beilinson} A. Beilinson, \emph{Constructible sheaves are holonomic}, Selecta Math. \textbf{22} (2016), no. 4, 1797--1819. Springer.
\bibitem{braverman-gaitsgory} A. Braverman and D. Gaitsgory, \emph{Deformations of Local Systems and Eisenstein Series}, Geom. Funct. Anal., 2008. Springer Science+Business Media.
\bibitem{benzvi-nadler} D. Ben-Zvi and D. Nadler, \emph{Betti geometric Langlands}, in Proceedings of Symposia in Pure Mathematics, vol. 97, 2018, pp. 3--41.
\bibitem{drinfeld-gaitsgory1} V. Drinfeld and D. Gaitsgory, \emph{On some finiteness questions for algebraic stacks}, Geom. Funct. Anal. \textbf{23} (2013), no. 1, 149--294. Springer.
\bibitem{drinfeld-gaitsgory2} V. Drinfeld and D. Gaitsgory, \emph{Geometric constant term functor(s)}, Selecta Math. \textbf{22} (2016), no. 4, 1881--1951. Springer.
\bibitem{faergeman-raskin} J. Faergeman and S. Raskin, \emph{Non-vanishing of geometric Whittaker coefficients for reductive groups}, J. Amer. Math. Soc., 2025.
\bibitem{fargues-scholze} L. Fargues and P. Scholze, \emph{Geometrization of the local Langlands correspondence}, preprint, 2021. arXiv:2102.13459.
\bibitem{gaitsgory-indcoh} D. Gaitsgory, \emph{Ind-coherent sheaves}, Mosc. Math. J. \textbf{13} (2013), no. 3, 399--528.
\bibitem{gaitsgory-outline} D. Gaitsgory, \emph{Outline of the proof of the geometric Langlands conjecture for $\mathrm{GL}(2)$}, Ast\'erisque, 2015. Centre National de la Recherche Scientifique.
\bibitem{gaitsgory-sheaves} D. Gaitsgory, \emph{Sheaves of categories and the notion of 1-affineness}, in Stacks and categories in geometry, topology, and algebra, vol. 643, 2015, pp. 127--225. American Mathematical Society, Providence, RI.
\bibitem{GKRV} D. Gaitsgory, D. Kazhdan, N. Rozenblyum, and Y. Varshavsky, \emph{A toy model for the Drinfeld--Lafforgue shtuka construction}, Indag. Math., 2021. Elsevier.
\bibitem{gaitsgory-raskin} D. Gaitsgory and S. Raskin, \emph{Proof of the geometric Langlands conjecture I: construction of the functor}, arXiv preprint arXiv:2405.03599, 2024.
\bibitem{gaitsgory-rozenblyum1} D. Gaitsgory and N. Rozenblyum, \emph{A study in derived algebraic geometry: Volume I: correspondences and duality}, Mathematical Surveys and Monographs, vol. 221, American Mathematical Society, 2019.
\bibitem{gaitsgory-rozenblyum2} D. Gaitsgory and N. Rozenblyum, \emph{A study in derived algebraic geometry: Volume II: deformations, Lie theory and formal geometry}, Mathematical Surveys and Monographs, vol. 221, American Mathematical Society, 2020.
\bibitem{harder} G. Harder, \emph{\"Uber die Galoiskohomologie halbeinfacher algebraischer Gruppen. III}, Walter de Gruyter, Berlin/New York, 1975.
\bibitem{lafforgue-l} L. Lafforgue, \emph{Chtoucas de Drinfeld et correspondance de Langlands}, Invent. Math. \textbf{147} (2002), no. 1, 1--241. Springer.
\bibitem{lafforgue-v1} V. Lafforgue, \emph{Quelques calculs reli\'es \`a la correspondance de Langlands g\'eom\'etrique pour $\mathbf{P}^1$}, preprint, 2009. Available at \url{http://www.math.jussieu.fr/~vlafforg}.
\bibitem{lafforgue-v2} V. Lafforgue, \emph{Chtoucas pour les groupes r\'eductifs et param\'etrisation de Langlands globale}, J. Amer. Math. Soc. \textbf{31} (2018), no. 3, 719--891.
\bibitem{laumon1} G. Laumon, \emph{Correspondance de Langlands g\'eom\'etrique pour les corps de fonctions}, Duke Math. J. \textbf{54} (1987), no. 2, 309--359. Duke University Press.
\bibitem{laumon2} G. Laumon, \emph{Transformation de Fourier g\'en\'eralis\'ee}, arXiv preprint alg-geom/9603004, 1996.
\bibitem{lurie} J. Lurie, \emph{Spectral algebraic geometry}, preprint, 2018. Available at \url{https://www.math.ias.edu/~lurie/papers/moduli.pdf}.
\bibitem{moeglin-waldspurger} C. Moeglin and J.-L. Waldspurger, \emph{Spectral decomposition and Eisenstein series}, Cambridge University Press, 1995. London Mathematical Society Monographs, Number 113.
\bibitem{nadler-yun} D. Nadler and Z. Yun, \emph{Spectral action in Betti Geometric Langlands}, Israel J. Math. \textbf{232} (2019), no. 1.
\bibitem{deligne} P. Deligne, \emph{S\'eminaire de G\'eom\'etrie Alg\'ebrique du Bois Marie--Cohomologie \'etale--(SGA 4 1/2)}, Springer, Berlin/New York, 1977.
\bibitem{varshavsky} Y. Varshavsky, \emph{Moduli spaces of principal F-bundles}, Selecta Math. (N.S.) \textbf{10} (2004), no. 1, 131. Springer.
\bibitem{xue-finiteness} C. Xue, \emph{Finiteness of cohomology groups of stacks of shtukas as modules over Hecke algebras, and applications}, \'Epijournal de G\'eom\'etrie Alg\'ebrique \textbf{4}, 2020. Episciences.org.
\bibitem{xue-smoothness} C. Xue, \emph{Smoothness of cohomology sheaves of stacks of shtukas}, preprint, 2020. arXiv:2012.12833.
\bibitem{zhu} X. Zhu, \emph{Coherent sheaves on the stack of Langlands parameters}, preprint, 2020. arXiv:2008.02998.

\medskip
\noindent Yale University, Department of Mathematics, 219 Prospect St, New Haven, CT 06511\\
\noindent\emph{Email address:} \texttt{sam.raskin@yale.edu}
\end{thebibliography}


\end{document}
